\documentclass[11pt]{article}
\usepackage[T1]{fontenc}
\usepackage{lmodern}
\usepackage[letterpaper,margin=1in,headheight=15pt]{geometry}
\usepackage{amsmath,amssymb,amsthm,amscd,mathtools}
\usepackage{microtype,booktabs,array,longtable}
\usepackage{tikz-cd}
\usetikzlibrary{arrows.meta,positioning,calc}
\usepackage{xcolor}
\definecolor{linkink}{HTML}{735632}
\definecolor{softink}{HTML}{5E675D}
\usepackage[colorlinks=true,linkcolor=linkink,urlcolor=linkink,citecolor=linkink]{hyperref}
\usepackage{bookmark,fancyhdr}
\pagestyle{fancy}\fancyhf{}
\fancyhead[L]{\small\textsc{Quotients and Reconstruction}}
\fancyhead[R]{\small\textsc{Tensorgami}}
\fancyfoot[C]{\thepage}
\renewcommand{\headrulewidth}{0.2pt}
\setlength{\parskip}{4pt plus 1pt minus 1pt}
\setlength{\parindent}{0pt}
\setlength{\emergencystretch}{2em}
\linespread{1.035}\allowdisplaybreaks[1]\raggedbottom
\newtheorem{theorem}{Theorem}[section]
\newtheorem{proposition}[theorem]{Proposition}
\newtheorem{lemma}[theorem]{Lemma}
\newtheorem{corollary}[theorem]{Corollary}
\theoremstyle{definition}
\newtheorem{definition}[theorem]{Definition}
\newtheorem{example}[theorem]{Example}
\theoremstyle{remark}\newtheorem{remark}[theorem]{Remark}
\newcommand{\id}{\mathrm{id}}
\newcommand{\kk}{k}
\newcommand{\nml}{\trianglelefteq}
\newcommand{\iso}{\xrightarrow{\sim}}
\DeclareMathOperator{\Ker}{ker}
\DeclareMathOperator{\im}{im}
\DeclareMathOperator{\Span}{span}
\DeclareMathOperator{\Hom}{Hom}
\DeclareMathOperator{\Ext}{Ext}
\DeclareMathOperator{\Aut}{Aut}
\DeclareMathOperator{\rank}{rank}
\setcounter{tocdepth}{2}
\hypersetup{pdftitle={Quotients, Refinement, and Reconstruction},pdfsubject={Exact Malcev categories, Jordan--Holder, coherent quotients, descent and derived structure},pdfauthor={}}
\title{\textbf{Quotients, Refinement,\\and Reconstruction}\\[9pt]\large From exact Mal'cev categories to homotopy and linearization}
\author{}\date{\small October 22, 2026}
\begin{document}
\maketitle\thispagestyle{empty}
\begin{abstract}
A quotient makes some distinctions disappear. A relation remembers which
distinctions are to disappear, and exactness makes that record into an
object. The Mal'cev condition makes compatible relations complete their
rectangles. From these constructions we derive a single refinement theorem
and a single Jordan--H\"older theorem. Groups, modules, and other algebraic
structures enter by verifying the hypotheses; invariance of dimension
then follows from the flags of two bases.

We next examine what the resulting factors retain and what they lose.
Comparison pages pair simple layers, additive invariants discard extension
data, and explicit probes recover nilpotent operators. Counting flags
produces an associative Hall multiplication. Finally, homotopy quotients
retain automorphisms and coherent gluing, descent recovers actions, and
linearization supplies coefficient categories, derived extensions and
cotangent information. The ordinary comparison is proved from the stated
categorical axioms. The higher chapters distinguish imported foundational
theorems from the deductions proved here. Each enlargement answers a
specific question that the preceding construction leaves open.
\end{abstract}
\clearpage
% BEGIN INPUT 00-reading-guide.tex
\section*{How to read the construction}
\addcontentsline{toc}{section}{How to read the construction}

The starting question is elementary: why do two finite bases have the same
number of vectors? One answer exchanges vectors by triangular substitutions.
We begin by explaining those substitutions as invertible maps, and the
familiar algebraic maneuvers as commutative diagrams. We then ask for a
construction that can compare layers even when there are no vectors and
no addition of maps.

The answer uses five ingredients with distinct jobs. Finite limits impose
compatible equations. Regular images forget intermediate witnesses in a
controlled way. Exactness turns an equivalence relation into a quotient.
The Mal'cev condition makes every compatible relation difunctional:
three related corners force the fourth. Pointedness selects the zero class
needed to organize a filtration. These hypotheses will be defined and used
one at a time; ``exact'' here means Barr-exact, not Quillen-exact.

The recurring construction is
\[
 \text{one correspondence}
 \quad\longmapsto\quad
 \text{two ambiguity relations}
 \quad\longmapsto\quad
 \text{one canonically identified quotient}.
\]
Applying it at every crossing of two filtrations produces the refinement
rectangle. Jordan--H\"older follows when simple factors prevent a
nontrivial further decomposition. Groups and modules are verified
instances of this one proof. We do not assume a prior dimension theorem,
an invariant length, or a Grothendieck group to justify it.

For a first reading focused on that result, read
Sections~\ref{f:begin}--\ref{e:section}. The relation notation requires
some attention, but the convention never changes: $SR$ means first $R$,
then $S$. Each abstract formula is accompanied by its meaning and its
typed maps. The elementary algebra supplies the concrete models that make
the common construction intelligible.

For the spectral viewpoint, continue through the module toolbox and the
comparison pages. The groups-with-operators calculation works with actual
nonabelian quotients because its differential images have proved
normality. The subsequent module chapter establishes the more general
two-term filtered-map calculus. Neither calculation is used to patch a
gap in the categorical refinement proof.

For reconstruction, follow the next four chapters. Simple factors give
the universal additive record, but cannot distinguish all extensions.
Explicit attachment parameters, maps from small probes, and counts of
flags recover progressively different information. The nilpotent
reconstruction theorem is proved by constructing chain bases; it is not
inferred from a classification theorem left unstated.

The final part changes the kind of quotient being taken. A homotopy
quotient retains automorphisms and coherent transport that an orbit set
forgets. Its descent theorem, coefficient categories, and derived
extension and cotangent constructions offer additional capabilities.
There the foundational equivalences from higher category theory are
\emph{explicitly quoted inputs}. Our deductions from them are proved,
with their assumptions and maps. This makes the logical boundary visible:
the text is self-contained for its ordinary categorical comparison,
relative to the categorical axioms it defines; it is not a reconstruction
of all of higher topos theory from set theory.

No claim of historical novelty is made. The purpose is a connected
exposition in which constructions retain, forget, and recover specified
data. The dependency record at the end separates theorems proved here,
framework results imported from the literature, illustrative calculations,
and the limits of each conclusion.

% END INPUT 00-reading-guide.tex

\clearpage\begingroup\setlength{\parskip}{0pt}\tableofcontents\endgroup
\clearpage\part{One comparison framework}

% BEGIN INPUT 01-maps-relations-exactness.tex
\section{From algebraic maneuvers to maps and relations}
\label{f:begin}

When we write $a=b+c$ and then $b=a-c$, something more organized than
symbol manipulation is happening. We have applied an invertible map. When
we then write $d(a-c)=da-dc$, we have used a map that preserves an
operation. These two ideas---invertibility and preservation---will remain
visible as the objects become more general.

\subsection{An elementary maneuver is a diagram}

Let $(V,+,0)$ be an abelian group. For $c\in V$, translation is the map of
sets $\tau_c:V\to V$, $x\mapsto x+c$. Associativity, the identity law,
and the inverse law give
\[
\tau_{-c}\tau_c(x)=(x+c)+(-c)=x+(c-c)=x.
\]
The same calculation in the other order proves
$\tau_{-c}=\tau_c^{-1}$. Thus $a=b+c$ means $a=\tau_c(b)$, and solving
for $b$ means applying the inverse. Translation need not preserve the
zero element, so it is generally an isomorphism of sets, not an
automorphism of abelian groups.

Associativity itself says that two composites are the same map:
\[
\begin{tikzcd}[column sep=large]
V\times V\times V \arrow[r,"+\,\times\,\id"]
  \arrow[d,"\id\,\times\,+"'] & V\times V \arrow[d,"+"]\\
V\times V \arrow[r,"+"'] & V.
\end{tikzcd}
\]
Here the usual canonical identification of the two parenthesizations of
the triple product is understood. The diagram says that making one
intermediate sum rather than the other does not change the result.

If $W$ is another abelian group and $h:V\to W$ preserves addition,
then $h(0)=0$: the equality
$h(0)=h(0)+h(0)$ can be translated by $-h(0)$. Also
$h(-c)=-h(c)$, because $h(c)+h(-c)=h(0)=0$. Consequently
\[
h(a-c)=h(a)-h(c),\qquad h\tau_c=\tau_{h(c)}h.
\]
For a module, multiplication by a scalar is an additive homomorphism,
so distributivity is exactly this preservation statement. In a group,
the same discussion uses multiplication and left or right translations,
with the order retained. Its multiplication is a map of underlying
sets; it is not generally a homomorphism from the product group.

There is also a structural explanation of the triangular substitutions
used in basis exchange. Given an additive homomorphism $u:V\to W$, the
\emph{shear}
\[
S_u:V\times W\longrightarrow V\times W,
\qquad (v,w)\longmapsto(v,w+u(v))
\]
has inverse $S_{-u}$. This follows by composing the maps, without counting
any basis. Over a field, if $t_1\ne0$, the map
\[
(x_1,\ldots,x_n)\longmapsto
(t_1x_1,t_2x_1+x_2,\ldots,t_nx_1+x_n)
\]
is a scaling in the first coordinate followed by the shear $S_u$ on
$\kk\times\kk^{n-1}$, where
\[
u:\kk\longrightarrow\kk^{n-1},\qquad
u(z)=(t_2t_1^{-1}z,\ldots,t_nt_1^{-1}z).
\]
Both maps are invertible. The inverse of the composite sends $y_1$ to
$t_1^{-1}y_1$ and $y_i$ to
$y_i-t_i t_1^{-1}y_1$ for $i>1$. This establishes the elementary maneuver
without presupposing invariance of dimension.

The aim of the following constructions is similar: replace a list of
rearrangements by maps whose behavior is controlled once, and then reuse
that control.

\subsection{What is retained by a universal construction?}

A \emph{category} consists of objects, sets of maps between objects,
identity maps, and associative composition. A diagram commutes when
every indicated pair of composites with the same endpoints agrees.
A map $m$ is a \emph{monomorphism} if $mu=mv$ implies $u=v$; a map $e$
is an \emph{epimorphism} if $ue=ve$ implies $u=v$. These are cancellation
conditions on maps, not definitions in terms of elements.

A \emph{subobject} of $X$ is a monomorphism into $X$, with two such
monomorphisms identified if an isomorphism over $X$ relates them. We
write $A\leq B\leq X$ when the map representing $A$ factors through the
one representing $B$. Products are characterized by their projections:
a map to $X\times Y$ is uniquely a pair of maps to $X$ and $Y$.

\begin{definition}[Pullback]
For $f:X\to Z$ and $g:Y\to Z$, a pullback is an object $P$ with maps
$a:P\to X$, $b:P\to Y$ such that $fa=gb$, universal with that property:
every pair $u:T\to X$, $v:T\to Y$ with $fu=gv$ has a unique map
$t:T\to P$ satisfying $at=u$ and $bt=v$.
\[
\begin{tikzcd}
P \arrow[r,"b"] \arrow[d,"a"'] & Y \arrow[d,"g"]\\
X \arrow[r,"f"'] & Z.
\end{tikzcd}
\]
We write $P=X\times_ZY$.
\end{definition}

In sets this is the set of pairs with the same image in $Z$. In groups
and modules it has the same description with coordinatewise operations.
The universal property retains the equation $fu=gv$, without selecting
coordinates. A terminal object receives exactly one map from each object.
An initial object has exactly one map to each object.
An equalizer of $u,v:X\to Y$ is universal among maps $a:T\to X$ with
$ua=va$. A finite-limit category has a terminal object, finite products,
and equalizers, hence all the pullbacks used here: $X\times_ZY$ is the
equalizer of $f\pi_X,g\pi_Y:X\times Y\to Z$. The
\emph{kernel pair} of $f:X\to Y$ is
$X\times_YX\rightrightarrows X$: it records exactly which inputs $f$
identifies.

A \emph{coequalizer} of $u,v:K\rightrightarrows X$ is a map $q:X\to Q$
with $qu=qv$ such that every map $h:X\to Z$ satisfying $hu=hv$ factors
uniquely as $h=\bar hq$. A \emph{regular epimorphism} is a coequalizer
of some pair. Every regular epimorphism is an epimorphism, by the
uniqueness part of this property. We draw it as $X\twoheadrightarrow Q$.

\begin{lemma}[The elementary descent rule]\label{f:descent}
In a category with pullbacks, a regular epimorphism $q:X\twoheadrightarrow Q$
is the coequalizer of its kernel pair. It has the following lifting
property against any monomorphism $m:A\hookrightarrow B$: in a
commutative square
\[
\begin{tikzcd}
X \arrow[r,"a"] \arrow[d,two heads,"q"'] & A \arrow[d,hook,"m"]\\
Q \arrow[r,"b"'] \arrow[ur,dashed,"d"] & B,
\end{tikzcd}
\]
there is a unique $d$ with $dq=a$ and $md=b$.
\end{lemma}
\begin{proof}
Choose a pair $u,v$ of which $q$ is a coequalizer. Since $qu=qv$, this
pair factors through the kernel pair of $q$. A map equal on that kernel
pair is therefore equal on $u,v$, and descends uniquely through $q$.
This proves the first claim.

For the square, $ma u=bq u=bq v=ma v$, and monicity of $m$ gives
$a u=a v$. Hence $a=dq$ for a unique $d$. Now $mdq=ma=bq$, so
$md=b$ because $q$ is epic. The equation $dq=a$ also proves uniqueness.
\end{proof}

In particular, a regular epimorphism that is also a monomorphism is an
isomorphism: apply the lifting property to its square with itself and
identity maps. This argument uses neither a choice of representatives nor
a splitting assumption.

\begin{definition}[Regular category]\label{f:regular}
A category is \emph{regular} if it has finite limits, every map factors
as a regular epimorphism followed by a monomorphism, and pullbacks of
regular epimorphisms are regular epimorphisms. The monomorphism in
\[
X\twoheadrightarrow \im(f)\hookrightarrow Y
\]
is called the \emph{image} of $f:X\to Y$.
\end{definition}

\begin{lemma}[Image calculus]\label{f:images}
In a regular category, images are unique as subobjects, regular
epimorphisms compose, and pulling back an image factorization gives an
image factorization.
\end{lemma}
\begin{proof}
Given two regular-epimorphism--monomorphism factorizations of a map,
Lemma~\ref{f:descent} gives a map between their middle objects in each
direction. Its uniqueness says the two composites are identities.
This proves uniqueness of images.

The lifting property in Lemma~\ref{f:descent} is preserved under
composition: first lift against the first regular epimorphism and then
against the second. Factor their composite as $mp$, where $p$ is a
regular epimorphism and $m$ is monic. The lifting property of the
composite, applied against $m$, gives a right inverse of $m$. A monic
map with a right inverse is an isomorphism. The composite is therefore
a regular epimorphism.

Pull back the mono part of a factorization, and then pull back its
regular epi part. The resulting square is a pullback of the composite,
by the defining universal properties. Pullback stability makes the
first part a regular epimorphism; pullbacks of monomorphisms are monic
because maps into the pullback are determined by their two projections.
This is the required image factorization.
\end{proof}

Images express what survives a map. Pullback stability says that this
description remains valid after imposing a compatible condition at the
target. It is the reason local witnesses can safely be used below.

\subsection{A precise license to reason with witnesses}

An arbitrary category may have too few maps from its terminal object to
detect equality. In particular, a group has just one group homomorphism
from the trivial group. Thus an argument about ordinary global points
would miss nearly everything. We instead use \emph{generalized elements}:
an input $x$ of $X$ over a test object $T$ means a map $x:T\to X$.

For a subobject $M\hookrightarrow X$, saying that $x$ belongs to $M$
means that $x$ factors through that monomorphism. A \emph{cover} below
means a regular epimorphism, not an assumed topological covering.

\begin{lemma}[Local witness calculus]\label{f:local-calculus}
Let $f:A\to X$ be a map in a regular category. A generalized element
$x:T\to X$ belongs to $\im(f)$ if and only if there is a cover
$e:T'\twoheadrightarrow T$ and a map $a:T'\to A$ such that $fa=xe$.
Finite lists of such witnesses can be chosen over a common cover.
They can be pulled back along any map of test objects.
\end{lemma}
\begin{proof}
Factor $f=mp$ with $p:A\twoheadrightarrow I$ regular epic and $m:I\hookrightarrow X$
monic. If $x=mi$, pull back $p$ along $i:T\to I$. The pullback cover
has a map $a$ with $fa=xe$. Conversely, such $e,a$ give a square
against $m$; Lemma~\ref{f:descent} makes $x$ factor through $m$.

Pulling back a cover gives a cover. To combine two choices over $T$,
take their fiber product over $T$; either projection is a cover, and
its composite to $T$ is a cover by Lemma~\ref{f:images}. Repeat for
finitely many choices. Nested choices are handled by composing covers.
\end{proof}

This is the precise meaning of ``there exists'' in our relation
calculations. A witness may exist only after a cover of the test object.
Every proof using witnesses below permits that replacement. Equalities
of subobjects are checked by inclusions in both directions: to check
$M\leq N$, one may test the generalized element $M\hookrightarrow X$
itself. The method therefore proves actual categorical statements,
not merely assertions about points.

\subsection{Relations compose by matching, then forgetting}

\begin{definition}[Internal relation]
A relation $R:X\rightsquigarrow Y$ is a subobject of $X\times Y$.
Choose a representative with projections
\[
\begin{tikzcd}[column sep=large]
& M \arrow[dl,"a"'] \arrow[dr,"b"] &\\
X && Y.
\end{tikzcd}
\]
The pair $(a,b)$ is jointly monic. We may write $x\,R\,y$ for membership
in this subobject. Its opposite $R^{\mathrm{op}}:Y\rightsquigarrow X$
interchanges the projections. The identity relation $\Delta_X$ is the
diagonal $X\hookrightarrow X\times X$.
\end{definition}

If $S:Y\rightsquigarrow Z$ has carrier $N$, form $M\times_YN$ and take
its image in $X\times Z$. This is the composite $SR$. Thus
\[
x\,(SR)\,z
\quad\Longleftrightarrow\quad
\text{locally there is }y\text{ with }x\,R\,y\text{ and }y\,S\,z.
\]
\emph{Our convention is always that $SR$ means first $R$, then $S$.}
The pullback retains the matching intermediate value; the image forgets
which matching value was used.

\begin{proposition}[The calculus of relations]\label{f:relation-calculus}
In a regular category, relation composition is associative and has
the diagonal relations as identities. It preserves inclusions, and
\[
(SR)^{\mathrm{op}}=R^{\mathrm{op}}S^{\mathrm{op}}.
\]
\end{proposition}
\begin{proof}
For $R:X\rightsquigarrow Y$, $S:Y\rightsquigarrow Z$, and
$T:Z\rightsquigarrow W$, membership in either $(TS)R$ or $T(SR)$
means that, over a cover, there are $y,z$ with
$x\,R\,y$, $y\,S\,z$, $z\,T\,w$. Lemma~\ref{f:local-calculus}
combines any nested witnesses into one cover and also allows the same
triple witness to be grouped in either order. This proves equality of
the two image subobjects. Equivalently, both are the image of
$M_R\times_YM_S\times_ZM_T$ in $X\times W$.

A diagonal intermediate relation forces the two matched coordinates
to agree and introduces no further condition; hence the unit laws.
An inclusion of relations sends every witness to a witness, proving
monotonicity. Reversing a witness exchanges its endpoints and reverses
the order of the two conditions, proving the last identity.
\end{proof}

This associativity has the same shape as elementary reassociation, but
it applies to correspondences that may not be functions. Its proof
identifies one object of complete witness data and shows that the order
of forgetting intermediate information does not matter.

A map $f:X\to Y$ is identified with its graph relation, represented by
$(\id_X,f):X\hookrightarrow X\times Y$. The opposite graph is generally
not a map: it remembers all possible preimages.

\begin{lemma}[Graphs, kernels, and supports]\label{f:graph-kernel}
For a map $f:X\to Y$,
\[
f^{\mathrm{op}}f=\operatorname{Eq}(f),
\qquad
ff^{\mathrm{op}}=\Delta_{\im(f)}
\quad\text{as a relation on }Y.
\]
Here $\operatorname{Eq}(f)=X\times_YX$ is the kernel pair, and the second
diagonal is supported on the image subobject of $Y$. Consequently
\[
ff^{\mathrm{op}}f=f,
\qquad ff^{\mathrm{op}}=\Delta_Y
\text{ if }f\text{ is regular epic}.
\]
Every relation with projections $a:M\to X$, $b:M\to Y$ satisfies
\[
R=ba^{\mathrm{op}}.
\]
\end{lemma}
\begin{proof}
The first composite relates $x,x'$ exactly when $f(x)=f(x')$, which is
the defining pullback. The second relates $y,y'$ exactly when, locally,
$y=f(x)=y'$ for some $x$; this is the image diagonal by
Lemma~\ref{f:local-calculus}. These descriptions give the two
consequences. Finally $ba^{\mathrm{op}}$ is the image of $(a,b)$ in
$X\times Y$; that pair already represents the subobject $R$.
\end{proof}

\begin{lemma}[When a relation is a map]\label{f:functional}
A relation $R:X\rightsquigarrow Y$ is the graph of a unique map
$f:X\to Y$ if
\[
\Delta_X\leq R^{\mathrm{op}}R,
\qquad RR^{\mathrm{op}}\leq\Delta_Y.
\]
The first condition says that every input has a locally available
output; the second says that the output is unique.
\end{lemma}
\begin{proof}
Write $R=ba^{\mathrm{op}}$ with carrier $M$. The first condition says
that the image of $a:M\to X$ is all of $X$, so $a$ is regular epic.
If $u,v:T\to M$ have $au=av$, then $bu$ and $bv$ are two outputs for
the same input. The second condition gives $bu=bv$. Joint monicity
of $(a,b)$ gives $u=v$, so $a$ is also monic and hence an isomorphism.
Now $f=ba^{-1}$ has graph $R$. Its value is forced by that graph,
which proves uniqueness.
\end{proof}

\subsection{Exactness makes an identification into an object}

A relation $E:X\rightsquigarrow X$ is \emph{reflexive} if
$\Delta_X\leq E$, \emph{symmetric} if $E^{\mathrm{op}}=E$, and
\emph{transitive} if $EE\leq E$. An \emph{internal equivalence relation}
has all three properties. Reflexivity also gives $E\leq EE$, so an
equivalence relation satisfies $EE=E$.

Kernel pairs are equivalence relations. For example, transitivity says
that $f(x)=f(y)$ and $f(y)=f(z)$ imply $f(x)=f(z)$; this is equally an
equality of maps over any test object.

\begin{definition}[Barr exactness]\label{f:exact}
A regular category is \emph{Barr-exact} if every internal equivalence
relation is the kernel pair of a map. Equivalently, for every such
$E$ on $X$ there is a regular epimorphism $q:X\twoheadrightarrow Q$
with $E=q^{\mathrm{op}}q$. We call $Q$ the quotient $X/E$.
\end{definition}

For the equivalence of these formulations, factor a map exhibiting the
kernel pair into its image. The monomorphism part does not change its
kernel pair, and the other part is regular epic. Lemma~\ref{f:descent}
then gives the quotient's full universal property: maps from $X$ that
identify the pairs in $E$ descend uniquely to $Q$. Two such quotients
are uniquely isomorphic over $X$, by applying this property twice.
Exactness concerns these equivalence-relation quotients; it does not
assert that the category is abelian or that its maps can be added.

A category is \emph{pointed} if it has an object $0$ that is both initial
and terminal. There is then a zero map between every two objects, namely
the composite through $0$. A \emph{kernel} of $f:X\to Y$ is the
pullback of $0\to Y$ along $f$:
\[
\begin{tikzcd}
\Ker(f) \arrow[r] \arrow[d,hook] & 0 \arrow[d]\\
X \arrow[r,"f"'] & Y.
\end{tikzcd}
\]
The left map is monic since $0\to Y$ is monic. In a pointed category
it is a distinguished subobject, but it is only one class of the whole
kernel pair.

If $E=q^{\mathrm{op}}q$, its zero class is the relation $E0_X$,
viewed as a subobject of $X$, where $0_X:0\to X$. Since $q0_X=0_Q$,
\[
E0_X=q^{\mathrm{op}}0_Q=\Ker(q).
\]
It is saturated: $E(E0_X)=E0_X$. We will retain both the full
equivalence relation and its zero class. In groups the zero class
determines the normal quotient, but we will not build that additional
fact into the general framework.

The roles of the assumptions are now visible. Finite limits impose
compatible equations. Images forget witnesses. Regularity makes that
forgetting compatible with substitution. Exactness realizes consistent
identifications as quotient objects. Pointedness selects a zero class.
The next condition will make different identifications compatible with
one another.

% END INPUT 01-maps-relations-exactness.tex


% BEGIN INPUT 02-malcev-and-comparison.tex
\section{The Mal'cev condition and the universal comparison}
\label{d:begin}

Suppose a relation connects inputs in $X$ with outputs in $Y$. Three
corners of a rectangle may already force the fourth:
\[
x\,R\,y,\qquad x'\,R\,y,\qquad x'\,R\,y'
\quad\Longrightarrow\quad x\,R\,y'.
\]
This property is called \emph{difunctionality}. It says that the
relation breaks into complete rectangular blocks. It does not say that
each input has just one output. Several inputs and several outputs can
belong to the same block.

The structural comparison is already suggested by this picture. Group
the inputs according to which block they occupy, and group the outputs
in the same way. The two resulting quotient objects should agree. We
will construct that agreement, including its maps, before using it to
compare any filtrations.

\subsection{Why relative difference is enough}

\begin{definition}[Mal'cev category]\label{d:malcev}
A regular category is a \emph{Mal'cev category} if every reflexive
internal relation is an equivalence relation.
\end{definition}

This categorical condition has an elementary source. Suppose we have
algebraic structures with a ternary term $p$ satisfying
\[
p(x,y,y)=x,\qquad p(x,x,y)=y.
\]
If a compatible relation $R\subseteq X\times X$ is reflexive and
$x\,R\,y$, its closure under $p$ gives
\[
p\bigl((y,y),(x,y),(x,x)\bigr)=(y,x)\in R.
\]
Thus it is symmetric. If also $y\,R\,z$, closure gives
\[
p\bigl((x,y),(y,y),(y,z)\bigr)=(x,z)\in R,
\]
so it is transitive. The operation is applied coordinatewise to pairs.
For groups $p(x,y,z)=xy^{-1}z$; for modules it is $x-y+z$. No
commutativity of group multiplication was used. A term is an operation
on the underlying algebra; in the group case it need not be a group
homomorphism from $X^3$ to $X$.

This gives the central common resource: a way to take a relative
difference and then reattach it at another point. The categorical
definition retains its effect on compatible relations without requiring
a chosen formula for the operation.

\begin{proposition}[Three forms of the Mal'cev condition]\label{d:difunctional}
For a regular category the following conditions are equivalent:
\begin{enumerate}
\item every reflexive relation is an equivalence relation;
\item equivalence relations on the same object commute under
relation composition;
\item every relation $R:X\rightsquigarrow Y$ is difunctional, that is,
\[
RR^{\mathrm{op}}R=R.
\]
\end{enumerate}
\end{proposition}
\begin{proof}
Assume the first condition. If $E,F$ are equivalence relations on $X$,
their composite $EF$ is reflexive, since both contain $\Delta_X$.
It is therefore symmetric, and
\[
EF=(EF)^{\mathrm{op}}=FE.
\]
This proves the second condition.

Assume the second. Give $R$ a carrier with projections
$a:M\to X$ and $b:M\to Y$. The kernel pairs
\[
E_a=a^{\mathrm{op}}a,\qquad E_b=b^{\mathrm{op}}b
\]
are equivalence relations on the \emph{same} object $M$, so they
commute. Using Lemma~\ref{f:graph-kernel}, we obtain
\begin{align*}
RR^{\mathrm{op}}R
&=ba^{\mathrm{op}}ab^{\mathrm{op}}ba^{\mathrm{op}}\\
&=bE_aE_ba^{\mathrm{op}}
 =bE_bE_aa^{\mathrm{op}}\\
&=(bb^{\mathrm{op}}b)(a^{\mathrm{op}}aa^{\mathrm{op}})
 =ba^{\mathrm{op}}=R.
\end{align*}
The equalities are typed on $M$; we have not interchanged unrelated
relations on $X$ and $Y$. This proves the third condition.

Finally assume the third and let $R$ be reflexive. Monotonicity gives
\[
R^{\mathrm{op}}=\Delta_XR^{\mathrm{op}}\Delta_X
\leq RR^{\mathrm{op}}R=R.
\]
Taking opposites gives the reverse inclusion, so $R$ is symmetric.
Also $\Delta_X\leq R^{\mathrm{op}}$, and hence
$RR\leq RR^{\mathrm{op}}R=R$. Thus $R$ is transitive and is an
equivalence relation.
\end{proof}

The local-witness interpretation of $RR^{\mathrm{op}}R\leq R$ is
exactly the three-corners implication above. The reverse inclusion
always holds: a related pair can supply all three edges of a degenerate
rectangle. The proposition explains why the rectangle-completion law
holds simultaneously for every relation, rather than having to be
verified anew for each quotient comparison.

\subsection{Partial identifications and their supports}

Relations used in a filtration may live on a subobject rather than all
of the ambient object. We must record that support instead of silently
assuming reflexivity everywhere.

\begin{lemma}[Partial equivalence relations]\label{d:partial-equivalence}
Let $E:X\rightsquigarrow X$ be symmetric and transitive in a regular
category. The two projections of its carrier have the same image
$m:D\hookrightarrow X$. Restricted to $D$, the relation is an
equivalence relation $E_D$. In a Barr-exact category there is a regular
epimorphism $q:D\twoheadrightarrow Q$ with
\[
E=mq^{\mathrm{op}}qm^{\mathrm{op}}.
\]
Conversely, every relation of this displayed form is symmetric and
transitive, with support $D$.
\end{lemma}
\begin{proof}
Symmetry interchanges the two projections, so their images agree.
Their factorizations through $D$ give a relation $E_D\hookrightarrow
D\times D$; the pair is still jointly monic. Both of its projections
are regular epic. If $x$ is an input in $D$, a cover supplies a witness
$x\,E_D\,y$. Symmetry gives $y\,E_D\,x$, and transitivity gives
$x\,E_D\,x$. Lemma~\ref{f:local-calculus} descends this membership,
so $E_D$ is reflexive. Symmetry and transitivity are inherited by
restriction. Exactness gives $E_D=q^{\mathrm{op}}q$, proving the formula.

Conversely, $q^{\mathrm{op}}q$ is an equivalence relation on $D$.
Its extension along $m$ is symmetric; since $m^{\mathrm{op}}m=\Delta_D$,
its square is itself. The projections onto $D$ are regular epic
because the relation is reflexive, so their images in $X$ are $D$.
\end{proof}

We call $Q$ the quotient of the partial equivalence relation. It is
the quotient of its support, not a quotient of all of $X$. This
distinction permits the same relation language to describe every step
of a chain whose terms are only normal in their successors.

For example, a quotient map $q:A\twoheadrightarrow S$ with
$A\hookrightarrow G$ determines a relation on $G$: two inputs are
related precisely when they both lie in $A$ and have the same image
in $S$. Its support is $A$; in a pointed category its zero class is
$\Ker(q)$. Thus one partial relation remembers the numerator, the
identifications, and the chosen quotient, up to its unique quotient
isomorphism.

\subsection{One relation, two quotients, one object}

Let $R:X\rightsquigarrow Y$ have carrier $M$ and factor its projections
through their images:
\[
\begin{tikzcd}[column sep=large]
&M \arrow[dl,two heads,"a"'] \arrow[dr,two heads,"b"]&\\
X' \arrow[d,hook,"m"'] &&Y' \arrow[d,hook,"n"]\\
X&&Y.
\end{tikzcd}
\]
The pair $(a,b)$ is jointly monic. We use $\widehat R:X'\rightsquigarrow Y'$
for the relation on these supports, so $R=n\widehat Rm^{\mathrm{op}}$.
Both projections of $\widehat R$ are regular epic.

\begin{theorem}[The universal support-quotient comparison]\label{d:comparison}
In a Barr-exact Mal'cev category, put
\[
E=\widehat R^{\mathrm{op}}\widehat R
\quad\text{on }X',\qquad
F=\widehat R\widehat R^{\mathrm{op}}
\quad\text{on }Y'.
\]
These are equivalence relations. For their regular quotient maps
$q:X'\twoheadrightarrow Q=X'/E$ and
$p:Y'\twoheadrightarrow P=Y'/F$, there is a unique isomorphism
$\theta:Q\iso P$ satisfying
\[
\theta qa=pb.
\]
In particular the two support quotients of any relation are canonically
isomorphic.
\end{theorem}
\begin{proof}
Because $a$ is regular epic, every input in $X'$ has locally some
$\widehat R$-related output. Using that same output twice proves
$\Delta_{X'}\leq E$. Likewise $\Delta_{Y'}\leq F$. Both relations are
symmetric by taking opposites. Difunctionality gives
\[
E^2=\widehat R^{\mathrm{op}}
 (\widehat R\widehat R^{\mathrm{op}}\widehat R)=E,
\qquad
F^2=(\widehat R\widehat R^{\mathrm{op}}\widehat R)
 \widehat R^{\mathrm{op}}=F.
\]
Thus they are equivalence relations, and exactness supplies $q,p$.

We construct the map rather than merely matching equivalence classes.
The map $pb:M\to P$ is equal on the kernel pair of $a$: if $au=av$,
then $bu$ and $bv$ share a $\widehat R$-related input, so they are
$F$-related and $pbu=pbv$. Lemma~\ref{f:descent} gives a unique
$g:X'\to P$ with $ga=pb$.

Next $g$ is equal on $E$. An $E$-related pair is locally of the form
$a(u),a(v)$ with $b(u)=b(v)$. Consequently
$ga(u)=pb(u)=pb(v)=ga(v)$. These equalities hold after a cover,
and covers are epimorphisms, so they hold on the relation itself.
Hence $g$ descends uniquely through $q$, giving $\theta:Q\to P$.

The same construction with the opposite relation gives $\psi:P\to Q$
such that $\psi pb=qa$. It follows that
\[
\psi\theta qa=qa,\qquad \theta\psi pb=pb.
\]
Both $qa$ and $pb$ are composites of regular epimorphisms, hence are
epic. Cancellation proves $\psi\theta=\id_Q$ and
$\theta\psi=\id_P$. Finally $qa$ is epic, so the equation specifying
$\theta$ proves its uniqueness.
\end{proof}

The construction is summarized by the commutative diagram
\[
\begin{tikzcd}[column sep=large,row sep=large]
&M \arrow[dl,two heads,"a"'] \arrow[dr,two heads,"b"]&\\
X' \arrow[d,two heads,"q"'] &&Y' \arrow[d,two heads,"p"]\\
Q \arrow[rr,"\theta", "\sim"'] &&P.
\end{tikzcd}
\]
There is generally no map $X'\to Y'$ to put across the middle row.
The relation supplies multiple possible outputs. The quotient removes
precisely their ambiguity, after which the comparison becomes a
uniquely determined isomorphism.

For a familiar model, take the additive group $\mathbb Z$ and the relation
$xRy$ meaning that $x-y$ is even. This is a subgroup of
$\mathbb Z\times\mathbb Z$, with both supports equal to $\mathbb Z$.
Two inputs share an output precisely when they have the same parity;
the same is true of two outputs sharing an input. Thus both support
quotients are $\mathbb Z/2\mathbb Z$, and the comparison sends the parity
class of an input to that of any related output. There are infinitely
many possible outputs before quotienting, but exactly one class afterwards.
This illustrates the construction; the categorical theorem did not need
a preferred output or a choice of representatives.

Here is another view of the same map, useful for computations. Its
graph is
\[
L=p\widehat Rq^{\mathrm{op}}:Q\rightsquigarrow P.
\]
Indeed,
\begin{align*}
L^{\mathrm{op}}L
 &=q\widehat R^{\mathrm{op}}p^{\mathrm{op}}p
       \widehat Rq^{\mathrm{op}}
   =qE^2q^{\mathrm{op}}=\Delta_Q,\\
LL^{\mathrm{op}}
 &=p\widehat Rq^{\mathrm{op}}q
       \widehat R^{\mathrm{op}}p^{\mathrm{op}}
   =pF^2p^{\mathrm{op}}=\Delta_P.
\end{align*}
We used $q^{\mathrm{op}}q=E$, $p^{\mathrm{op}}p=F$, their
idempotence, and $qq^{\mathrm{op}}=\Delta_Q$,
$pp^{\mathrm{op}}=\Delta_P$. Lemma~\ref{f:functional} makes $L$ a
map, and makes $L^{\mathrm{op}}$ its inverse. This is the same $\theta$
because its graph relates the quotient images of every pair in
$\widehat R$.

This theorem is the comparison mechanism from which the refinement
rectangle will be built. It is stronger than a statement about lengths:
it constructs the actual isomorphism between the surviving objects.
No filtration, simplicity assumption, or Jordan--H\"older theorem has
entered its proof.

\subsection{Which relations descend through a quotient?}

There is one further compatibility to make explicit. If we identify
inputs by a quotient map, a relation descends only when changing an
input within its equivalence class does not change the relation. This
is saturation.

\begin{proposition}[Saturation and descent]\label{d:saturated-descent}
Let $q:X\twoheadrightarrow Q$ and $p:Y\twoheadrightarrow P$ be
regular epimorphisms, with kernel pairs
$E=q^{\mathrm{op}}q$ and $F=p^{\mathrm{op}}p$. There is an
order-preserving bijection
\[
\left\{\text{relations }L:Q\rightsquigarrow P\right\}
\longleftrightarrow
\left\{R:X\rightsquigarrow Y\mid FRE=R\right\}.
\]
It sends $L$ to $p^{\mathrm{op}}Lq$, and its inverse sends $R$ to
$pRq^{\mathrm{op}}$. The first construction is precisely inverse image
of the relation along $q\times p$.
\end{proposition}
\begin{proof}
Inverse image consists of pairs $(x,y)$ whose quotient images satisfy
$q(x)\,L\,p(y)$, which is $p^{\mathrm{op}}Lq$ by the graph formulas.
Since $qq^{\mathrm{op}}=\Delta_Q$ and
$pp^{\mathrm{op}}=\Delta_P$, the two constructions satisfy
\[
p(p^{\mathrm{op}}Lq)q^{\mathrm{op}}=L,
\qquad
p^{\mathrm{op}}(pRq^{\mathrm{op}})q=FRE.
\]
An inverse-image relation is saturated because
$Fp^{\mathrm{op}}=p^{\mathrm{op}}$ and $qE=q$. Thus the formulas are
mutually inverse exactly on the stated saturated relations. Both are
order preserving by monotonicity of relation composition.
\end{proof}

For a relation $T$ on $X$ and a single quotient $q$, saturation is
$ETE=T$. Since $E$ contains the diagonal, this implies separately
$ET=T$ and $TE=T$. Consequently the correspondence respects
composition of saturated relations:
\[
(qSq^{\mathrm{op}})(qTq^{\mathrm{op}})
=qSETq^{\mathrm{op}}=qSTq^{\mathrm{op}}.
\]
It also respects opposites. Thus a saturated symmetric transitive
relation descends to a symmetric transitive relation, including its
possibly proper support. Conversely, inverse image preserves these
properties: for example
\[
(q^{\mathrm{op}}Lq)^2
=q^{\mathrm{op}}L^2q\leq q^{\mathrm{op}}Lq
\]
when $L$ is transitive. These statements hold for partial equivalence
relations as well as for equivalence relations on all of $X$.

\begin{corollary}[Quotients of quotients]\label{d:iterated-quotient}
In a Barr-exact category, fix $q:X\twoheadrightarrow Q$ with kernel
pair $E$. Equivalence relations $L$ on $Q$ correspond to equivalence
relations $T$ on $X$ containing $E$, by
\[
T=q^{\mathrm{op}}Lq,\qquad L=qTq^{\mathrm{op}}.
\]
If $t:X\twoheadrightarrow Z$ is the quotient of $T$, it factors
uniquely as $t=uq$, where $u:Q\twoheadrightarrow Z$ is the quotient
of $L$. Thus $(X/E)/L\cong X/T$, with the isomorphism determined by
the quotient maps.
\end{corollary}
\begin{proof}
If $T$ is an equivalence relation containing $E$, then
$T\leq ETE\leq TTT=T$, so it is saturated. Its descent is reflexive
because $q\Delta_Xq^{\mathrm{op}}=\Delta_Q$, and is symmetric and
transitive by the preceding formulas. Conversely, the inverse image
of an equivalence relation is an equivalence relation containing
$q^{\mathrm{op}}\Delta_Qq=E$.

Since $E\leq T$, the map $t$ is equal on the kernel pair of $q$,
so it descends uniquely to $u$. Factor $u$ into its image. If its mono
part were a proper subobject of $Z$, then $t=uq$ would factor through
that same subobject, contradicting the lifting property of the regular
epimorphism $t$. Hence $u$ is regular epic. Its kernel pair is
\[
u^{\mathrm{op}}u
=qq^{\mathrm{op}}u^{\mathrm{op}}uqq^{\mathrm{op}}
=qt^{\mathrm{op}}tq^{\mathrm{op}}=qTq^{\mathrm{op}}=L.
\]
Lemma~\ref{f:descent} makes it the quotient of $L$. The universal
property proves the claimed canonical isomorphism.
\end{proof}

This last result needs exactness and regularity, but not the Mal'cev
condition. Recording that distinction matters: Mal'cev is used to make
\emph{every} relation difunctional and hence to produce the universal
comparison. The ordinary law for compatible iterated quotients has
already been supplied by exactness. Later, in a higher setting, we
will see why retaining such compatibility data remains essential.

The relation calculus and support comparison are classical. For
context, see \href{https://arxiv.org/pdf/math/0610552}{Knop,
\emph{Tensor envelopes of regular categories}}, Section~5 and
Appendices~A--B. The arguments here have been written out so that the
refinement theorem can depend on explicit constructions rather than on
an unnamed isomorphism theorem.

% END INPUT 02-malcev-and-comparison.tex


% BEGIN INPUT 03-refinement-and-factors.tex
\section{A quotient remembers which identifications were made}
\label{r:chapter}

We now have a language in which a quotient is an operation on a relation,
and a comparison of quotients is an operation on a correspondence.  We will
use that language once, in a pointed Barr-exact Mal'cev category
$\mathcal C$.  Groups and modules will enter afterwards by verifying these
axioms.  Neither commutative addition of arrows nor an invariant called
dimension is used in this chapter.

There are three stages.  First we record a normal series without forgetting
its quotient maps.  Next we compare two such records by a rectangular array
of correspondences.  Finally simplicity prevents the comparison from
breaking a factor into smaller nonzero factors.  The last step is
Jordan--H\"older; the rectangle, rather than a numerical count, does the work.

\subsection{The data in a normal series}

\begin{definition}[Normal series with its quotient maps]
\label{r:normal-series}
A finite normal series of $X$ consists of a chain of subobjects and chosen
regular epimorphisms
\[
  0=X_0\lhook\joinrel\longrightarrow X_1
  \lhook\joinrel\longrightarrow\cdots
  \lhook\joinrel\longrightarrow X_n=X,
  \qquad q_i:X_i\twoheadrightarrow S_i,
\]
such that the inclusion $X_{i-1}\hookrightarrow X_i$ is a kernel of $q_i$.
The objects $S_i$ are its \emph{factors}.  Repetitions of subobjects and zero
factors are permitted.  Two series are \emph{equivalent} when their lists of
factors are isomorphic after a permutation.  We also use equivalence
\emph{after deleting zero factors}, stating explicitly when we do so.
\end{definition}

The kernel square in this definition is
\[
\begin{tikzcd}[column sep=large,row sep=large]
X_{i-1}\arrow[r,hook]\arrow[d] & X_i\arrow[d,two heads,"q_i"]\\
0\arrow[r] & S_i .
\end{tikzcd}
\tag{\ref{r:normal-series}a}
\]
It is a pullback.  In particular, the normality assertion is relative to
$X_i$, not to all of $X$.  For a group this allows a subgroup normal in its
successor without demanding that it be normal in the ambient group.

We keep $q_i$ because the kernel alone need not determine a quotient in
this generality.  A familiar implication from groups, such as ``a regular
epimorphism with zero kernel is an isomorphism'', is not inserted as an
unstated axiom.  All our comparisons retain the kernel pair of the actual
map, which is the full record of its identifications.

Here and below, composition of relations is function-like: $SR$ means
first $R$, then $S$.  For a relation $r:X\rightsquigarrow X$ and a subobject
$U\hookrightarrow X$, write $rU$ for the image of the restriction of $r$
whose input lies in $U$.  Thus $rX$ is its output support and $r0$ its zero
class.  A symmetric transitive relation is a \emph{partial equivalence
relation}; it is an equivalence relation on its support, by
Lemma~\ref{d:partial-equivalence}.  Its support quotient will be written
\[
                    Q(r)=(rX)/r.
\]
The slash is shorthand for the effective quotient of the restricted
equivalence relation, not a new operation on objects without relations.

\begin{lemma}[A series as a sequence of identifications]
\label{r:series-relations}
For $n\geq1$, the data of Definition~\ref{r:normal-series}, up to the unique isomorphisms
allowed for quotient objects, are equivalently a list of partial equivalence
relations $r_1,\ldots,r_n$ on $X$ satisfying
\begin{equation}
\label{r:boundary}
 r_1 0=0,\qquad r_iX=r_{i+1}0\ (1\leq i<n),\qquad r_nX=X.
\end{equation}
The associated factors are $Q(r_i)$. The empty series of $0$ is treated
separately; it has no endpoint relations to name.
\end{lemma}

\begin{proof}
Let $a_i:X_i\hookrightarrow X$ be the composite inclusion.  Regarding
arrows as their graphs, put
\[
       r_i=a_iq_i^{\mathrm{op}}q_i a_i^{\mathrm{op}}.
\]
On $X_i$, this is the kernel pair of $q_i$; it relates precisely the two
inputs that $q_i$ identifies.  Extending it along the monomorphism $a_i$
makes it a partial equivalence relation on $X$.  Its support is $X_i$.
The zero arrow into $X_i$ is sent to the zero arrow into $S_i$, so taking
the zero class of that kernel pair gives precisely the pullback defining
$\ker q_i$.  Therefore
\[
                 r_iX=X_i,\qquad r_i0=X_{i-1}.
\]
These equalities give~\eqref{r:boundary}.  Since $q_i$ is a regular
epimorphism, it is the effective quotient of its kernel pair, and
$Q(r_i)\cong S_i$ canonically.

Conversely, set $X_i=r_iX$.  The zero class $r_i0$ is a subobject of $r_iX$.
Thus the identities~\eqref{r:boundary} make the $X_i$ a chain, beginning
with $X_0=0$ and ending with $X_n=X$.  By exactness the restricted relation
$r_i$ has an effective quotient $q_i:X_i\twoheadrightarrow Q(r_i)$.
The same kernel-pair calculation says that $\ker q_i=r_i0=X_{i-1}$.
These two constructions are inverse because an effective quotient is
determined, up to a unique compatible isomorphism, by its kernel pair.
\end{proof}

This change of language has retained more than the chain.  It retains
which pairs of points become equal at each step.  That is the information
needed to compare two quotient constructions in a category where the zero
class alone may not determine the entire relation.

\subsection{What a genuine refinement means}

To refine a series, we refine one of its factors and pull that refinement
back.  Suppose $q:D\twoheadrightarrow Q$ is one quotient step, with kernel
$K\hookrightarrow D$, and suppose
\[
 0=Q_0\hookrightarrow Q_1\hookrightarrow\cdots\hookrightarrow Q_t=Q,
 \qquad e_j:Q_j\twoheadrightarrow Z_j,\quad\ker e_j=Q_{j-1},
\]
is a normal series of its factor.  Form $P_j=D\times_Q Q_j$.  We obtain
\[
\begin{tikzcd}[column sep=large,row sep=large]
K=P_0\arrow[r,hook]\arrow[d,two heads]
 &P_1\arrow[r,hook]\arrow[d,two heads,"p_1"]
 &\cdots\arrow[r,hook]&P_t=D\arrow[d,two heads,"q"]\\
0=Q_0\arrow[r,hook]&Q_1\arrow[r,hook]
 &\cdots\arrow[r,hook]&Q_t=Q .
\end{tikzcd}
\]
Each displayed comparison square is a pullback.  Pullback stability of
regular epimorphisms makes $p_j:P_j\twoheadrightarrow Q_j$ a regular
epimorphism.  The refined quotient map at the $j$th step is
\begin{equation}
\label{r:refined-map}
                  e_jp_j:P_j\twoheadrightarrow Z_j.
\end{equation}
Its kernel is $P_{j-1}$: its pullback over $0$ is
$P_j\times_{Q_j}Q_{j-1}=D\times_Q Q_{j-1}$.  Regular epimorphisms are
closed under composition, so this really is a normal series between $K$
and $D$.  Replacing any or all of the original steps in this way is what
we mean by a \emph{refinement}.
For the empty series of the zero object we adopt the harmless convention
that inserting a finite list of zero steps is a refinement.  This
convention concerns only how empty lists are represented; it introduces
no nonzero factor.

It would be insufficient to insert arbitrary kernel subobjects between
$K$ and $D$ while forgetting their relation to $q$.  We need the new
quotients to have arisen from the original quotient.  The following
saturation condition expresses exactly that requirement.

\begin{proposition}[Saturation is descent of a refinement]
\label{r:refinement-descent}
Let $r$ be a partial equivalence relation on $X$, let $D=rX$, and let
$q:D\twoheadrightarrow Q(r)$ be its quotient.  For $k\geq1$, lists of
partial equivalence relations $t_1,\ldots,t_k$ on $X$ satisfying
\begin{align}
 rt_jr&=t_j \quad(1\leq j\leq k), \label{r:saturation}\\
 t_1 0&=r0,\qquad t_jX=t_{j+1}0\ (1\leq j<k),\qquad
 t_kX=rX \label{r:refinement-boundary}
\end{align}
correspond to normal series of $Q(r)$ with exactly $k$ steps. Under this correspondence, their
support quotients $Q(t_j)$ are the factors of the pulled-back refinement
of $q$.
\end{proposition}

\begin{proof}
The equality $rt_jr=t_j$ first implies that the support of $t_j$ lies
inside $D$, so we can work on $D$.  Write $E=q^{\mathrm{op}}q$ for the
kernel pair on $D$; it is the restriction of $r$.  The saturated descent
lemma, Proposition~\ref{d:saturated-descent}, gives mutually inverse
operations
\[
       T_j=q\,t_j\,q^{\mathrm{op}},\qquad
       t_j=q^{\mathrm{op}}T_jq.
\]
The first is the direct image relation on $Q(r)$; the second is its
pullback along $q\times q$.  Their being inverse uses $qq^{\mathrm{op}}
=1_{Q(r)}$ and $Et_jE=t_j$.  That lemma also proves that symmetry and
transitivity descend, so each $T_j$ is a partial equivalence relation.

We check the endpoints, since those checks are what make this a descent
of \emph{series}.  Saturation and reflexivity of $E$ on $D$ give
$t_jE=t_j=Et_j$.  The inverse image of the zero subobject of $Q(r)$ is
$E0$.  Therefore, as subobjects obtained by regular images,
\[
             T_j0=q\,t_jE0=q\,t_j0,
       \qquad T_jQ(r)=q(t_jD).
\]
Now $q(E0)=0$ and $qD=Q(r)$.  Applying these formulas to
\eqref{r:refinement-boundary} gives
\[
 T_10=0,\qquad T_jQ(r)=T_{j+1}0,\qquad T_kQ(r)=Q(r).
\]
Lemma~\ref{r:series-relations} therefore supplies a normal series of
$Q(r)$.

In the other direction, begin with a normal series of $Q(r)$ and its
relations $T_j$.  Pulling back the relations gives $t_j=q^{\mathrm{op}}T_jq$,
which are saturated and have precisely the support and zero-class
boundaries in~\eqref{r:refinement-boundary}.  Concretely, if $Q_j$ is the
support of $T_j$, the support of $t_j$ is $P_j=D\times_{Q(r)}Q_j$.
If $e_j:Q_j\twoheadrightarrow Z_j$ is its factor quotient, the kernel pair
of $e_jp_j$ on $P_j$ is $t_j$.  As $e_jp_j$ is a regular epimorphism,
its codomain $Z_j$ is the effective quotient of $t_j$.  This gives the
unique compatible isomorphism $Q(t_j)\cong Z_j$ and identifies the
quotient maps as well as the objects.
\end{proof}

The proposition is a small descent theorem.  The equality $rt_jr=t_j$
says that changing either input by an identification already imposed by
$r$ does not change the new relation.  Exactly then the relation can be
defined downstairs, independently of the choice of a representative.
The proof above replaces that verbal test by pullbacks, images, and
unique quotient maps.

\subsection{The comparison rectangle}

Let $r,s$ be partial equivalence relations on the same object.  Their
composite $rs$ need not be an equivalence relation on all of $X$.  It is
nevertheless a relation, so the Mal'cev theorem applies to it.  This
produces the identities that will make our rectangle fit together.

\begin{lemma}[Two endpoint identities]
\label{r:endpoint-identities}
For partial equivalence relations $r,s$ on $X$,
\begin{equation}
\label{r:rs-idempotent}
                  rsrs=rs,
\end{equation}
and
\begin{equation}
\label{r:endpoint-equalities}
                  rsr0=rs0,\qquad rsrX=rsX.
\end{equation}
\end{lemma}

\begin{proof}
The relation $P=rs$ is difunctional by Proposition~\ref{d:difunctional}, so
$PP^{\mathrm{op}}P=P$.  Symmetry and idempotence of $r,s$ reduce the left
side to
\[
           rssrrs=rsrs,
\]
which proves~\eqref{r:rs-idempotent}.

Every subobject in a pointed category contains its zero arrow.  In
particular $0\subseteq r0$ and $0\subseteq s0$.  Monotonicity of relational
image and~\eqref{r:rs-idempotent} now give
\[
 rs0\subseteq rsr0\subseteq rsrs0=rs0.
\]
Here the second inclusion comes from $r0\subseteq rs0$, itself obtained
by applying $r$ to $0\subseteq s0$.  For supports, $sX\subseteq X$ gives
\[
        rsX=rsrsX\subseteq rsrX\subseteq rsX;
\]
the final inclusion uses $rX\subseteq X$.  Both chains give equality.
\end{proof}

Now take two normal series of $X$ of positive lengths $m,n$, encoded by
$r_1,\ldots,r_m$ and $s_1,\ldots,s_n$.  At the crossing of the $i$th row
and $j$th column, define
\begin{equation}
\label{r:cells}
       r_{ij}=r_is_jr_i,\qquad
       s_{ji}=s_jr_is_j,\qquad P_{ij}=r_is_j.
\end{equation}
The two cell relations are the two equivalence relations associated with
the correspondence $P_{ij}$:
\[
       P_{ij}P_{ij}^{\mathrm{op}}=r_{ij},\qquad
       P_{ij}^{\mathrm{op}}P_{ij}=s_{ji}.
\]
Our composition convention matters here: $s_{ji}$ is on the input support
of $P_{ij}$, and $r_{ij}$ on its output support.

\begin{proposition}[Every row and column is a genuine refinement]
\label{r:rectangle-refines}
For fixed $i$, the list $r_{i1},\ldots,r_{in}$ satisfies
Proposition~\ref{r:refinement-descent} for $r_i$.  For fixed $j$, the list
$s_{j1},\ldots,s_{jm}$ satisfies that proposition for $s_j$.
\end{proposition}

\begin{proof}
We verify every condition for a row.  Symmetry follows by reversing the
word $r_is_jr_i$.  Idempotence follows from
Lemma~\ref{r:endpoint-identities}:
\[
 r_{ij}r_{ij}=r_is_jr_is_jr_i=r_is_jr_i=r_{ij}.
\]
Thus each cell is symmetric and transitive.  Its saturation is
\[
       r_i r_{ij}r_i=r_i^2s_jr_i^2=r_{ij}.
\]
For the first zero class, use $s_10=0$:
\[
            r_{i1}0=r_is_1r_i0=r_is_10=r_i0.
\]
For an intermediate join, use $s_j0=s_{j-1}X$:
\[
\begin{split}
 r_{ij}0
 &=r_is_jr_i0=r_is_j0
  =r_is_{j-1}X\\
 &=r_is_{j-1}r_iX=r_{i,j-1}X
                  \qquad(2\leq j\leq n).
\end{split}
\]
For the last support, use $s_nX=X$:
\[
                 r_{in}X=r_is_nr_iX=r_is_nX=r_iX.
\]
Every equality that removes the trailing $r_i$ is one of
\eqref{r:endpoint-equalities}.  We have proved saturation and all boundary
conditions, so Proposition~\ref{r:refinement-descent} supplies a
pulled-back refinement of the original quotient.  Interchanging $r$ and
$s$ proves the column assertion.
\end{proof}

The boundary identities are the gluing conditions.  They say that the
support of one cell is precisely the zero class required to start the
next cell.  Without these identities we would have an assortment of
quotients, not a refinement of a filtration.

\begin{proposition}[The two readings of a cell have the same quotient]
\label{r:cell-isomorphism}
For every $(i,j)$ there is a canonical isomorphism
\begin{equation}
\label{r:cell-iso}
                       Q(s_{ji})\xrightarrow{\ \sim\ }Q(r_{ij}).
\end{equation}
It is induced by the correspondence $P_{ij}$ and characterized by its
two projections, not chosen by comparing lengths or other invariants.
\end{proposition}

\begin{proof}
Let $C_{ij}$ be the object representing $P_{ij}\hookrightarrow X\times X$.
Factor its two projections through their image supports:
\[
 C_{ij}\xrightarrow{a_{ij}}L_{ij}\hookrightarrow X,
 \qquad C_{ij}\xrightarrow{b_{ij}}R_{ij}\hookrightarrow X.
\]
Both $a_{ij}$ and $b_{ij}$ are regular epimorphisms.  The input relation
$P_{ij}^{\mathrm{op}}P_{ij}$ is $s_{ji}$, so $L_{ij}=s_{ji}X$; similarly
$R_{ij}=r_{ij}X$.  Write their effective quotient maps as
$\lambda_{ij}:L_{ij}\twoheadrightarrow Q(s_{ji})$ and
$\rho_{ij}:R_{ij}\twoheadrightarrow Q(r_{ij})$.

The support-quotient comparison theorem, Theorem~\ref{d:comparison},
applies to $P_{ij}$ and gives the unique isomorphism $\theta_{ij}$ for
which
\[
\begin{tikzcd}[column sep=large,row sep=large]
L_{ij}\arrow[d,two heads,"\lambda_{ij}"']
 &C_{ij}\arrow[l,two heads,"a_{ij}"']
          \arrow[r,two heads,"b_{ij}"]
 &R_{ij}\arrow[d,two heads,"\rho_{ij}"]\\
Q(s_{ji})\arrow[rr,"\theta_{ij}"',"\sim"]
 &&Q(r_{ij})
\end{tikzcd}
\]
commutes, meaning
$\theta_{ij}\lambda_{ij}a_{ij}=\rho_{ij}b_{ij}$.
The cited theorem proves existence and invertibility by effective
quotients and difunctionality.  Its uniqueness also follows directly
because $\lambda_{ij}a_{ij}$ is an epimorphism.  The displayed computation
of the two kernel relations identifies its conclusion with
\eqref{r:cell-iso}.
\end{proof}

This is the categorical butterfly.  It is the same correspondence read
from its two ends, with each end quotiented by precisely the ambiguity
introduced by that correspondence.  No picture of four wings is needed
to remember the construction: form $P$, compute $P^{\mathrm{op}}P$ and
$PP^{\mathrm{op}}$, and take their support quotients.

\begin{theorem}[Schreier refinement in an exact Mal'cev category]
\label{r:schreier}
\label{j:schreier}
In a pointed Barr-exact Mal'cev category, any two finite normal series of
an object have equivalent refinements in the pullback sense specified
above.  The refinements may contain repeated subobjects and zero factors.
\end{theorem}

\begin{proof}
Suppose first that both series have positive numbers $m,n$ of steps.
For the first series, concatenate the $m$ rows supplied by
Proposition~\ref{r:rectangle-refines}.  Their endpoints are the original
successive subobjects, so they fit together into a normal series of $X$.
For the second series, concatenate the $n$ columns.  Both refined lists
have $mn$ factors.  Their indexing orders differ, but the transposition
$(i,j)\leftrightarrow(j,i)$ pairs them.  Proposition~\ref{r:cell-isomorphism}
gives the isomorphism for each paired cell.

An empty series can occur only for $X=0$.  Every subobject of $0$ is
isomorphic to $0$: a monomorphism $A\to0$ has its unique map $0\to A$ as
a right inverse, and monicity makes it a left inverse as well.  A regular
epimorphism out of $0$ also has zero codomain, since it is at once monic
and a regular epimorphism.  Thus every factor of every series of $0$ is
zero.  We may insert zero steps to give both series the same positive
number of zero factors, or treat two empty lists as already equivalent.
These are refinements by the zero object's zero series.  This resolves
the degenerate case without assuming a nonempty comparison rectangle.
\end{proof}

The theorem does not produce a single nested chain containing the terms
of both original chains.  Two different lines in a two-dimensional vector
space already prevent that: neither line contains the other.  It produces
one refinement of each chain and a matching between their quotient
factors.  This distinction matters even in the most elementary example.

\subsection{Simplicity removes the remaining freedom}

\begin{definition}[Quotient-simple object]
\label{j:simple}
A nonzero object $S$ is \emph{quotient-simple} if every regular
epimorphism $e:S\twoheadrightarrow T$ is either an isomorphism or has
zero codomain.  A composition series is a finite normal series all of
whose factors are quotient-simple.
\end{definition}

The condition concerns the arrow $e$.  The weaker phrase ``its codomain
is isomorphic to $S$ or to $0$'' would not exclude a proper self-quotient.
The definition above does exclude it.  In the group and module instances,
it agrees with the usual definition of a nontrivial simple object.

\begin{lemma}[A simple factor cannot be broken into smaller factors]
\label{j:simple-refinement}
Every finite normal series of a quotient-simple object $S$ has exactly
one nonzero factor, isomorphic to $S$.  Its other factors are zero.
\end{lemma}

\begin{proof}
The last quotient is a regular epimorphism $e:S\twoheadrightarrow T$.
If $T=0$, its kernel is all of $S$, so the preceding subobject is $S$.
Remove this last, zero-factor step and consider the shorter series.
Because the series is finite and begins at $0\neq S$, this removal
process must eventually reach a last step whose preceding subobject is
proper in $S$.

For that step the quotient cannot have zero codomain: a map to $0$ has
kernel $S$.  Simplicity therefore makes the quotient map an isomorphism.
Its kernel is $0$, so every preceding subobject of the chain is $0$.
Every regular quotient of $0$ is $0$, as checked in the proof of
Theorem~\ref{r:schreier}.  We have one quotient isomorphic to $S$, zero
quotients before it, and the zero quotients stripped from its end.
\end{proof}

\begin{theorem}[Jordan--H\"older]
\label{j:jordan-holder}
Let $\mathcal C$ be a pointed Barr-exact Mal'cev category.  If an object
$X$ has two finite composition series, then their factor lists agree,
including multiplicity, up to permutation and isomorphism.  In particular,
the two series have the same number of factors.
\end{theorem}

\begin{proof}
Apply Theorem~\ref{r:schreier} to obtain equivalent refinements.  Each
refinement is obtained from normal series of the original factors, pulled
back along the original quotient maps.  By Lemma~\ref{j:simple-refinement},
each such factor series contributes exactly its original simple factor
and some zero factors.  Proposition~\ref{r:refinement-descent} identifies
these with the factors of the pulled-back refinement.

Consequently, deleting zero factors from the first refined list recovers
the first original list, and deleting them from the second recovers the
second original list.  An isomorphism cannot identify a nonzero object
with $0$, so the matching between refined lists restricts to a matching
between their nonzero lists.  This is the desired permutation and set of
factor isomorphisms.  The equality of lengths is a consequence of that
bijection, not an input to its construction.
\end{proof}

No existence assertion is hidden in this theorem.  We assumed two finite
composition series and compared them.  A category may contain objects
with no finite composition series.  Even when maximal or minimal
subobjects can be selected, a termination argument is still required
before one can assert existence.  The finite sequences used here are
already supplied, so neither an ascending nor a descending chain
condition enters the uniqueness proof.

\subsection{Factor data becomes an invariant only after the comparison}

Now, and only now, a numerical or algebraic record of the factors becomes
legitimate.  Call an object \emph{composition-finite} when it has a finite
composition series, allowing $0$ with its empty series.  Write $\mu(X)$
for the finite multiset of isomorphism classes of its simple factors.
Theorem~\ref{j:jordan-holder} proves that $\mu(X)$ is independent of the
chosen series.  It does not reconstruct $X$: distinct extension data can
give the same multiset.

\begin{proposition}[Splicing an extension]
\label{j:extension-additivity}
\label{j:splicing}
Suppose $e:E\twoheadrightarrow Q$ is a regular epimorphism with kernel
$K\hookrightarrow E$.  If $K$ and $Q$ are composition-finite, then so is
$E$, and
\begin{equation}
\label{j:mu-additive}
                         \mu(E)=\mu(K)+\mu(Q).
\end{equation}
Here $+$ means disjoint union of finite multisets, including multiplicity.
\end{proposition}

\begin{proof}
Choose composition series for $K$ and $Q$.  Pull the series of $Q$ back
along $e$.  The construction~\eqref{r:refined-map} gives a normal series
from $K$ to $E$ whose factors are exactly those of $Q$.  Prepend the
chosen composition series from $0$ to $K$.  Every factor in the resulting
series is simple, so it is a composition series of $E$.  Its factor list
is the concatenation of the two chosen lists.  Jordan--H\"older makes
that calculation independent of all the choices and proves
\eqref{j:mu-additive}.
\end{proof}

This proposition establishes closure under these extensions.  It does
not assert, without an additional argument, that every subobject or every
regular quotient of a composition-finite object is composition-finite in
our weakest categorical setting.  The direction used in the proof is
precisely the direction stated.

For the next statement assume that the isomorphism classes of simple
objects under discussion form a set $\Sigma$, for example by working in
an essentially small category.  The free commutative monoid
$\mathbb N^{(\Sigma)}$ consists of finitely supported functions
$\Sigma\to\mathbb N$, with pointwise addition.  Thus $\mu(X)$ is an
element of this monoid.

\begin{corollary}[The universal additive record]
\label{j:universal-invariant}
\label{j:universal-monoid}
Let $M$ be a commutative monoid.  An arbitrary assignment $w:\Sigma\to M$
extends uniquely to an isomorphism invariant $v$ of composition-finite
objects satisfying
\[
 v(0)=0,\qquad v(E)=v(K)+v(Q)
\]
for every kernel--regular-quotient sequence $K\hookrightarrow
E\twoheadrightarrow Q$ whose terms are composition-finite.  The formula is
\[
                 v(X)=\sum_{[S]\in\Sigma}\mu(X)_{[S]}\,w([S]).
\]
\end{corollary}

\begin{proof}
Only finitely many summands are nonzero.  Jordan--H\"older makes the
formula well defined, and Proposition~\ref{j:extension-additivity} gives
the required additivity.  For uniqueness, every intermediate term of a
composition series is composition-finite, because its initial segment
is itself a composition series.  Repeatedly applying additivity to these
successive quotient maps forces $v(X)$ to equal the displayed sum.
\end{proof}

Taking every $w([S])=1$ defines length.  Taking values in a group gives
the corresponding universal group of additive invariants.  These are
consequences of factor comparison; using them to justify that comparison
would reverse the dependency and create circular reasoning.  The order
here is deliberate: quotient construction, refinement, factor matching,
then invariant.

\subsection{What is natural in the construction}

\begin{proposition}[Transport by a regular functor]
\label{r:functoriality}
Let $F:\mathcal C\to\mathcal D$ be a functor between pointed Barr-exact
Mal'cev categories that preserves finite limits, regular epimorphisms,
and a zero object.  Then $F$ transports the relation rectangle, its
refinements, and its cell isomorphisms to the corresponding constructions
in $\mathcal D$, up to their canonical structural isomorphisms.
\end{proposition}

\begin{proof}
Finite-limit preservation gives preservation of monomorphisms, products,
kernel pairs, kernels, and pullbacks.  If a map is factored as a regular
epimorphism followed by a monomorphism, its image under $F$ is again
such a factorization.  Uniqueness of image factorizations therefore shows
that $F$ preserves images.  Relational composition is an image of a
pullback projection, so $F(SR)$ identifies with $F(S)F(R)$.  Converse
relations are preserved by preservation of the product symmetry.

If $q$ is the effective quotient of an equivalence relation $E$, then
$F(q)$ is a regular epimorphism and its kernel pair is $F(E)$.  Hence it
is the effective quotient of $F(E)$.  These observations preserve every
operation in~\eqref{r:cells}, the support and zero-class equations, and
the pullback squares defining refinements.  Finally the image of the
commuting diagram in Proposition~\ref{r:cell-isomorphism} satisfies the
same characterization as the cell comparison in $\mathcal D$.  Its
uniqueness identifies the two isomorphisms.
\end{proof}

A regular functor need not preserve simple objects.  For instance,
forgetting a module's scalar action can reveal additive subgroups that
were not submodules.  The proposition transports the comparison
construction; a statement about composition series additionally requires
that the relevant simple factors remain nonzero and quotient-simple.
This distinction prevents a change of framework from silently forgetting
part of the theorem's information.

\subsection{Dependency and scope check}

The ordinary categorical proof is now complete.  The foundation needed
finite limits and regular images for relation composition, exactness for
the effective quotients, and the Mal'cev condition for difunctionality.
Pointedness supplied zero classes and kernel boundaries.  Saturated
descent made each row a refinement of the \emph{chosen} quotient map.
The support-quotient comparison matched each cell.  Simplicity then
removed the zero cells.  No invariant length, dimension, Grothendieck
group, or pre-existing Jordan--H\"older theorem entered the chain of
dependencies.

The comparison pattern is the classical categorical refinement argument
associated with Lambek; Knop's Appendix~B2 gives its exact Mal'cev
formulation and emphasizes the saturation requirement.  The proof here
has supplied the descent, boundary, and quotient-map checks explicitly;
it is not a claim of a new theorem.  The detailed primary reference is
\href{https://arxiv.org/pdf/math/0610552}{Friedrich Knop, \emph{Tensor
envelopes of regular categories}, Appendix B2}.

What has become general is a construction with stated hypotheses, not
the assertion that every category has composition factors.  The later
homotopical chapters will ask which quotient and refinement data survive
when equalities become coherent identifications.  The present theorem
is available there whenever its ordinary exact Mal'cev hypotheses are
actually verified; it does not, by terminology alone, assert a
Jordan--H\"older theorem for arbitrary higher objects.

% END INPUT 03-refinement-and-factors.tex

\clearpage\part{Instances and comparison pages}

% BEGIN INPUT 04-instances-and-dimension.tex
\section{Groups, modules, and dimension as instances}\label{e:section}

The comparison theorem has been proved before choosing a multiplication
table, a scalar field, or a basis. To use it, we must now check its hypotheses.
This is an essential part of the argument: a category name does not by itself
verify an image factorization or construct a quotient.

\subsection{One verification for algebraic structures}\label{e:varieties}

An algebraic theory specifies operations of finite arity and equations between
expressions built from them. A model is a set with those operations satisfying
the equations; a morphism preserves every operation. A subalgebra is a subset
closed under them. An internal relation on a model $A$ is a subalgebra of
$A\times A$. An equivalence relation of this kind is a \emph{congruence}:
equivalent inputs give equivalent outputs for every operation.

\begin{proposition}\label{e:algebraic-instances}
The category of models of an algebraic theory is Barr-exact. If its one-element
model is also initial, it is pointed. If the theory has a ternary term $p$
satisfying
\[
p(x,y,y)=x,\qquad p(x,x,z)=z,
\]
then it is a Mal'cev category. Under these last two conditions the general
comparison theorem applies to its finite composition series.
\end{proposition}
\begin{proof}
Products carry the operations coordinatewise. An equalizer is the subset
where two homomorphisms agree, closed under every operation. These give finite
limits, with the one-element algebra terminal. The image of a homomorphism
is a subalgebra, so every map factors as a surjection followed by an inclusion.
A pullback of a surjection is surjective: an element of the other object has
a lift of its image, giving a pair in the pullback.

Let $q:A\to B$ be surjective. A homomorphism $h:A\to C$ constant on its
fibers defines $\bar h(b)=h(a)$ for any $a$ with $q(a)=b$. This is independent
of $a$, and choosing lifts of the finitely many inputs of an operation proves
that $\bar h$ preserves that operation. Surjectivity proves uniqueness.
Thus $q$ is the coequalizer of its kernel pair.

Conversely a regular epimorphism is surjective. Factor a coequalizer
$q:A\to B$ through its set-theoretic image as $q=me$. If $q$ coequalizes
$u,v$, injectivity of $m$ shows $eu=ev$. Its universal property supplies
$d:B\to\im(q)$ with $dq=e$. Then $mdq=q$, so $md=\id_B$ because
coequalizers are epimorphisms. Also $dme=e$, and the surjectivity of $e$
gives $dm=\id_{\im(q)}$. Hence $m$ is an isomorphism.
We have established regularity with exactly the stated quotient maps.
Monomorphisms are precisely the injective homomorphisms: injectivity gives
left cancellation, while factoring a monomorphism into its image makes
the surjective part both monic and regular epic, hence an isomorphism by
Lemma~\ref{f:descent}. Thus the internal subobjects and relations really
are the subalgebras used in the next step.

For a congruence $E$ on $A$, form the set $A/E$ and define every operation
on classes by applying it to representatives. Closure of $E$ under the
operation makes this well-defined; the equations descend because they hold
on representatives. The quotient map has kernel pair exactly $E$ and has
the same coequalizer property just proved. Thus every internal equivalence
relation is effective, proving Barr-exactness.

Finally let $R\subseteq A\times A$ be reflexive and an internal relation.
It is closed under the term $p$. If $xRy$, apply $p$ to
$(y,y),(x,y),(x,x)$ to obtain $(y,x)$, proving symmetry. If $xRy$ and $yRz$,
apply $p$ to $(x,y),(y,y),(y,z)$ to obtain $(x,z)$, proving transitivity.
Thus every reflexive internal relation is an equivalence relation.
The pointed assertion is precisely the extra initial-object hypothesis.
\end{proof}

The displayed terms are operations in a presentation. For a nonabelian group,
$G^3\to G$, $(x,y,z)\mapsto xy^{-1}z$, need not be a group homomorphism.
What we used is that a subalgebra is closed under terms, including that term.
This distinction lets algebraic expressions justify a categorical property
without confusing them with arrows in the category.

\subsection{Reading the axioms in familiar examples}\label{e:examples}

For groups use multiplication, inverse and identity, with their axioms.
The one-element group is initial and terminal, and
$p(x,y,z)=xy^{-1}z$ has the required identities. For left modules over a fixed
unital ring $R$, use addition, negative, zero, and an operation $x\mapsto rx$
for every $r\in R$, with the module equations. The zero module is initial
and terminal, and $p(x,y,z)=x-y+z$ works. Neither argument uses a dimension,
a composition series, or Jordan--H\"older.

We check that the categorical quotient data are the familiar ones. In a
group $G$, the kernel $N$ of a homomorphism is normal, since
$q(gng^{-1})=q(g)q(n)q(g)^{-1}=1$. Conversely, for a normal subgroup $N$,
the relation $x\sim y$ if $x^{-1}y\in N$ is an equivalence relation
compatible with multiplication and inverses. For example, if $x'=xn$ and
$y'=ym$, then $x'y'=xy(y^{-1}ny)m$ differs from $xy$ by an element of $N$.
The congruence quotient just constructed is the usual group $G/N$.

More generally, if $E$ is a group congruence, its identity class $N$ is a
normal subgroup: closure gives products and inverses, and conjugating a pair
$(n,1)\in E$ by $(g,g)\in E$ gives $(gng^{-1},1)\in E$. Translating
$(x,y)$ by $(x^{-1},x^{-1})$ shows $xEy$ exactly when $x^{-1}y\in N$.
Thus normal subgroups and these quotient relations agree.

For modules, a congruence has zero class a submodule $N$, and translating
pairs gives $x\sim y$ exactly when $y-x\in N$. The quotient operations are
$(x+N)+(y+N)=x+y+N$ and $r(x+N)=rx+N$. Consequently all isomorphisms supplied
by the categorical theorem are $R$-linear. A quotient-simple group is exactly
a nontrivial group with no proper nontrivial normal subgroup; a quotient-simple
module is exactly a nonzero module with no proper nonzero submodule.

We therefore obtain both classical theorems from Theorem~\ref{j:jordan-holder}:
two composition series of a group have the same group factors, and two of an
$R$-module have the same module factors. This is one specialization step for
each category, not a second proof of the comparison theorem.

The same verification covers groups with specified endomorphism operators:
include those operators and their preservation equations in the theory.
Subobjects and quotients retain them. It also covers Lie algebras over a field,
using the additive term $x-y+z$, and nonunital associative algebras over a
commutative ring. Their kernels are respectively ideals and two-sided ideals.
For Lie algebras our categorical use of simple includes a one-dimensional
abelian algebra. Unital rings with identity-preserving maps do not have a zero
object and are not a direct instance of this pointed statement.

\subsection{Reading one cell as the familiar group butterfly}
\label{e:butterfly-dictionary}

Take $A_1\nml A_2\leq G$ and $B_1\nml B_2\leq G$.
Let $r,s$ be the partial equivalence relations on $G$ supplied by
$A_2\twoheadrightarrow A_2/A_1$ and $B_2\twoheadrightarrow B_2/B_1$.
Put
\[
 C=A_2\cap B_2,\qquad U=A_1\cap B_2,\qquad W=A_2\cap B_1.
\]
Applying $s$ and then $r$ to a subset means taking the appropriate
congruence classes inside $B_2$ and then inside $A_2$. In particular,
$sG=B_2$ and $s0=B_1$, so the endpoint identities give
\[
 (rsr)G=rsG=A_1C,\qquad (rsr)0=rs0=A_1W.
\]
Here $0$ denotes the identity subgroup. Interchanging the two relations gives
support $B_1C$ and identity class $B_1U$ for $srs$. Since a group
congruence is determined by its identity class, the two cell quotients are
\[
 Q(rsr)=\frac{A_1C}{A_1W},\qquad
 Q(srs)=\frac{B_1C}{B_1U}.
\]

We can see their common quotient and every normality assertion directly.
The groups $U,W$ are normal in $C$: conjugation by an element of $C$
preserves $A_1$ or $B_1$, respectively, and remains in the other ambient
subgroup. Consequently $UW=WU$ is a normal subgroup of $C$.
The product $A_1C$ is a subgroup because $A_1\nml A_2$ and $C\leq A_2$.
In $A_1C/A_1$, the image of $C$ is the whole group and the image of
$W\nml C$ is normal. Its inverse image is $A_1W$, proving
$A_1W\nml A_1C$. The same argument proves $B_1U\nml B_1C$.

Now use the two homomorphisms
\[
 C\longrightarrow A_1C/A_1W,\quad c\longmapsto c(A_1W),
 \qquad
 C\longrightarrow B_1C/B_1U,\quad c\longmapsto c(B_1U).
\]
Both are onto, since the extra factor $A_1$ or $B_1$ is killed.
The first kernel is $C\cap A_1W=UW$: if $c=aw$ lies in $C$, then
$a=cw^{-1}\in A_1\cap C=U$; the converse inclusion is immediate.
The second kernel is $C\cap B_1U=WU=UW$, by the same calculation.
The effective quotient property therefore gives
\[
\begin{tikzcd}[column sep=large,row sep=large]
 &C\arrow[dl,two heads]\arrow[d,two heads]\arrow[dr,two heads]&\\
 A_1C/A_1W & C/UW\arrow[l,"\sim"']\arrow[r,"\sim"] & B_1C/B_1U .
\end{tikzcd}
\]
These isomorphisms give the usual butterfly, with its maps specified.
They are also the previously constructed cell comparison: for every
element of the correspondence $rs$, its two ends can be represented as
$bc$ and $ac$, with $b\in B_1$, $a\in A_1$ and $c\in C$.
Both quotient maps then use the same class of $c$. The characterization
in Proposition~\ref{r:cell-isomorphism} identifies the maps.
Thus the subgroup products are the concrete supports and identity
classes of one categorical construction.

\subsection{A comparison that changes the order}\label{e:cyclic-six}

In $G=C_6=\mathbb Z/6\mathbb Z$ put $A_1=\{0,3\}$ and
$B_1=\{0,2,4\}$. Reduction modulo $3$ has kernel $A_1$ and reduction
modulo $2$ has kernel $B_1$. Hence
\[
0<A_1<G,\qquad 0<B_1<G
\]
have factors $(C_2,C_3)$ and $(C_3,C_2)$. Every nonzero element of $C_2$
or $C_3$ generates the group, so the factors are simple.
The relation comparison becomes the following concrete rectangle:
\[
\begin{array}{c|cc}
 & B_1/0&G/B_1\\\hline
A_1/0&0&C_2\\
G/A_1&C_3&0
\end{array}
\]
Indeed $A_1\cap B_1=0$ and $A_1+B_1=G$, since $3+4=1$ modulo $6$.
The two nonzero comparisons are induced by $a\mapsto a+B_1$ and
$b\mapsto b+A_1$. Each is onto by the sum equality and injective by the
intersection equality. These are actual isomorphisms, not just equalities
of cardinalities. A refined chain need not contain both $A_1$ and $B_1$;
the matching occurs between its factors.

Normality is always in the next term. In $A_4$, for example,
$1<V_4<A_4$ has both terms normal in the ambient group, but $V_4$ is not
simple. Insert one of its order-two subgroups to obtain a composition series.
That subgroup is normal in the abelian group $V_4$ and is not normal in
$A_4$: conjugation by a three-cycle permutes the three double transpositions.
The lattice of subgroups normal in the ambient group describes a different
filtration problem.

\subsection{Finite basis size, with no dimension in the prerequisites}
\label{e:dimension}

Let $k$ be a field. A list $v_1,\ldots,v_n$ is a basis if it spans the vector
space $V$ and every relation $\sum_i a_i v_i=0$ has all $a_i=0$.
The module construction above provides subspaces, quotient spaces and their
linear maps before any notion of dimension is introduced.

\begin{theorem}[Invariance of finite basis size]\label{e:dimension-theorem}
If $V$ has finite bases $v_1,\ldots,v_n$ and $w_1,\ldots,w_m$, then $n=m$.
\end{theorem}
\begin{proof}
Put $V_i=\Span(v_1,\ldots,v_i)$, with $V_0=0$. Each inclusion is strict:
otherwise $v_i$ would be a linear combination of the preceding vectors,
contradicting independence. The map
\[
k\longrightarrow V_i/V_{i-1},\qquad a\longmapsto av_i+V_{i-1}
\]
is onto by the definition of span. If its image is zero, a relation with
coefficient $a$ on $v_i$ shows $a=0$. It is therefore an isomorphism of
$k$-modules. The module $k$ is simple: a submodule containing $a\ne0$
contains $a^{-1}a=1$ and every scalar multiple of $1$.
The basis thus supplies a composition series with $n$ factors isomorphic
to $k$. The other basis supplies one with $m$ factors. Apply the general
Jordan--H\"older theorem in $k$-modules. If $V=0$, independence forces
both lists empty, giving the same conclusion.
\end{proof}

We may now define $\dim_k V$ to be this invariant basis size. We have not
used bases of arbitrary subspaces, rank--nullity, determinants, exterior
powers, or a pre-existing dimension function. Those are possible later
constructions; they are not concealed inputs to this proof.

\subsection{What a factor list leaves unresolved}\label{e:extensions}

Both $C_4$ and $C_2\times C_2$ have factors $C_2,C_2$. Their exact sequences
can be written
\[
0\longrightarrow C_2\xrightarrow{\,a\mapsto2a\,}C_4
\longrightarrow C_2\longrightarrow0,
\qquad
0\longrightarrow C_2\longrightarrow C_2\times C_2
\longrightarrow C_2\longrightarrow0,
\]
where the last two maps in the second sequence are first-factor inclusion
and second-factor projection. The first quotient is reduction modulo $2$.
The displayed kernels and surjections verify exactness directly.
Yet $C_4$ has an element of order four, whereas every element of
$C_2\times C_2$ is killed by multiplication by two. They are not isomorphic.
The general theorem compares the simple pieces. It does not split the
extension joining them. The following parts develop measurements and
constructions that retain more of this attachment information.

% END INPUT 04-instances-and-dimension.tex


% BEGIN INPUT 05-modules-and-finite-tools.tex
\section{Keeping the operations when we change the object}
\label{m:section}

The common comparison has produced intrinsic simple factors. Its group example
also showed a limit: the factors of $C_4$ and $C_2\times C_2$ agree, although
the groups do not. We now develop that question in a setting where the
attachments can be calculated explicitly. The objects will be modules,
and our main examples will be vector spaces equipped with an operator.
We retain the operator in every map, just as the dimension argument
retained scalar multiplication in every quotient.

\subsection{Modules and their maps}

A ring $R$ has an abelian group under addition, an associative multiplication
with identity $1$, and the two distributive laws. Multiplication need not be
commutative. A left $R$-module is an abelian group $M$ with operations $rx$
for $r\in R$, $x\in M$, satisfying
\[
 r(x+y)=rx+ry,\quad (r+s)x=rx+sx,\quad
 r(sx)=(rs)x,\quad 1x=x.
\]
The same cancellation arguments as for vector spaces give $r0=0$,
$0x=0$ and $(-1)x=-x$. A submodule contains zero and is closed under
addition and all these scalar operations. A module map preserves addition
and satisfies $f(rx)=rf(x)$. In this part all unspecified maps between
modules are module maps. We work over a fixed ring whenever forming
extensions or comparing modules.

Intersections of submodules are submodules. Their sum $U+W$ consists of
the elements $u+w$; distributivity proves that it is a submodule, and it
is the smallest one containing both. The direct sum of two modules is
the set of pairs with coordinatewise addition and scalar action. Each
module identity holds in both coordinates. Its inclusions and projections
are module maps. A map out of $A\oplus B$ is determined by its restrictions
to $(a,0)$ and $(0,b)$, since $(a,b)=(a,0)+(0,b)$.

For a submodule $N\subseteq M$, the additive quotient $M/N$ is already
defined by the group construction. Set $r(x+N)=rx+N$. Replacing $x$ by
$x+n$ changes $rx$ by $rn\in N$, so this action is well-defined;
the module identities descend on representatives. The projection is a
module map. If a module map $f:M\to X$ kills $N$, its unique additive
factorization $\bar f:M/N\to X$ is a module map, because
\[
 \bar f(r(x+N))=f(rx)=rf(x)=r\bar f(x+N).
\]
Thus the quotient universal property retains every operation required here.

Kernels, images and preimages of submodules are submodules: apply a map
to a sum or a scalar multiple, and use the corresponding closure in the
target. An image of a submodule is the image of the map restricted to it.
Consequently the first isomorphism theorem supplies a module isomorphism
\[
 M/\Ker f\iso\im f.
\]
The inverse of a bijective module map is a module map: writing $y=f(x)$
gives $f^{-1}(ry)=rx$, and the analogous calculation handles sums.
All the short exact sequences below mean injectivity, surjectivity, and
image equal to kernel for these maps.

\subsection{Preimages, quotients, and the local comparison}

For a surjection $q:M\to Q$ with kernel $N$, inverse image and direct
image give an inclusion-preserving correspondence
between submodules of $Q$ and submodules of $M$ containing $N$.
Indeed, for a submodule $H$ of $M$,
\[
 q^{-1}(q(H))=N+H:
 \quad q(x)=q(h)\ \Longleftrightarrow\ x-h\in N.
\]
Every quotient map and induced isomorphism in this correspondence is
linear over $R$. For $N\subseteq L\subseteq M$ this also proves
$(M/N)/(L/N)\cong M/L$ by $(x+N)+(L/N)\mapsto x+L$.

The kernel calculation at the heart of the butterfly remains valid:
if $A_1\subseteq A_2$ and $B_1\subseteq B_2$ are submodules of $M$,
write $C=A_2\cap B_2$, $U=A_1\cap B_2$, $W=A_2\cap B_1$.
The map
\[
 C\longrightarrow (A_1+C)/(A_1+W),\qquad c\longmapsto c+(A_1+W)
\]
is onto, and its kernel is $U+W$. To check the nontrivial inclusion,
if $c=a+w$ with $c,w\in C$ and $a\in A_1$, then $a=c-w\in U$.
The reverse inclusion follows by substitution. Exchanging $A$ and $B$
gives the same kernel and the other wing. This is a concrete calculation
of the common support quotient, retaining the $R$-action. It is available independently
of any theorem about composition series.

\subsection{Finite-dimensional tools, after dimension has been proved}
\label{m:finite-tools}

The later operator examples use subspace bases and lifts. We establish
their existence here, after Theorem~\ref{e:dimension-theorem}; none is
being fed back into the earlier proof of dimension.

\begin{lemma}\label{m:bases}
If a vector space $V$ has a finite basis of size $n$, every subspace has
a finite basis. Every independent list in $V$ extends to a basis of $V$.
For a subspace $W$, a basis of $W$ followed by lifts of a basis of $V/W$
is a basis of $V$.
\end{lemma}
\begin{proof}
Start with any finite independent list. Pass through a fixed finite
spanning list of $V$, appending a vector exactly when it is outside
the span of the vectors already present. Each append preserves
independence: a relation with nonzero coefficient on the new vector
would put it in the old span. At the end the list spans $V$, so it is
a basis. The already proved invariance of basis size shows that the
original independent list had at most $n$ members.
An arbitrary independent set cannot contain $n+1$ distinct vectors,
so it too is finite and the same conclusion applies.

To construct a basis of $W$, start with the empty list and, as long
as it does not span $W$, choose an element outside its span. This
cannot produce $n+1$ independent vectors. Thus it stops after finitely
many steps with a basis of $W$. This requires only finitely many
choices, bounded by $n$.

The images in $V/W$ of a finite basis of $V$ span the quotient.
Applying the append procedure starting with the empty list to those
images gives a quotient basis. Choose one lift of each of its finitely
many members. The lifts together with a basis of $W$ span $V$:
subtract the lifted expression of a vector's quotient class, leaving
an element of $W$. In a relation among this list, projection to $V/W$
first forces all lift coefficients to vanish; independence in $W$
then forces the remaining coefficients to vanish.
\end{proof}

In particular, $\dim V=\dim W+\dim(V/W)$. Applying the first
isomorphism theorem to a linear map gives
$\dim V=\dim\Ker f+\dim\im f$. These are consequences of the
dimension theorem already proved, and will be used only from here on.
The same basis-lifting proof gives a vector-space section of any
surjective linear map onto a finite-dimensional space. Such a section
need not preserve a further module or operator structure.

\subsection{An operator is a polynomial action}
\label{m:operators}

For a field $k$, $k[t]$ consists of finite formal sums
$\sum_{i\geq0}a_it^i$. Addition is coefficientwise and multiplication
is convolution: the coefficient of $t^r$ in a product is
$\sum_{i+j=r}a_ib_j$. All sums involved are finite. Associativity
holds because either parenthesization of a triple product has coefficient
$\sum_{i+j+\ell=r}a_ib_jc_\ell$; distributivity is the corresponding
distributivity of finite sums. The constant polynomial $1$ is the identity.

Given a $k$-linear endomorphism $T$ of $V$, define
\[
 p(t)v=p(T)v=\sum_i a_iT^iv,\qquad T^0=\id.
\]
Induction gives $T^iT^j=T^{i+j}$, so convolution proves
$p(T)q(T)=(pq)(T)$. The other module axioms follow from linearity.
Conversely, a $k[t]$-module has its underlying $k$-vector space and
the linear operator $T(v)=tv$; the module axioms recover the displayed
polynomial action. These constructions are inverse.

A linear map $f:(V,T)\to(V',T')$ is a module map exactly when
$fT=T'f$. Necessity uses the polynomial $t$; sufficiency follows by
induction on powers and addition. Submodules are exactly the
$T$-stable subspaces. Thus the square
\[
\begin{tikzcd}[column sep=large,row sep=large]
V\arrow[r,"T"]\arrow[d,"f"']&V\arrow[d,"f"]\\
V'\arrow[r,"T'"']&V'
\end{tikzcd}
\]
records the extra condition that distinguishes an operator-preserving
map from an arbitrary linear map.

For $\lambda\in k$, let $S_\lambda$ be $k$ with $t$ acting by
$\lambda$. It is simple as a module because its only $k$-subspaces
are zero and itself. It is also $k[t]/(t-\lambda)$: evaluation
$p\mapsto p(\lambda)$ is onto, and the identity
$t^i-\lambda^i=(t-\lambda)\sum_{j=0}^{i-1}t^{i-1-j}\lambda^j$
shows that its kernel is $(t-\lambda)$.

We also use $P_r=k[t]/(t^r)$ for $r\geq1$. Its vectors have unique
representatives of degree less than $r$, by separating the terms of
degree at least $r$. Thus $1,t,\ldots,t^{r-1}$ is a vector-space
basis, and multiplication by $t$ shifts along a nilpotent chain.
No polynomial factorization theorem is required for these modules.

\subsection{Simple layers and finite length}
\label{m:length-foundations}

A module is simple if it is nonzero and has only zero and itself as
submodules. A composition series is a finite strictly increasing
submodule filtration from zero to the whole module with simple
successive quotients. A module has finite length if it admits such
a series; at this stage that phrase is an existence condition.
The zero module has the series with one term and no factors.

\begin{lemma}\label{m:simple-map}
A nonzero module map between simple modules is an isomorphism.
\end{lemma}
\begin{proof}
Its image is nonzero, hence is the whole target. Its kernel cannot be
the whole source, hence is zero. The kernel test gives injectivity,
and the inverse of a bijective module map preserves the module operations.
\end{proof}

Here a $k$-algebra is a ring $R$ with a unital map from $k$ into its
center. Restricting scalars makes each $R$-module a $k$-vector space,
and every $R$-submodule is a $k$-subspace.
Every finite-dimensional module over a $k$-algebra has a composition
series. If it is nonzero, choose a proper submodule of largest possible
vector-space dimension; zero is one candidate. The possible dimensions
form a finite set by Lemma~\ref{m:bases}. Its quotient is simple:
an intermediate nonzero proper submodule would lift to a larger proper
submodule. Repeat inside the chosen proper submodule. Dimensions
strictly decrease, so the process ends. This proves existence without
using uniqueness of composition factors.

We will also need closure of finite-length modules under submodules
and quotients over a general ring. Let
$0=M_0\subsetneq\cdots\subsetneq M_n=M$ be a composition series
and $N\subseteq M$ a submodule. Set
$N_i=N\cap M_i$ and $Q_i=(M_i+N)/N$.
For each $i$, inclusion and projection give a short exact sequence
\begin{equation}\label{m:layer-exact}
0\longrightarrow N_i/N_{i-1}\longrightarrow M_i/M_{i-1}
\longrightarrow Q_i/Q_{i-1}\longrightarrow0.
\end{equation}
The first map sends $n+N_{i-1}$ to $n+M_{i-1}$ and is injective.
The second sends $m+M_{i-1}$ to $(m+N)+Q_{i-1}$ and is onto.
Its kernel consists of $m=a+n$ with $a\in M_{i-1}$ and $n\in N$;
then $n=m-a\in N_i$, proving equality with the first image.
The simple middle module forces precisely one outer module to be zero
and the other to be isomorphic to it. Deleting repetitions from the
$N_i$ and $Q_i$ therefore gives composition series for $N$ and $M/N$.

Conversely, if $A$ and $B$ have composition series and
$0\to A\to E\xrightarrow{q}B\to0$ is short exact, identify $A$
with its image in $E$. Append to its series the preimages under $q$
of a composition series of $B$. The quotient correspondence proves
strictness, and the maps
\[
 q^{-1}(B_j)/q^{-1}(B_{j-1})\longrightarrow B_j/B_{j-1},\qquad
 e+q^{-1}(B_{j-1})\longmapsto q(e)+B_{j-1}
\]
are isomorphisms by their kernels and surjectivity. This constructs
a composition series for $E$. Its factors are exactly the two given
lists together. No uniqueness theorem is used in this construction.

We now have the objects and operations needed for an independent
spectral comparison of two module filtrations. Length and multiplicities
are already intrinsic by Theorem~\ref{j:jordan-holder}. The following
page calculations will compare the factors directly, without using that
theorem as a premise.

% END INPUT 05-modules-and-finite-tools.tex


% BEGIN INPUT 06a-common-comparison-pages.tex
\section{Comparison pages with the operations retained}
\label{p:chapter}

The relation rectangle gives one proof for every pointed Barr-exact
Mal'cev category. There is another useful way to watch the comparison
proceed: let simple factors persist through a sequence of pages until
maps pair them. We can carry this calculation out simultaneously for
groups and modules, without pretending that arbitrary group maps behave
like linear maps.

For this chapter a group with operators means a group equipped with a
fixed family of endomorphisms. Subgroups and normal subgroups under
discussion must be invariant under all those endomorphisms, and maps
commute with them. Modules are included by their additive groups and the
operators $x\mapsto rx$ for every scalar $r$. Quotient simplicity is
simplicity with this structure retained. Ordinary groups are the case
of an empty family of extra operators. This is a concrete common instance
of the preceding framework; the page construction below is asserted for
this instance, with composition series as input.

\subsection{The normality that makes the pages possible}

Write two composition series as
\[
1=A_0\nml A_1\nml\cdots\nml A_m=G,
\qquad
1=B_0\nml B_1\nml\cdots\nml B_n=G.
\]
We use multiplicative notation, including for the identity subgroup $1$.
Extend both filtrations constantly below zero and beyond their last
indices. Put $S_i=A_i/A_{i-1}$ and $T_j=B_j/B_{j-1}$ for the real
indices $1\leq i\leq m$, $1\leq j\leq n$.

\begin{lemma}\label{p:normality}
A subnormal subgroup of a simple group with operators is either trivial
or the whole group. For every integer $k$ and every real index $j$,
\[
 B_{j-1}(A_k\cap B_j)
 \quad\text{is either }B_{j-1}\text{ or }B_j.
\]
In particular this product is an invariant normal subgroup of $B_j$.
Similarly the image of $A_i\cap B_k$ in $S_i$ is either trivial or all
of $S_i$.
\end{lemma}
\begin{proof}
Take a finite normal chain from a proper subnormal subgroup $H$ to a
simple group $S$. The term immediately before the first occurrence of
$S$ is proper and normal in $S$, so it is trivial. The earlier terms,
including $H$, must be trivial. This proves the first claim with
invariance under the operators included.

Intersecting a normal chain with a fixed subgroup preserves adjacent
normality: conjugating an element of $U\cap K$ by an element of
$V\cap K$ stays in both $U$ and $K$ when $U\nml V$.
Taking images under a homomorphism also preserves adjacent normality,
since $f(v)f(u)f(v)^{-1}=f(vuv^{-1})$. Intersections and images here
remain invariant under every operator. Thus $A_k\cap B_j$ is subnormal
in $B_j$, by intersecting the tail of the $A$-chain with $B_j$; the
constant extensions at the endpoints cause no change to this assertion.
Its image in the simple factor $T_j$ is trivial or everything. The
preimage of that image is exactly $B_{j-1}(A_k\cap B_j)$, proving
the stated alternatives. The same argument with the roles exchanged
proves the last assertion.
\end{proof}

This check is essential. A nonzero homomorphism between simple groups
need not be onto: an inclusion $C_2\hookrightarrow A_5$ is an immediate
counterexample. The maps constructed below have subnormal, and in fact
normal, images for a specified reason. That is what makes the cancellation
argument valid.

\subsection{Terms, maps, and the next page}

For $r\geq1$ define two rows
\begin{align}
 U_r(i)&=
 \frac{A_i\cap B_{n+i-r}}{A_{i-1}\cap B_{n+i-r}},
       &&1\leq i\leq m,\label{p:source}\\
 V_r(j)&=
 \frac{B_j}{B_{j-1}(A_{j+r-n-1}\cap B_j)},
       &&1\leq j\leq n.\label{p:target}
\end{align}
The same formulas give trivial groups outside the real index ranges.
Every source denominator is invariant and normal in its numerator by
intersection. Every target denominator is invariant and normal by
Lemma~\ref{p:normality}. Consequently these are actual quotient groups
with all their operators, not just sets of cosets.

For $j=n+i-r$, inclusion followed by quotient defines
\begin{equation}\label{p:differential}
 d_{r,i}:U_r(i)\longrightarrow V_r(j),\qquad [x]\longmapsto[x].
\end{equation}
Its source denominator is $A_{i-1}\cap B_j$, and this is contained in
the target denominator $B_{j-1}(A_{i-1}\cap B_j)$. Thus the map is
well-defined, is a homomorphism, and preserves the operators. Each source
has one designated target; each target has one designated incoming source.

\begin{proposition}[The page transition]\label{p:transition}
With $j=n+i-r$, the inclusion and quotient maps give canonical
isomorphisms
\[
 U_{r+1}(i)\cong\ker d_{r,i},\qquad
 V_{r+1}(j)\cong V_r(j)/\im d_{r,i}.
\]
The image in the second formula is an invariant normal subgroup.
\end{proposition}
\begin{proof}
Put $H=A_i\cap B_j$, $D=A_{i-1}\cap B_j$, and $K=B_{j-1}$.
Thus the map in question is $H/D\to B_j/KD$. The subgroup $D$ is
normal in $H$, $K$ is normal in $B_j$, and both $KD$ and $KH$ are
normal in $B_j$ by Lemma~\ref{p:normality}. If $j$ is outside the real
target range, the constant extension gives $B_{j-1}=B_j$, so these
normality assertions hold automatically. This also covers transitions
whose designated source or target lies outside the initial rows.

We first compute its kernel. Since $D\leq H$,
\[
                     H\cap KD=(H\cap K)D.
\]
Indeed, if $h=kd\in H$ with $k\in K$ and $d\in D$, then
$k=hd^{-1}\in H\cap K$. The reverse inclusion follows by multiplying
elements already in $H$ and in $KD$. Restricting the quotient map to
$H\cap K$ therefore gives a surjection onto the kernel with kernel
$D\cap K$. Hence
\[
 \ker d_{r,i}\cong\frac{H\cap K}{D\cap K}
 =\frac{A_i\cap B_{j-1}}{A_{i-1}\cap B_{j-1}}
 =U_{r+1}(i).
\]
This isomorphism is induced by inclusion; all its operations are retained.

The image is $KH/KD$, a normal subgroup of $B_j/KD$. The next target
projection has precisely this kernel, so the compatible iterated quotient
map gives
\[
 \frac{V_r(j)}{\im d_{r,i}}
 \cong \frac{B_j}{KH}
 =\frac{B_j}{B_{j-1}(A_i\cap B_j)}
 =V_{r+1}(j).
\]
The last equality uses $i=j+r-n$. These are explicit maps between
quotients, so their inverses preserve the operators as well.
\end{proof}

The transition can be displayed as two short exact sequences:
\[
\begin{tikzcd}[column sep=small]
1\arrow[r]&U_{r+1}(i)\arrow[r,hook]&U_r(i)
 \arrow[r,two heads]&\im d_{r,i}\arrow[r]&1,\\
1\arrow[r]&\im d_{r,i}\arrow[r,hook]&V_r(j)
 \arrow[r,two heads]&V_{r+1}(j)\arrow[r]&1.
\end{tikzcd}
\]
The first row retains the source part not yet seen in the target. The
second removes the target part accounted for by that map. This is the
two-row page mechanism $E_{r+1}\cong H(E_r,d_r)$, with the group
quotients justified by the normality proof. We need no addition of maps
to perform this particular calculation.

\begin{theorem}[Every layer eventually finds its partner]\label{p:cancellation}
The pages above give a bijection between the simple factors $S_i$ and
$T_j$, together with isomorphisms preserving all the operators. The proof
uses neither Schreier refinement nor Jordan--H\"older as an input.
\end{theorem}
\begin{proof}
The source $U_r(i)$ is isomorphic to the image of $A_i\cap B_{n+i-r}$
in $S_i$, by the kernel of the restricted quotient map. By
Lemma~\ref{p:normality} it is trivial or isomorphic to $S_i$.
The target denominator in~\eqref{p:target} is either $B_{j-1}$ or
$B_j$, so $V_r(j)$ is trivial or isomorphic to $T_j$.

The image of $d_{r,i}$ is normal by Proposition~\ref{p:transition}.
A nontrivial map therefore surjects onto its simple target. Its kernel
is an invariant normal subgroup of its simple source; it cannot be the
whole source, so it is trivial. Consequently every nontrivial
differential is an isomorphism. The transition then removes precisely
that source and target. When the differential is trivial, the transition
retains both terms by its canonical isomorphisms.

On the first page, $n+i-1\geq n$ for every real source index, so
$U_1(i)=S_i$. For the targets $j-n\leq0$, giving $V_1(j)=T_j$.
If $G\neq1$, then $m,n\geq1$. On every page $r\geq m+n$ we have
$n+i-r\leq0$ and $j+r-n-1\geq m$ at real indices. Every term is
therefore trivial. Each initial factor must have disappeared exactly
once, by an isomorphism differential, and each such differential removes
one source and one target. Following the transition isomorphisms to that
page gives the required pairwise factor isomorphisms. For $G=1$, both
lists are empty and the conclusion is immediate.
\end{proof}

Notice the dependency: simplicity is part of the supplied filtrations;
invariance of their lengths is the conclusion. No page count assumes
that the two rows start with the same number of terms.

\subsection{A nonabelian crossing and the module specialization}

Let $G=A_5\times C_2$ and use the two series with intermediate subgroups
$A_5\times1$ and $1\times C_2$. Their first-page factors have opposite
orders. The first differential at $(i,j)=(1,2)$ identifies
$A_5\times1$ with $G/(1\times C_2)$. It removes the two $A_5$ factors.
The second page has no nontrivial differential, but the process has not
terminated. On the third page, at $(i,j)=(2,1)$, the actual differential
is the identity on $1\times C_2$, since both $U_3(2)$ and $V_3(1)$
are that group. Its source is identified with the original factor
$G/(A_5\times1)$ by the restricted quotient map
$1\times C_2\to G/(A_5\times1)$. Thus it pairs the two remaining
$C_2$ factors. The fourth page is trivial. A quiet page need not be the
last page.

For completeness, the simplicity of $A_5$ used only in this illustration
can be checked by its conjugacy classes. Its nonidentity elements are
the $15$ double transpositions, $20$ three-cycles, and $24$ five-cycles.
Their centralizers in $A_5$ have orders $4,3,5$, respectively: commuting
permutations permute cycles of equal length; on a three- or five-cycle they
are powers of that cycle, with the odd transposition on the two unused
letters excluded in the three-cycle case; the even permutations commuting
with a double transposition form its four-element Klein group.
Counting conjugates by cosets of the centralizer therefore gives class
sizes $15,20,12,12$. A normal subgroup is a union of these classes and
the identity, and its order divides $60$ because its cosets partition
$A_5$. None of the proper nontrivial sums of $1$ and a selection from
$15,20,12,12$ divides $60$. Hence no proper nontrivial normal subgroup
exists. This finite check is not a prerequisite of the page theorem.

For modules the products in~\eqref{p:target} are sums, and the first
row of~\eqref{p:source} is a subquotient of modules. Put
$F_pX=A_{p-n}$, $F_pY=B_p$, and take $f=\id_G:X\to Y$ on the
underlying module. The page formulas in the next chapter recover exactly
$U_r(i)$ at $p=n+i$ and $V_r(j)$ at $q=j$, with the differential
landing at $q=p-r$. The next chapter proves that calculus for an
\emph{arbitrary filtered module map}, where the pages may have nonsimple
terms and nonzero limiting kernel or cokernel. The present calculation
explains why its composition-series cancellation can also be carried
out for groups with operators.

The scope is precise. We proved the relation refinement theorem for all
pointed Barr-exact Mal'cev categories. We proved this explicit page
calculation for composition series of groups with operators. We have not
asserted that every filtered group complex supports the ordinary abelian
spectral-sequence machinery. General nonabelian exact-couple theories
require additional constructions; one primary reference is
\href{https://arxiv.org/pdf/1007.0632}{Grandis, \emph{Homotopy spectral
sequences}}, Sections~2--4.

% END INPUT 06a-common-comparison-pages.tex


% BEGIN INPUT 06b-module-spectral-calculus.tex
\section{Filtered maps and their successive obstructions}\label{s:chapter}

The common Jordan--H\"older theorem is already complete. A reader following
only that theorem may pass to Section~\ref{k:section}; the purpose here
is to extend the page construction to arbitrary filtered module maps.
We now develop the module page calculus beyond composition-series input.
Its simplest acyclic example will recover the preceding comparison.
We can put the two copies of an object into the exact complex
\[
0\longrightarrow M\xrightarrow{\id}M\longrightarrow0,
\]
give its two terms different filtrations, and calculate what happens to
their layers. The complex has no homology: the identity has zero kernel
and zero cokernel. Its initial layers need not cancel immediately,
because they occur at different filtration levels. A spectral sequence
records the successive cancellations.

When the layers are simple, every nonzero cancellation map is an
isomorphism. Since nothing survives, these maps pair every factor on one
side with exactly one factor on the other. This will prove the module
Jordan--H\"older theorem without assuming composition length, composition
multiplicities, or their invariance.

We shall construct all the pages needed for this argument. The setting is
left modules over a fixed ring with identity. Only submodules, quotient
modules, and module homomorphisms are used. The argument does not invoke
the earlier categorical Jordan--H\"older theorem. The preceding chapter
has already supplied the group and operator version of the particular
simple-factor calculation. Here additive module structure permits
arbitrary filtered maps, nonsimple layers, and nonzero limiting homology.

\subsection{A map seen one layer at a time}

Let $f:X\to Y$ be a module homomorphism. Give both modules finite
increasing filtrations indexed by the integers. Thus
\[
F_pX\subseteq F_{p+1}X,\qquad F_pY\subseteq F_{p+1}Y,
\]
the filtrations are zero for sufficiently negative $p$, and they equal
the whole modules for sufficiently positive $p$. Assume that $f$ preserves
them:
\[
f(F_pX)\subseteq F_pY\quad\text{for every }p.
\]
One may regard $X\xrightarrow{f}Y$ as a chain complex in degrees $1$ and
$0$. Its homology is
\[
H_1=\Ker f,\qquad H_0=Y/f(X).
\]
There are no other terms or differentials, so the condition that two
successive differentials compose to zero is automatic.

At the first level of resolution, keep only the associated layers:
\[
X_0(p)=F_pX/F_{p-1}X,\qquad
Y_0(p)=F_pY/F_{p-1}Y.
\]
Filtration preservation gives a well-defined module map
\[
d_{0,p}:X_0(p)\longrightarrow Y_0(p),\qquad [x]\longmapsto[f(x)].
\]
If $x$ is changed by an element of $F_{p-1}X$, its image changes by an
element of $F_{p-1}Y$. This is the entire well-definedness check.

Taking the kernel and cokernel of these maps retains the source classes
whose images have fallen to a lower layer, and the target classes not yet
accounted for by images. We then ask whether the retained source classes
reach a target one layer lower, two layers lower, and so on. The page
number measures how far down the filtration the map is being examined.

\subsection{The pages and their maps}

For $r\geq1$ and integers $p,q$, define
\begin{align}
X_r(p)
&=
\frac{F_pX\cap f^{-1}(F_{p-r}Y)}
     {F_{p-1}X\cap f^{-1}(F_{p-r}Y)},
\label{s:source-page}\\[3pt]
Y_r(q)
&=
\frac{F_qY}
     {F_{q-1}Y+\bigl(F_qY\cap f(F_{q+r-1}X)\bigr)}.
\label{s:target-page}
\end{align}
All the displayed sets are submodules. In the first formula the
denominator lies in the numerator because $F_{p-1}X\subseteq F_pX$;
in the second both summands of the denominator lie in $F_qY$.

The numerator of $X_r(p)$ consists of representatives in source layer $p$
whose image has reached target level $p-r$ or lower. Its denominator
identifies representatives that were already in a lower source layer
and satisfy the same condition. On the target side, $Y_r(q)$ discards
lower target layers and images already detected from source levels up
to $q+r-1$.

More precisely, $[x]\mapsto x+F_{p-1}X$ embeds $X_r(p)$ into
$F_pX/F_{p-1}X$: the map on its numerator has exactly the displayed
denominator as kernel.
The target $Y_r(q)$ is a quotient of $F_qY/F_{q-1}Y$. Thus a surviving
source entry is retained as a submodule of its original layer, while a
surviving target entry is retained as a quotient of its original layer.

There is a module map
\begin{equation}\label{s:page-differential}
d_{r,p}:X_r(p)\longrightarrow Y_r(p-r),\qquad[x]\longmapsto[f(x)].
\end{equation}
To check it, put $q=p-r$. A numerator representative has image in $F_qY$.
Changing it by
$z\in F_{p-1}X\cap f^{-1}(F_qY)$ changes the image by
\[
f(z)\in F_qY\cap f(F_{p-1}X)
       =F_qY\cap f(F_{q+r-1}X),
\]
which is in the target denominator. The formula is therefore independent
of representatives and is linear.

For each page, the arrows go from the source row to the target row:
\[
\begin{tikzcd}[column sep=large,row sep=large]
X_r(p) \arrow[d,"d_{r,p}"] & \text{source degree }1\\
Y_r(p-r) & \text{target degree }0.
\end{tikzcd}
\]
There are no arrows out of the target row or into the source row. Hence
the square of the differential is zero.

\begin{proposition}[The next page is homology]\label{s:transition}
There are canonical isomorphisms
\begin{align*}
X_{r+1}(p)&\cong\Ker\bigl(d_{r,p}:X_r(p)\to Y_r(p-r)\bigr),\\
Y_{r+1}(q)&\cong
Y_r(q)/\im\bigl(d_{r,q+r}:X_r(q+r)\to Y_r(q)\bigr)
\end{align*}
for every $r\geq0$.
\end{proposition}

\begin{proof}
For $r=0$, the kernel of $d_{0,p}$ consists of the classes represented by
$x\in F_pX$ with $f(x)\in F_{p-1}Y$. Since
$f(F_{p-1}X)\subseteq F_{p-1}Y$, its denominator is $F_{p-1}X$.
This is exactly $X_1(p)$. Quotienting $Y_0(q)$ by the image of $d_{0,q}$
gives
\[
F_qY/\bigl(F_{q-1}Y+f(F_qX)\bigr)=Y_1(q).
\]

Now take $r\geq1$ and put $q=p-r$. The numerator for $X_{r+1}(p)$ is
contained in the numerator for $X_r(p)$. Its intersection with the old
denominator is
\[
\bigl(F_pX\cap f^{-1}(F_{q-1}Y)\bigr)
\cap\bigl(F_{p-1}X\cap f^{-1}(F_qY)\bigr)
=F_{p-1}X\cap f^{-1}(F_{q-1}Y).
\]
That is precisely the new denominator. Thus inclusion of representatives
induces an injection $X_{r+1}(p)\hookrightarrow X_r(p)$.
Every class in its image maps to zero under $d_{r,p}$, because its image
under $f$ lies in $F_{q-1}Y$.

Conversely, suppose $[x]\in X_r(p)$ maps to zero. Then
\[
f(x)\in F_{q-1}Y+\bigl(F_qY\cap f(F_{p-1}X)\bigr).
\]
Choose $z\in F_{p-1}X$ with $f(z)\in F_qY$ such that
$f(x)-f(z)\in F_{q-1}Y$. The element $z$ is in the denominator of
$X_r(p)$, so $x-z$ represents the same class there. But $x-z$ is in the
numerator of $X_{r+1}(p)$. This proves that the injection identifies
$X_{r+1}(p)$ with the kernel.
The choice of $z$ proves that a suitable representative exists; the
canonical map itself is induced by inclusion and does not depend on that
choice.

On the target side, the denominators increase with $r$, giving a natural
surjection $Y_r(q)\twoheadrightarrow Y_{r+1}(q)$. The additional elements
killed are the classes of
\[
F_qY\cap f(F_{q+r}X).
\]
Every such element is $f(x)$ for an $x\in F_{q+r}X$ whose image lies in
$F_qY$, so it comes from $X_r(q+r)$. Conversely, every image of
$d_{r,q+r}$ has this form. The kernel of the surjection is therefore
exactly that image, proving the asserted cokernel formula.
\end{proof}

This calculation is a spectral sequence, in the ordinary algebraic
sense. To see the conventional indexing, set
\[
E_r^{p,q}=
\begin{cases}
X_r(p),&p+q=1,\\
Y_r(p),&p+q=0,\\
0,&\text{otherwise}.
\end{cases}
\]
Then $d_r$ has bidegree $(-r,r-1)$, the convention for an increasing
filtration on a chain complex, and Proposition~\ref{s:transition} says
exactly
\[
E_{r+1}=H(E_r,d_r).
\]
The explicit formulas establish this relation; no general spectral
sequence existence theorem is being assumed.
The support lies on the two total-degree diagonals $p+q=1$ and $p+q=0$.
The index $q$ may be negative; no first-quadrant assumption is needed.
The finite filtration support will supply the convergence bound.

The construction also respects maps. Suppose
\[
\begin{tikzcd}
X\arrow[r,"f"]\arrow[d,"a"']&Y\arrow[d,"b"]\\
X'\arrow[r,"f'"']&Y'
\end{tikzcd}
\]
commutes and $a,b$ preserve filtrations. The formula $[x]\mapsto[a(x)]$
defines $X_r(p)\to X'_r(p)$: the equality $f'a=bf$ carries both numerator
and denominator into their primed versions. Similarly $[y]\mapsto[b(y)]$
defines $Y_r(q)\to Y'_r(q)$, since
$b(f(F_jX))=f'(a(F_jX))\subseteq f'(F_jX')$. These maps commute with
$d_r$ by the same equality $f'a=bf$. The source inclusions and target
quotient maps defining the next page commute with $a,b$ as well, so
the identifications with homology are natural. Thus every page and differential
belongs to the given filtered map, rather than to choices of bases.

\subsection{What eventually survives}

\begin{proposition}[Convergence in this two-term setting]\label{s:convergence}
For sufficiently large $r$, the pages stabilize to
\begin{align*}
X_r(p)&=
\frac{F_pX\cap\Ker f}{F_{p-1}X\cap\Ker f},\\
Y_r(q)&=
\frac{F_qY}{F_{q-1}Y+(F_qY\cap f(X))}.
\end{align*}
The first is the associated layer of the induced filtration on $\Ker f$;
the second is canonically the associated layer of the quotient filtration
on $Y/f(X)$. In particular, if $f$ is an isomorphism, every page entry
is eventually zero.
\end{proposition}

\begin{proof}
Fix $p$. For sufficiently large $r$, $F_{p-r}Y=0$, so the two inverse
images in \eqref{s:source-page} become $\Ker f$. Fix $q$. For sufficiently
large $r$, $F_{q+r-1}X=X$, giving the second formula from
\eqref{s:target-page}.

Let $\pi:Y\to Y/f(X)$ be the quotient map. The quotient filtration is
$F_q(Y/f(X))=\pi(F_qY)$. The map
\[
F_qY\longrightarrow\pi(F_qY)/\pi(F_{q-1}Y),\qquad
y\longmapsto\pi(y)+\pi(F_{q-1}Y)
\]
is surjective. An element $y\in F_qY$ is in its kernel exactly when
$y=y_0+f(x)$ for $y_0\in F_{q-1}Y$. Then $f(x)=y-y_0\in F_qY$,
so the kernel is $F_{q-1}Y+(F_qY\cap f(X))$, as asserted.

Only finitely many filtration columns are nonzero on page $0$.
Proposition~\ref{s:transition} shows that a zero entry stays zero.
We may therefore take a single sufficiently large $r$ for all relevant
columns. If $f$ is an isomorphism, its kernel and cokernel are zero,
so all the stabilized entries are zero.
\end{proof}

Notice what the proof used: bounded filtrations, kernels, images, and
explicit quotient maps. It did not count dimensions or composition
factors. This is the point that makes a noncircular application to
Jordan--H\"older possible.

\subsection{Two composition series inside an acyclic complex}

Suppose that a module $M$ has two composition series
\[
0=A_0\subsetneq A_1\subsetneq\cdots\subsetneq A_m=M,
\qquad
0=B_0\subsetneq B_1\subsetneq\cdots\subsetneq B_n=M.
\]
Thus each $A_i/A_{i-1}$ and each $B_j/B_{j-1}$ is a nonzero simple
module. The numbers $m$ and $n$ are not assumed equal.

Extend the notation by $A_i=0$ for $i\leq0$, $A_i=M$ for $i\geq m$,
and similarly for $B$. Take $X=Y=M$ and $f=\id_M$, with
\begin{equation}\label{s:staggered-filtrations}
F_pX=A_{p-n},\qquad F_pY=B_p.
\end{equation}
The entire $B$ filtration occurs before the nonzero stages of the
$A$ filtration. This makes the identity filtration preserving:
if $p\leq n$, then $F_pX=0$; if $p>n$, then $F_pY=M$.
We have a filtered complex even though the inverse identity need not
preserve these two filtrations.

Its nonzero page-$0$ entries are precisely
\[
X_0(n+i)=A_i/A_{i-1}\quad(1\leq i\leq m),
\qquad
Y_0(j)=B_j/B_{j-1}\quad(1\leq j\leq n).
\]
Their column ranges are disjoint, so $d_0=0$ and $E_1=E_0$.
The positions can be summarized as follows:
\[
\begin{array}{c|c|c}
\text{row}&\text{nonzero columns}&\text{entry}\\ \hline
\text{source}&p=n+i,\quad1\leq i\leq m&A_i/A_{i-1}\\
\text{target}&p=j,\quad1\leq j\leq n&B_j/B_{j-1}
\end{array}
\]
Each $d_r$ travels from a source column to the target column $r$ places
to its left.

\begin{theorem}[Module Jordan--H\"older by spectral cancellation]
\label{s:jordan-holder}
The two composition series have equal length, $m=n$, and their factors
can be paired by module isomorphisms. The nonzero differentials in the
spectral sequence of \eqref{s:staggered-filtrations} supply such a pairing.
\end{theorem}

\begin{proof}
Recall first why a nonzero map $u:S\to T$ between simple modules is an
isomorphism. Its kernel is a submodule of $S$, and cannot be all of $S$
when $u\neq0$, so it is zero. Its image is a nonzero submodule of $T$,
so it is $T$. The first isomorphism theorem now gives the assertion.

Initially every nonzero page entry is simple. Inductively suppose this
is true on page $r$. A differential is either zero, leaving its source
and target unchanged on the next page, or is an isomorphism between two
simple entries, making both zero on the next page. There is only one
possible outgoing differential from each source position and one
possible incoming differential to each target position on a fixed page.
There are no incoming differentials to source positions and no outgoing
differentials from target positions. Proposition~\ref{s:transition}
therefore proves the induction.

Keep each original factor as a label on its position until that position
vanishes. A zero differential carries the label forward by the canonical
kernel or quotient identification. A nonzero differential pairs its two
labels by an isomorphism and removes them. Removed positions stay zero,
so no label is used twice. Proposition~\ref{s:convergence} applies to
$f=\id_M$ and says that eventually every position is zero. Consequently
every original label is paired exactly once, one $A$ factor with one
$B$ factor. This is a bijection of their finite index sets, proving
$m=n$ and the required isomorphisms.

If $M=0$, both composition series have no factors and the conclusion
holds directly.
\end{proof}

For $M\ne0$, termination also has an explicit bound: $E_{m+n}=0$.
Indeed, in a source column $p=n+i$, formula~\eqref{s:source-page}
then uses $F_{p-(m+n)}Y=B_{i-m}=0$. In a target column $j\geq1$,
formula~\eqref{s:target-page} uses
$F_{j+(m+n)-1}X=A_{j+m-1}=M$, so its denominator is all of $B_j$.
These bounds use the lengths of the two given lists, without assuming
that those lengths agree or are intrinsic to $M$.

This proof obtains an actual pairing. An equality of alternating sums
would give less information unless one had already established the
independence of simple classes. No such independence, or Grothendieck
group calculation, is needed here.

\subsection{A calculation with nonzero differentials}

Let $M=\kk e_1\oplus\kk e_2$, with
\[
A_1=\kk e_1,\qquad B_1=\kk e_2,\qquad A_2=B_2=M.
\]
Take $n=2$ in \eqref{s:staggered-filtrations}. Thus the source enters at
columns $3,4$, while the target enters at columns $1,2$:
\[
\begin{array}{c|cccc}
&1&2&3&4\\ \hline
X_1&0&0&\kk e_1&M/\kk e_1\\
Y_1&\kk e_2&M/\kk e_2&0&0.
\end{array}
\]
The only potentially nonzero $d_1$ is
\[
d_{1,3}:\kk e_1\longrightarrow M/\kk e_2,
\qquad ae_1\longmapsto ae_1+\kk e_2.
\]
It is an isomorphism: every target class has such a representative, and
the only vector common to the two coordinate lines is zero. Those two
positions disappear. Page $2$ therefore has just
\[
X_2(4)=M/\kk e_1,\qquad Y_2(1)=\kk e_2.
\]
There is no target in column $2$, so $d_{2,4}=0$; there is no source
in column $3$, so no $d_2$ reaches column $1$. Page $3$ has the same two
surviving factors, but its source formula gives a useful representative:
\[
X_3(4)=\frac{M\cap\kk e_2}{\kk e_1\cap\kk e_2}
       =\kk e_2,
\qquad
Y_3(1)=\frac{\kk e_2}{\kk e_2\cap\kk e_1}=\kk e_2.
\]
The identification of this $X_3(4)$ with $X_2(4)$ sends
$ae_2$ to $ae_2+\kk e_1$. Now
\[
d_{3,4}:\kk e_2\xrightarrow{\id}\kk e_2
\]
removes the remaining pair, and $E_4=0$.
In particular, the zero differential $d_2$ did not mean that the
calculation had stabilized: $d_3$ was still nonzero. Equality of two
successive pages does not by itself imply that every later differential
vanishes.

The cancellations occur at different distances:
\[
\begin{tikzcd}[column sep=large,row sep=large]
\kk e_1\arrow[r,"d_{1,3}","\sim"']&M/\kk e_2
&\text{columns }3\to2\\
\kk e_2\arrow[r,"d_{3,4}","\sim"']&\kk e_2
&\text{columns }4\to1.
\end{tikzcd}
\]
This last picture records the two nonzero arrows on their respective
pages; it is not a single page. The source label $\kk e_2$ in the second row
is the page-$3$ representative just calculated.

For comparison, if $A_1=B_1=\kk e_1$, then $d_1=0$ and the two nonzero
maps occur on page $2$:
\[
d_{2,3}:\kk e_1\iso\kk e_1,
\qquad
d_{2,4}:M/\kk e_1\iso M/\kk e_1.
\]
The page of cancellation records relative filtration position; it is not
an invariant of the unfiltered module. One may also give $M$ the operator
$Te_1=0$, $Te_2=e_2$ and regard it as a $\kk[t]$-module. The same crossed
filtration calculation then pairs the simple type with eigenvalue $0$
on page $1$, and the different simple type with eigenvalue $1$ on page
$3$. The argument matches module types, as well as the number of terms.

\subsection{The butterfly is visible inside the differentials}

The construction has a precise relation to the earlier refinement
rectangle. Take any two finite filtrations $A$ and $B$ of a module $M$;
their factors need not be simple. Use the same shifted identity complex.
For $1\leq i\leq m$ and $1\leq j\leq n$, put
\[
r=n+i-j,\qquad p=n+i.
\]
Then $r\geq1$, and the explicit pages become
\begin{align*}
X_r(n+i)&=\frac{A_i\cap B_j}{A_{i-1}\cap B_j},\\
Y_r(j)&=\frac{B_j}{B_{j-1}+(A_{i-1}\cap B_j)}.
\end{align*}
The differential is induced by the inclusion $A_i\cap B_j\hookrightarrow
B_j$. Its image is canonically
\begin{equation}\label{s:differential-cell}
\im d_{n+i-j,n+i}\cong
Q_{ij}:=
\frac{A_i\cap B_j}
{(A_{i-1}\cap B_j)+(A_i\cap B_{j-1})}.
\end{equation}
Here is the kernel check. If $x\in A_i\cap B_j$ has zero target class,
write $x=b+z$ with $b\in B_{j-1}$ and $z\in A_{i-1}\cap B_j$.
Since both $x$ and $z$ lie in $A_i$, so does $b=x-z$; hence
$b\in A_i\cap B_{j-1}$. The converse is immediate. Thus the kernel of
the map from $A_i\cap B_j$ to its image is exactly the denominator
in \eqref{s:differential-cell}, and the first isomorphism theorem applies.

The same $Q_{ij}$ is what appears on taking associated layers in either
order. To state this precisely, filter $A_i/A_{i-1}$ by
\[
\overline B_j=\frac{A_{i-1}+(A_i\cap B_j)}{A_{i-1}}.
\]
The map $A_i\cap B_j\to\overline B_j/\overline B_{j-1}$ is surjective.
Its kernel is
$(A_{i-1}\cap B_j)+(A_i\cap B_{j-1})$: if
$x=a+y$ with $a\in A_{i-1}$ and $y\in A_i\cap B_{j-1}$, then
$a=x-y\in B_j$; the reverse inclusion is immediate. Consequently
\[
\operatorname{gr}^{B}_j\operatorname{gr}^{A}_i M
\cong Q_{ij}
\cong\operatorname{gr}^{A}_i\operatorname{gr}^{B}_j M.
\]
The second isomorphism follows by the same calculation after interchanging
$A$ and $B$. These are isomorphisms of individual indexed subquotients;
no simultaneous choice of a basis or a splitting of the filtrations has
been made.

Thus the static comparison rectangle and the spectral calculation share
the same subquotients. The rectangle displays all $Q_{ij}$ together.
The spectral sequence encounters them in order of the distance
$r=n+i-j$, as images of actual differentials. The kernel and quotient
calculations proved above explain the connection. Calling the rectangle
alone a spectral sequence would omit the page transitions and the
differentials that we have now supplied.

\subsection{What this adds, and what it does not assert}

The initial framework is a filtered map. Its successive pages separate
what the map already accounts for from what remains. The stable pages
retain the layers of its kernel and cokernel. When the map is an identity,
these stable pages are zero; when the initial layers are simple, the
necessary cancellations become isomorphisms. Jordan--H\"older emerges
from those two facts together.

A single filtration on a module concentrated in degree zero has no
nonzero differentials. Its spectral sequence leaves its own associated
graded object in place and does not by itself compare a second
filtration. The two-term identity complex is what creates a comparison
with nonzero differentials. Likewise, the bare statement that an
alternating sum vanishes would not supply our pairing proof.

With finite filtrations, the stabilized pages describe associated
layers of homology. They do not, without further data, specify how those
layers are attached. That limitation returns us to the questions about
extensions and reconstruction in the following chapters. Spectral
sequences organize successive information; they do not make the extension
problem disappear.

For the general theory, the Stacks Project constructs filtered-complex
spectral sequences and proves finite-filtration convergence in
\href{https://stacks.math.columbia.edu/tag/012K}{Section 12.24,
especially Lemmas 12.24.2 and 12.24.11}. Its convention is cohomological
with decreasing filtrations, whereas ours is homological with increasing
filtrations. The two-filtration interchange is explicitly identified
with Zassenhaus's lemma in Deligne's
\href{https://www.numdam.org/item/10.1007/BF02684692.pdf}{\emph{Th\'eorie
de Hodge II}, \S1.2.1}. These references place the constructions in their
classical setting; the proof in this chapter is given directly, and no
originality claim is made for this presentation.

% END INPUT 06b-module-spectral-calculus.tex

\clearpage\part{What the factors retain and what probes recover}

% BEGIN INPUT 07-additive-invariants.tex
\section{What every additive measurement must retain}
\label{k:section}

The comparison of composition series gives more than an integer.
For each isomorphism class of simple modules $S$, let $m_S(M)$ be
the number of its occurrences in a composition series of $M$.
The common Jordan--H\"older theorem (Theorem~\ref{j:jordan-holder})
makes this well-defined. Only finitely many of these numbers are
nonzero for a fixed finite-length module $M$. The universal monoid
statement was proved in Corollary~\ref{j:universal-monoid}. We now
construct its group-valued recipient explicitly for modules, including
the size and closure checks needed to use it.

\subsection{Splicing gives additivity before we name a group}

The last construction in Section~\ref{m:section} spliced composition
series across a short exact sequence. Uniqueness now turns that
particular construction into an intrinsic assertion:
\begin{equation}\label{k:multiplicity-additivity}
0\to A\to E\to B\to0
\quad\Longrightarrow\quad
m_S(E)=m_S(A)+m_S(B)\quad\hbox{for every simple }S.
\end{equation}
The existence arguments there also show that the category of
finite-length $R$-modules is closed under submodules, quotients and
extensions. Every object subsequently used in such a sequence stays
inside that category.

An additive measurement is an assignment $a(M)$ in an abelian group
$D$, constant on isomorphic modules, such that
$a(E)=a(A)+a(B)$ for every short exact sequence of finite-length
modules. The zero sequence forces $a(0)=0$. Repeatedly using the
short exact sequences of a composition series gives
\begin{equation}\label{k:all-measurements}
 a(M)=\sum_S m_S(M)a(S).
\end{equation}
Thus a choice of values on simple modules determines every additive
measurement. Conversely any such choice defines an additive measurement
by this formula and \eqref{k:multiplicity-additivity}. This already
answers the universal question in concrete terms.

\subsection{Constructing the universal recipient}

For completeness we construct the group that receives every additive
measurement. Take one formal generator $e_M$ for each isomorphism
class of finite-length $R$-modules. These classes form a set: every
simple module is cyclic, since the map $R\to S$, $r\mapsto rx$, has
nonzero image for $x\ne0$; inductively every finite-length module is
finitely generated by generators of a submodule and lifts of generators
of a quotient. Such a module is a quotient of some $R^n$, and the
submodules of all these $R^n$ form a set.

Let $F$ be the abelian group of finitely supported integer combinations
of the symbols $e_M$, with coefficientwise addition. Let $J$ be the
subgroup generated by
\[
 e_E-e_A-e_B
 \quad\text{for all }0\to A\to E\to B\to0.
\]
Define $K_0(R\text{-}\mathrm{fl})=F/J$, and write $[M]$ for the
image of $e_M$. Here $R\text{-}\mathrm{fl}$ means finite-length
left $R$-modules. We specify that category because enlarging it can
introduce new exact sequences and therefore new relations.

Given an additive measurement $a$ into $D$, the formula
$\sum n_Me_M\mapsto\sum n_Ma(M)$ is a homomorphism $F\to D$.
It kills $J$ exactly by additivity, so the quotient universal property
gives a unique homomorphism $\bar a:K_0\to D$ with $\bar a([M])=a(M)$.
Conversely such a homomorphism defines an additive measurement.
This is the universal property of the Grothendieck group.

\begin{theorem}\label{k:free-simples}
The homomorphism
\[
 \bigoplus_{[S]\text{ simple}}\mathbb Z e_S\longrightarrow
 K_0(R\text{-}\mathrm{fl}),\qquad e_S\longmapsto[S]
\]
is an isomorphism. Its inverse sends $[M]$ to
$\sum_Sm_S(M)e_S$.
\end{theorem}
\begin{proof}
Additivity of every $m_S$ shows that the proposed inverse kills each
defining relation and is therefore well-defined. A simple module has
one factor, itself, so the composite on each $e_S$ is the identity.
For the other composite, the defining relations applied to a composition
series give $[M]=\sum_Sm_S(M)[S]$. Since the classes $[M]$ generate
$K_0$, this composite too is the identity.
\end{proof}

The direct sum in the theorem means finite support, even if there are
infinitely many simple isomorphism classes. The theorem proves exactly
which information all additive measurements retain. It says nothing
about measurements that fail additivity on short exact sequences.

Length is the homomorphism sending every $[S]$ to $1$. For
finite-dimensional vector spaces there is one simple class $[k]$,
so $K_0\cong\mathbb Z$ and $[V]=(\dim_kV)[k]$.
For finite-dimensional modules over a $k$-algebra, dimension instead
assigns the weight $\dim_k S$ to a simple class. To use the same proof
in that category, observe that submodules and quotients are finite-dimensional
by Lemma~\ref{m:bases}, extensions have a basis made from a submodule
basis and quotient lifts, and all these modules have composition series
by Section~\ref{m:length-foundations}.

\subsection{Why the chosen category matters}

Let $S=k[t]/(t)$. In the category of finite-length $k[t]$-modules,
$[S]$ is a nonzero generator of the free group in the theorem.
If we instead impose exact-sequence relations for all finitely generated
$k[t]$-modules, the sequence
\[
0\longrightarrow k[t]\xrightarrow{\,t\,}k[t]
\longrightarrow S\longrightarrow0
\]
forces $[S]=0$. Multiplication by $t$ is injective by inspection of
coefficients, and its image is exactly the polynomials with zero
constant term, proving exactness. The module $k[t]$ does not have
finite length: the strictly descending submodules $(t^n)$ would each
force a strictly smaller nonnegative length if it did, by
\eqref{k:multiplicity-additivity}. Thus this new relation was not
available in the previous category.

The invariant depends on which objects and exact sequences we permit.
Keeping that context visible prevents a numerical formula from
silently changing its meaning.

\subsection{A short bridge to Euler characteristic}
\label{k:euler}

Let $C^\bullet$ be a bounded complex of finite-length modules: it
has maps $d^i:C^i\to C^{i+1}$ with $d^{i+1}d^i=0$, and only
finitely many $C^i$ are nonzero. Put
\[
 Z^i=\Ker d^i,\quad B^i=\im d^{i-1},\quad H^i=Z^i/B^i.
\]
The complex condition gives $B^i\subseteq Z^i$. Kernels, images and
quotients are finite-length by the results already proved. The two
short exact sequences
\[
0\to Z^i\to C^i\xrightarrow{d^i}B^{i+1}\to0,
\qquad 0\to B^i\to Z^i\to H^i\to0
\]
give $[C^i]=[B^i]+[H^i]+[B^{i+1}]$. Alternating addition cancels
every boundary class with its next occurrence. Boundedness makes all
the sums finite, and proves
\begin{equation}\label{k:euler-formula}
 \sum_i(-1)^i[C^i]=\sum_i(-1)^i[H^i].
\end{equation}
Applying any additive measurement gives its numerical version.
In particular, vector-space dimension gives the usual Euler
characteristic formula.

This identity follows directly from the exact-sequence relations.
We have introduced it only after the earlier comparison. In
particular, we did not infer Jordan--H\"older from
an equality in a Grothendieck group whose simple generators had
already been justified by Jordan--H\"older. The preceding spectral
argument compared actual simple modules and actual differentials.

\subsection{The information that remains to be recovered}

If $0\to A\to E\to B\to0$ is exact, its class is always
$[A]+[B]$, whether or not the sequence splits. Every additive
measurement therefore agrees on $E$ and $A\oplus B$.
This is an exact limitation of additivity. The next chapter retains
the maps of the extension, and measures how the middle object can
depend on them.

% END INPUT 07-additive-invariants.tex


% BEGIN INPUT 08-extensions-and-lifting.tex
\section{What attaches one layer to another}\label{x:extensions}

The additive invariant has reached its limit. For finite-length modules,
an exact sequence
\[
0\longrightarrow A\xrightarrow{i}E\xrightarrow{q}B\longrightarrow0
\]
always gives $[E]=[A]+[B]$, whether or not $E$ is a direct sum. To learn
more, we must retain the way $A$ sits in $E$ and the way $E$ maps onto $B$.
The maps are part of the object we now study.

Throughout this section $R$ is a ring with identity, modules are left
$R$-modules, and every map between modules is $R$-linear. Direct sums carry
componentwise operations. An \emph{extension of $B$ by $A$} means the whole
displayed short exact sequence, with $A$ and $B$ fixed. We will first
construct extensions, then calculate a family of them, and finally use
them to detect when a map cannot be lifted.

\subsection{Keeping the endpoints fixed}

Two extensions $(i,E,q)$ and $(i',E',q')$ are \emph{equivalent} if there is
a map $h:E\to E'$ making this diagram commute:
\[
\begin{tikzcd}[column sep=large]
0\arrow[r]&A\arrow[r,"i"]\arrow[d,equal]
 &E\arrow[r,"q"]\arrow[d,"h"]&B\arrow[r]\arrow[d,equal]&0\\
0\arrow[r]&A\arrow[r,"i'"]&E'\arrow[r,"q'"]&B\arrow[r]&0.
\end{tikzcd}
\]
Thus $hi=i'$ and $q'h=q$. No separate invertibility assumption on $h$ is
needed.

\begin{lemma}\label{x:marked-iso}
Every such $h$ is an isomorphism.
\end{lemma}
\begin{proof}
If $h(e)=0$, then $q(e)=q'h(e)=0$, so $e=i(a)$ for some $a\in A$.
Now $0=hi(a)=i'(a)$, and injectivity of $i'$ gives $a=0$ and $e=0$.
Hence $h$ is injective.

Given $e'\in E'$, choose $e\in E$ with $q(e)=q'(e')$. This is possible
because $q$ is surjective. The difference $e'-h(e)$ lies in $\Ker q'$, so
it is $i'(a)$ for some $a\in A$. Then $h(e+i(a))=e'$. Thus $h$ is
surjective. The inverse of a bijective linear map is linear: apply $h$ to
the proposed addition and scalar identities and use its injectivity.
\end{proof}

Identity maps, compositions, and inverses now show that equivalence is
indeed an equivalence relation. The word \emph{fixed} matters: an
isomorphism between the middle modules that changes an endpoint
identification need not give an equivalence of extensions.

The \emph{split extension} is
\[
0\longrightarrow A\xrightarrow{a\mapsto(a,0)}A\oplus B
 \xrightarrow{(a,b)\mapsto b}B\longrightarrow0.
\]
An extension is called split when it is equivalent to this one.

\begin{proposition}\label{x:splitting}
An extension is split if and only if there is a linear map $s:B\to E$
with $qs=\id_B$.
\end{proposition}
\begin{proof}
Such a map is a \emph{section} of $q$. Given a section, the map
\[
A\oplus B\longrightarrow E,\qquad (a,b)\longmapsto i(a)+s(b)
\]
commutes with both endpoint maps and is therefore an isomorphism by
Lemma~\ref{x:marked-iso}. Conversely, an equivalence with the split
extension transports its section $b\mapsto(0,b)$ to a section of $q$.
\end{proof}

A choice of a lift of each individual element of $B$ is insufficient.
The lifts must form a linear map. Failure of that compatibility is
precisely what a nonsplit extension records.

\subsection{Change the quotient: pull back}

Suppose we want to compare an extension with a map $f:P\to B$. We need
an object that simultaneously records an element of $E$ and an element
of $P$ which have the same image in $B$. Define
\[
f^*E=E\times_B P
 =\{(e,p)\in E\oplus P:q(e)=f(p)\}.
\]
This is a submodule: it is the kernel of the linear map
$(e,p)\mapsto q(e)-f(p)$. Put
\[
j(a)=(i(a),0),\qquad \pi(e,p)=p,\qquad u(e,p)=e.
\]
These maps form the diagram
\[
\begin{tikzcd}[column sep=large]
0\arrow[r]&A\arrow[r,"j"]\arrow[d,equal]
 &f^*E\arrow[r,"\pi"]\arrow[d,"u"]&P\arrow[r]\arrow[d,"f"]&0\\
0\arrow[r]&A\arrow[r,"i"]&E\arrow[r,"q"]&B\arrow[r]&0.
\end{tikzcd}
\]
The left square commutes because $uj=i$. The right one commutes because
membership in $f^*E$ says exactly that $qu=f\pi$.

The top row is exact. The map $j$ is injective because $i$ is. If
$\pi(e,p)=0$, then $p=0$ and $q(e)=0$, so $(e,p)=(i(a),0)=j(a)$
for a unique $a$. Conversely every $j(a)$ maps to zero. Finally, for
any $p\in P$, surjectivity of $q$ supplies $e$ with $q(e)=f(p)$, so
$\pi(e,p)=p$. No assumption that $f$ is injective or surjective was used.

The right square is a pullback in the precise universal sense. If maps
$r:T\to E$ and $s:T\to P$ satisfy $qr=fs$, the unique map to $f^*E$
with coordinates $r$ and $s$ is
\[
T\longrightarrow f^*E,\qquad t\longmapsto(r(t),s(t)).
\]
The compatibility puts its image in $f^*E$; coordinatewise operations
make it linear; its two coordinates force uniqueness.

\subsection{Change the subobject: push out}

Now suppose instead that $g:A\to A'$ is given. We want to replace the
embedded $A$ by $A'$ while preserving the quotient $B$. In $A'\oplus E$
we therefore identify $(g(a),0)$ with $(0,i(a))$. The required relation
submodule is
\[
L_g=\{(g(a),-i(a)):a\in A\},
\qquad
g_*E=(A'\oplus E)/L_g.
\]
Here $L_g$ is the image of a linear map, so it is a submodule and the
quotient is defined. Brackets below denote classes modulo $L_g$. Define
\[
i_g(a')=[a',0],\qquad u_g(e)=[0,e],\qquad q_g([a',e])=q(e).
\]
The last map is well-defined because $q(i(a))=0$. The diagram is
\[
\begin{tikzcd}[column sep=large]
0\arrow[r]&A\arrow[r,"i"]\arrow[d,"g"]
 &E\arrow[r,"q"]\arrow[d,"u_g"]&B\arrow[r]\arrow[d,equal]&0\\
0\arrow[r]&A'\arrow[r,"i_g"]&g_*E\arrow[r,"q_g"]&B\arrow[r]&0.
\end{tikzcd}
\]
The left square commutes because $[0,i(a)]=[g(a),0]$. The right square
commutes by the formula for $q_g$.

The bottom row is exact. If $[a',0]=0$, then
$(a',0)=(g(a),-i(a))$ for some $a$; injectivity of $i$ gives $a=0$
and then $a'=0$. If $q_g([a',e])=0$, write $e=i(a)$. The relation gives
\[
[a',e]=[a',i(a)]=[a'+g(a),0],
\]
so the kernel is the image of $i_g$. That image is visibly killed by
$q_g$, and $q_g$ is surjective because $q$ is. Again no injectivity or
surjectivity of $g$ was required.

The left square is a pushout. Given maps $r:A'\to T$ and $s:E\to T$
with $rg=si$, the formula
\[
[a',e]\longmapsto r(a')+s(e)
\]
defines a linear map $g_*E\to T$: the map on the direct sum kills
$(g(a),-i(a))$. It has the prescribed composites with $i_g$ and $u_g$.
Uniqueness follows because every class is $[a',0]+[0,e]$.

These two constructions have opposite directions for a reason. A map
\emph{into} the quotient permits us to restrict the extension over a new
domain. A map \emph{out of} the subobject permits us to impose a new
identification in the kernel. Their universal properties specify exactly
which compatible maps survive the construction.

\subsection{Adding extensions by merging their kernels}

Suppose
\[
\xi_j:\quad 0\longrightarrow A\xrightarrow{i_j}E_j
 \xrightarrow{q_j}B\longrightarrow0
\qquad (j=1,2)
\]
are two extensions with the same endpoints. Taking their direct sum gives
two copies of $B$ and two copies of $A$. We first require the quotient
coordinates to agree, and then merge the kernel coordinates by addition.
Thus we pull back along the diagonal $B\to B\oplus B$ and push out along
the sum $A\oplus A\to A$.

Here is the resulting object in one formula. Let
\[
T_2=\{(e_1,e_2):q_1(e_1)=q_2(e_2)\},\qquad
L_2=\{(i_1(a),-i_2(a)):a\in A\},
\]
and set $Q_2=T_2/L_2$. Both are submodules of the relevant direct sum,
and $L_2\subseteq T_2$. Define
\[
i_+(a)=[i_1(a),0]=[0,i_2(a)],\qquad
q_+([e_1,e_2])=q_1(e_1)=q_2(e_2).
\]
The equality defining $i_+$ is exactly the relation imposed by $L_2$.
The map $i_+$ is injective: if $(i_1(a),0)=(i_1(c),-i_2(c))$, then
$c=0$ and $a=0$. Every pair in the kernel of $q_+$ has the form
$(i_1(a),i_2(b))$, and its class is $i_+(a+b)$. Surjectivity follows by
lifting a given $b\in B$ through both $q_1$ and $q_2$. Hence
\[
\xi_1+\xi_2:\quad
0\longrightarrow A\xrightarrow{i_+}Q_2\xrightarrow{q_+}B
\longrightarrow0
\]
is an extension. This operation is called the \emph{Baer sum}.

To check the asserted pullback--pushout description, the pullback is
$T_2$, with kernel $A\oplus A$. Its pushout along addition is
\[
\frac{A\oplus T_2}
 {\{(a+b,(-i_1(a),-i_2(b))):a,b\in A\}}.
\]
It maps to $Q_2$ by
\[
[c,(e_1,e_2)]\longmapsto[e_1+i_1(c),e_2].
\]
A displayed relation maps to $[i_1(b),-i_2(b)]=0$. Its inverse sends
$[e_1,e_2]$ to $[0,(e_1,e_2)]$; the relation defining $L_2$ maps to a
pushout relation with $b=-a$. The two composites are identities, since
$[c,(e_1,e_2)]=[0,(e_1+i_1(c),e_2)]$ in the pushout. This verifies the
construction, including its endpoint maps.

\begin{proposition}\label{x:baer-group}
Baer sum makes the equivalence classes of extensions of $B$ by $A$ into
an abelian group. Its zero is the split extension. An inverse of
$(i,E,q)$ is $(-i,E,q)$.
\end{proposition}
\begin{proof}
First the operation is well-defined on classes. Equivalences
$h_j:E_j\to E'_j$ send $T_2$ into $T'_2$ and $L_2$ into $L'_2$
coordinatewise. They induce a map of extensions, which is an isomorphism
by Lemma~\ref{x:marked-iso}.

For associativity, form the common-image tuples for any $n\geq1$:
\[
T_n=\{(e_1,\ldots,e_n):q_1(e_1)=\cdots=q_n(e_n)\},
\]
and the relation submodule
\[
L_n=\{(i_1(a_1),\ldots,i_n(a_n)):a_1+\cdots+a_n=0\}.
\]
The quotient $Q_n=T_n/L_n$ is an extension of $B$ by $A$. Its quotient
map takes the common image; its kernel map sends $a$ to
$[i_1(a),0,\ldots,0]$. This kernel map is injective, because a tuple of
that form belongs to $L_n$ only if $a=0$. A tuple in the kernel of the
quotient map has entries $i_j(a_j)$; modulo $L_n$ it equals the kernel
image of $\sum_j a_j$. Finally, a given element of $B$ can be lifted in
each coordinate. These statements prove exactness.

In the following calculation, $E_1+E_2$ denotes the middle module of
the Baer sum, with its displayed endpoint maps.
For $n=3$ the map
\[
T_3\longrightarrow (E_1+E_2)+E_3,
\qquad (e_1,e_2,e_3)\longmapsto[[e_1,e_2],e_3]
\]
is linear. If $a_1+a_2+a_3=0$, its value on
$(i_1(a_1),i_2(a_2),i_3(a_3))$ is the class of
$(i_{12}(a_1+a_2),i_3(a_3))$, which is zero by the outer sum relation.
Here $i_{12}$ denotes the kernel inclusion for $E_1+E_2$.
Thus it factors through $Q_3$. The induced map is the identity on the
fixed endpoints, so it is an equivalence. The map
\[
(e_1,e_2,e_3)\longmapsto[e_1,[e_2,e_3]]
\]
likewise identifies $Q_3$ with $E_1+(E_2+E_3)$: a kernel tuple maps to
the relation determined by $a_1+(a_2+a_3)=0$. Both bracketings therefore
represent the same extension class. This proves associativity without
assuming it for the new operation.

Swapping the two coordinates identifies $E_1+E_2$ with $E_2+E_1$.
It preserves the kernel inclusion because
$[i_1(a),0]=[0,i_2(a)]$, and preserves the common quotient image.
This proves commutativity.

For the split extension $S=A\oplus B$, the formula
\[
E+S\longrightarrow E,\qquad [e,(a,b)]\longmapsto e+i(a)
\]
is defined on pairs with $q(e)=b$. The relation
$(i(c),(-c,0))$ maps to zero, and both endpoint maps are preserved.
Lemma~\ref{x:marked-iso} makes it an equivalence, proving the identity law.

Finally, in $E+(-E)$, where $-E$ denotes $(-i,E,q)$, the relation
submodule consists of pairs $(i(a),i(a))$. The linear map
$e\mapsto[e,e]$ kills $i(A)$, so it factors uniquely through $E/i(A)$.
The quotient map identifies $E/i(A)$ with $B$; hence we obtain a linear
map $s:B\to E+(-E)$ with $q_+s=\id_B$. By
Proposition~\ref{x:splitting}, this sum is split. The inverse law follows.
\end{proof}

These equivalence classes form a set. Indeed, choosing a set-theoretic
lift of every element of $B$ identifies the underlying set of a middle
module with $A\times B$, by $(a,b)\mapsto i(a)+s(b)$; this uses only
a choice of representatives, not a linear section. Module operations
and endpoint maps on this fixed set range over sets of functions.
Thus the extension classes can be represented by a set of such data.
For arbitrary modules this set-theoretic observation uses the usual
axiom of choice; none of the finite-filtration comparison proofs needs it.

We denote this abelian group by $\Ext^1_R(B,A)$. The order of the
arguments records the sequence: $A$ is the kernel and $B$ the quotient.
The superscript $1$ refers to a single extension step. No higher
extension theory is needed for what follows. These are classical
constructions; the notation and the general pullback--pushout description
are also recorded in the \href{https://stacks.math.columbia.edu/tag/010I}
{Stacks Project, Section 12.6}.

The new associativity has a concrete source. Three extensions can be
added by retaining triples with a common image and then imposing the
single condition that the three kernel coordinates sum to zero. Either
pairwise construction gives that same object up to the specified
isomorphism. The quotient has removed the dependence on how the addition
was parenthesized.

\subsection{An extension that fits in one matrix}\label{x:operator-extensions}

Let $k$ be a field. For $\lambda\in k$, let $S_\lambda$ be the
one-dimensional $k[t]$-module on which $t$ acts by multiplication by
$\lambda$. It is simple because a nonzero vector spans its underlying
one-dimensional space. Consider
\[
0\longrightarrow S_\mu\xrightarrow{i}E\xrightarrow{q}S_\lambda
\longrightarrow0.
\]
Write $T$ for the action of $t$. Choose $u=i(1)$ and choose $v\in E$
with $q(v)=1$. These form a $k$-basis: every $e\in E$ can first have its
quotient coordinate removed by subtracting $q(e)v$, leaving an element
of $i(S_\mu)$; and applying $q$ to $cu+dv=0$ gives $d=0$, followed by
$c=0$ since $i$ is injective.

The endpoint maps commute with the operators, so
\[
Tu=\mu u,\qquad Tv=\lambda v+a u
\]
for a unique $a\in k$. In this basis,
\[
[T]_{(u,v)}=\begin{pmatrix}\mu&a\\0&\lambda\end{pmatrix}.
\]
Conversely, for any $a$, these formulas define a linear operator on
$ku\oplus kv$, hence a $k[t]$-module $E_a$, and an extension with
$i(1)=u$, $q(u)=0$, and $q(v)=1$. Thus every extension under discussion
has this form.

The choice of $u$ is fixed by the marked endpoint, whereas another lift
of $1$ has the form $v'=v+bu$. Its coefficient is found by one calculation:
\[
T(v+bu)=\lambda(v+bu)+\bigl(a+b(\mu-\lambda)\bigr)u.
\]
The change of lift acts on the attachment coefficient by
\[
a\longmapsto a+b(\mu-\lambda).
\]

\begin{proposition}\label{x:operator-ext-classification}
The preceding coefficient gives an isomorphism of abelian groups
\[
\Ext^1_{k[t]}(S_\lambda,S_\mu)
 \cong k/(\mu-\lambda)k.
\]
In particular, extensions split when $\lambda\ne\mu$. When
$\lambda=\mu$, each $a\in k$ gives a distinct marked extension class,
and its class is zero exactly when $a=0$.
\end{proposition}
\begin{proof}
Changing the lift changes $a$ by an element of $(\mu-\lambda)k$, so
its coset is independent of that choice. An equivalence fixes $u$ and
sends a lift of $1$ to a lift of $1$; thus it preserves the same coset.
Conversely, if $a'=a+b(\mu-\lambda)$, the map $E_{a'}\to E_a$ sending
$u'\mapsto u$ and $v'\mapsto v+bu$ commutes with $T$ by the displayed
calculation. It therefore commutes with every polynomial in $T$, is
$k[t]$-linear, and is an equivalence of extensions. The classification
is consequently bijective.

To compute its addition, take $E_a$ and $E_c$ with chosen vectors
$u_1,v_1$ and $u_2,v_2$. In their Baer sum put
\[
u_+=[u_1,0]=[0,u_2],\qquad v_+=[v_1,v_2].
\]
The first vector is the marked kernel image of $1$ and the second maps
to $1$ in the quotient. The same spanning and independence argument
used above makes them a basis. Their operator satisfies
\[
Tu_+=\mu u_+,\qquad
Tv_+=\lambda v_++[a u_1,c u_2]
       =\lambda v_++(a+c)u_+.
\]
Thus Baer sum corresponds to addition of coefficient classes.

The split extension has coefficient zero. If $\mu-\lambda\ne0$, it
is invertible in $k$, so $(\mu-\lambda)k=k$ and the quotient group is
zero. If $\mu=\lambda$, the ideal is zero, giving the stated result.
\end{proof}

When the eigenvalues differ, the explicit section can be read off as
well. Choose $b=-a/(\mu-\lambda)$. Then $v'=v+bu$ satisfies
$Tv'=\lambda v'$, and the map $S_\lambda\to E$, $c\mapsto cv'$, is
$k[t]$-linear and is a section of $q$.

When the eigenvalues agree, a nonzero $a$ cannot be removed by changing
the lift. Yet all nonzero $a$ give isomorphic \emph{middle modules}.
Indeed, put $w=au$; then
\[
Tw=\lambda w,\qquad Tv=\lambda v+w,
\]
which has the same matrix for every nonzero $a$. This rescaling changes
the marked kernel vector, so it is not an equivalence keeping the
endpoints fixed. The $a=0$ middle module is different: $T-\lambda\id$
is zero there, whereas it is nonzero for $a\ne0$, a property preserved
by every intertwining isomorphism.

For $\lambda=\mu=1$, the cases $a=0$ and $a=1$ give the identity
operator and the triangular shear $(x,y)\mapsto(x+y,y)$. Both layers carry the
identity action. The coefficient $a$ now records explicitly the
attachment that their common list of layers forgot.

There are therefore two different classification questions. One asks
for the middle operator up to change of basis. The other remembers a
specified subobject and a specified quotient map. The latter contains
more information, which is why it supports the extension group just
constructed.

\subsection{A boundary map measures failure to lift}\label{x:obstruction}

Return to an arbitrary extension
\[
\xi:\quad0\longrightarrow A\xrightarrow{i}E\xrightarrow{q}B
\longrightarrow0.
\]
Given $f:P\to B$, we can now form a precise obstruction to finding
$\widetilde f:P\to E$ with $q\widetilde f=f$: take the pullback
extension and its class. Define
\[
\delta_\xi:\Hom_R(P,B)\longrightarrow\Ext^1_R(P,A),
\qquad f\longmapsto[f^*E].
\]
If $\widetilde f$ exists, then $p\mapsto(\widetilde f(p),p)$ is a
section of the pullback extension. Conversely, a section of its
projection to $P$ has exactly this form, and its first coordinate is
a lift. Proposition~\ref{x:splitting} therefore proves
\[
\delta_\xi(f)=0
\quad\Longleftrightarrow\quad
f\text{ has a linear lift through }q.
\]

This is a homomorphism of abelian groups. To verify the assertion,
take $f,g:P\to B$. The map
\[
(f^*E)+(g^*E)\longrightarrow(f+g)^*E,
\qquad[(e_1,p),(e_2,p)]\longmapsto(e_1+e_2,p)
\]
is well-defined: the target condition is
$q(e_1+e_2)=f(p)+g(p)$, and a Baer relation
$((i(a),0),(-i(a),0))$ maps to zero. It is linear and preserves both
endpoint maps. Lemma~\ref{x:marked-iso} makes it an equivalence, proving
$\delta_\xi(f+g)=\delta_\xi(f)+\delta_\xi(g)$. The pullback along
the zero map has section $p\mapsto(0,p)$, so $\delta_\xi(0)=0$.

In particular, we have proved exactness of the sequence of abelian groups
\[
0\longrightarrow\Hom_R(P,A)\xrightarrow{i\circ-}\Hom_R(P,E)
\xrightarrow{q\circ-}\Hom_R(P,B)
\xrightarrow{\delta_\xi}\Ext^1_R(P,A).
\]
Here exactness is asserted at the three $\Hom$ groups. The first map
is injective because $i$ is injective. A map $h:P\to E$ is killed by
$q\circ-$ exactly when its image lies in $i(A)$; it then factors
uniquely as $h=i\bar h$, with $\bar h$ linear because $i$ is an
injective linear map. Exactness at $\Hom_R(P,B)$ is the lifting
criterion proved above. No surjectivity onto $\Ext^1_R(P,A)$ is claimed.

For example, if $P=B$ and $f=\id_B$, its pullback is equivalent to
the original extension by $e\mapsto(e,q(e))$. Thus
$\delta_\xi(\id_B)=[\xi]$: lifting the identity is exactly splitting
the extension. In the equal-eigenvalue example, this class is the
coefficient $a$. A section would require an eigenvector above $1$;
the calculation shows exactly when such a vector exists.

There is also a useful contrast. If $P=R$, every map $f:R\to B$ lifts:
choose $e$ with $q(e)=f(1)$ and define $\widetilde f(r)=re$. This is
linear and satisfies $q\widetilde f(r)=r f(1)=f(r)$. The obstruction
depends on the object from which we test the quotient. A test object
with relations can detect constraints that a single free generator
cannot.

\subsection{An attachment can skip an adjacent layer}\label{x:skip-layer}

It is tempting to recover a filtered object by listing its simple
layers and the extension between every adjacent pair. Even that list
can miss information. The following example identifies the loss without
using any classification theorem.

On $k^3$ fix the flag
\[
0=F_0\subset F_1=ke_1\subset F_2=ke_1+ke_2\subset F_3=k^3.
\]
Compare the zero operator with the operator $N$ given by
\[
N(e_1)=N(e_2)=0,\qquad N(e_3)=e_1.
\]
Both preserve every $F_i$, so these are filtrations by $k[t]$-submodules.
Each of the three successive one-dimensional quotients has zero induced
operator in both examples, hence is $S_0$.

Now retain the adjacent two-layer pieces as well. On $F_2/F_0=F_2$
the operator is zero in both examples. On $F_3/F_1$ it is also zero,
because $N(F_3)\subseteq F_1$. Thus the marked extensions
\[
0\to F_1/F_0\to F_2/F_0\to F_2/F_1\to0,
\qquad
0\to F_2/F_1\to F_3/F_1\to F_3/F_2\to0
\]
are split and agree, using the indicated coordinate identifications,
for the two operators. Nevertheless the full modules are not
isomorphic. An isomorphism intertwining $N$ with zero would satisfy
$hN=0$; injectivity of $h$ would force $N=0$, contrary to $N(e_3)=e_1$.

The map from the third layer into the first is invisible in either
adjacent two-layer quotient. It remains visible if we retain the whole
extension
\[
0\longrightarrow F_2\longrightarrow F_3\longrightarrow F_3/F_2
\longrightarrow0
\]
with its maps: in the $N$ example every lift of the quotient generator
has the form $e_3+x$ with $x\in F_2$, and $N(e_3+x)=e_1\ne0$,
so no $k[t]$-linear section exists. For the zero operator the evident
section does exist.

The example specifies what the proposed summary discarded. Complete
successive extension data includes each already assembled module and
the maps attaching the next layer. Replacing those data by isolated
adjacent two-layer summaries can erase a longer attachment. We now
seek observations that retain enough information to reconstruct all
finite-dimensional nilpotent operators.

% END INPUT 08-extensions-and-lifting.tex


% BEGIN INPUT 09-probes-and-reconstruction.tex
\section{Recovering a nilpotent operator from spaces of maps}
\label{p:probes}

Composition factors count the simple layers, and extensions record ways to
attach them. We now ask a reconstruction question with a complete answer:
can a suitable collection of maps into an object determine its isomorphism
class?

Our objects will be finite-dimensional vector spaces over an arbitrary field
$k$, equipped with a nilpotent operator. The answer will use only kernels,
images, quotients, and finite bases. In particular, we will prove the required
decomposition into chains. No classification of modules over a polynomial
ring, or Jordan normal form theorem, is assumed.

\subsection{A relation becomes a test object}
\label{p:relation-probe}

Let $V$ be a finite-dimensional $k$-vector space and let $N:V\to V$ be linear.
Recall how $N$ gives $V$ the structure of a $k[t]$-module $M$:
\[
 \left(\sum_i a_it^i\right)\!\cdot v=\sum_i a_iN^iv,
 \qquad N^0=\id_V.
\]
The identity $N^iN^j=N^{i+j}$ ensures that multiplication of polynomials
agrees with composition of their actions. A linear map $F:V\to V'$ is a
$k[t]$-module map precisely when $FN=N'F$: necessity follows by considering
$t$, and sufficiency follows by induction for powers of $t$ and then by
linearity for polynomials.

Assume henceforth that $N^R=0$ for some integer $R\geq1$. This is the
meaning of nilpotence. We allow $V=0$.

For every $r\geq1$, introduce the module
\[
 P_r=k[t]/(t^r).
\]
Its elements have unique representatives
$a_0+a_1t+\cdots+a_{r-1}t^{r-1}$: remove all terms divisible by $t^r$, and
observe that a polynomial of degree less than $r$ divisible by $t^r$ is zero.
Thus $1,t,\ldots,t^{r-1}$, understood as residue classes, form a $k$-basis.
Multiplication by $t$ runs along this basis and then gives zero.

The module $P_r$ has one generator with one relation, $t^r\cdot1=0$.
A map out of $P_r$ asks whether a proposed image of the generator satisfies
that relation.

\begin{proposition}[A kernel represented by a module]
\label{p:hom-kernel}
Evaluation at $1$ gives a $k$-linear isomorphism
\[
 \operatorname{ev}_1:\Hom_{k[t]}(P_r,M)\iso\Ker N^r,
 \qquad f\longmapsto f(1).
\]
Its inverse sends $v\in\Ker N^r$ to the map
\[
 f_v:P_r\longrightarrow M,\qquad \overline{p(t)}\longmapsto p(N)v.
\]
\end{proposition}

\begin{proof}
If $f$ is a module map, then
$N^rf(1)=f(t^r\cdot1)=0$. Moreover,
$f(\overline{p(t)})=p(N)f(1)$, so $f(1)$ determines $f$.

Conversely, if $N^rv=0$ and $p-p'=t^rq$, then
$(p-p')(N)v=q(N)N^rv=0$. Thus $f_v$ is well-defined on residue classes.
It is additive and $k$-linear, and
\[
 f_v(\overline{a(t)p(t)})=a(N)p(N)v
 =a(t)\cdot f_v(\overline{p(t)}),
\]
so it is a module map. The two constructions are inverse. Addition and
scalar multiplication of module maps are defined pointwise; the module
identities verify that these remain module maps. Evaluation and its inverse
respect these operations.
\end{proof}

This isomorphism respects maps between the objects being measured. If
$F:M\to M'$ is a module map, then $FN^r=(N')^rF$, so $F$ sends
$\Ker N^r$ into $\Ker(N')^r$. The square
\[
\begin{tikzcd}[column sep=large,row sep=large]
\Hom_{k[t]}(P_r,M) \arrow[r,"\operatorname{ev}_1","\sim"']
  \arrow[d,"f\mapsto F\circ f"']
  & \Ker N^r \arrow[d,"F|_{\Ker N^r}"] \\
\Hom_{k[t]}(P_r,M') \arrow[r,"\operatorname{ev}_1","\sim"']
  & \Ker (N')^r
\end{tikzcd}
\]
commutes because both paths send $f$ to $F(f(1))$. This compatibility is what
it means here for the identification to be \emph{natural} in $M$.

There is also compatibility as we enlarge the test object. The quotient
map $q_r:P_{r+1}\to P_r$ sends $\overline p$ to $\overline p$; it is
well-defined because $(t^{r+1})\subseteq(t^r)$. Precomposition by $q_r$
corresponds, under evaluation at $1$, to the inclusion
$\Ker N^r\subseteq\Ker N^{r+1}$. The growing sequence of kernels is thus
also a growing sequence of spaces of maps from explicit test objects.

\subsection{Which endpoints can be reached by a long chain?}
\label{p:bottom-spaces}

Set
\[
 K_r=\Ker N^r,\qquad d_r=\dim K_r\quad(r\geq0).
\]
Since $N^0=\id_V$, $K_0=0$ and $d_0=0$. For $r\geq R$, we have
$K_r=V$. The numbers $d_r$ are now well-defined by the dimension theorem
already proved. By Proposition~\ref{p:hom-kernel}, they are precisely the
dimensions of the spaces of maps from the $P_r$.

To understand these numbers, look at a second filtration, this time inside
$\Ker N$:
\[
 B_r=\Ker N\cap\im N^{r-1}\qquad(r\geq1).
\]
Every $B_r$ is a subspace, and
\[
 \Ker N=B_1\supseteq B_2\supseteq\cdots
 \supseteq B_R\supseteq B_{R+1}=0.
\]
An element of $B_r$ is killed by $N$ and can be reached by applying $N$
$r-1$ times. It can therefore be the endpoint of a chain of length $r$.
Membership in $B_{r+1}$ says that it can be reached by an even longer chain.

There is a short exact sequence of vector spaces
\begin{equation}
\label{p:kernel-bottom-ses}
 0\longrightarrow K_{r-1}\longrightarrow K_r
 \xrightarrow{\ N^{r-1}\ } B_r\longrightarrow0.
\end{equation}
The middle map has the stated target because $N^rx=0$ for $x\in K_r$.
Its kernel is $K_{r-1}$. It is surjective: if $y\in B_r$, write
$y=N^{r-1}x$; then $0=Ny=N^rx$, so $x\in K_r$. The first arrow is
inclusion, and the end arrows are the unique maps to or from the zero space.

Finite-dimensional dimension additivity applied to this sequence gives
\begin{equation}
\label{p:bottom-dim}
 \dim B_r=d_r-d_{r-1}.
\end{equation}
Consequently the successive increments in the sequence $(d_r)$ are
nonnegative and nonincreasing. Define
\begin{equation}
\label{p:canonical-multiplicity}
 E_r=B_r/B_{r+1},\qquad b_r=\dim E_r.
\end{equation}
Taking dimensions in the quotient gives
\begin{equation}
\label{p:second-difference}
 \boxed{\ b_r=2d_r-d_{r-1}-d_{r+1}.\ }
\end{equation}
This equation has been proved before any choice of chains. In particular,
its right side is nonnegative because it is the dimension of an actual
quotient space. We next prove why that quotient counts chains of exactly
length $r$.

\subsection{Lifting the endpoint filtration to a basis of chains}
\label{p:chain-construction}

For each $r$, choose vectors
\[
 u_{r,1},\ldots,u_{r,b_r}\in B_r
\]
whose residue classes form a basis of $B_r/B_{r+1}$. Such vectors exist by
finite-dimensional basis extension, proved after the dimension theorem.
Equivalently, start from the empty basis of $B_{R+1}=0$ and extend a basis
successively as we pass to $B_R,B_{R-1},\ldots,B_1$.

It follows that, for every $h\geq1$, the vectors
\begin{equation}
\label{p:endpoint-basis}
 \{u_{r,j}:h\leq r\leq R,\ 1\leq j\leq b_r\}
 \quad\text{form a basis of } B_h.
\end{equation}
Here is the elementary reason. If a subspace $W\subseteq U$ has a basis and
we lift a basis of $U/W$, the two lists together span $U$: subtract a
combination of the lifts to leave an element of $W$. They are independent:
reduce a relation modulo $W$ first, and then use independence in $W$.
Applying this observation down the filtration proves
\eqref{p:endpoint-basis}.

Because $u_{r,j}\in\im N^{r-1}$, choose $v_{r,j}\in V$ with
\[
 N^{r-1}v_{r,j}=u_{r,j}.
\]
Since $u_{r,j}\in\Ker N$, we also have $N^rv_{r,j}=0$. For
$1\leq h\leq r$, write
\[
 w_{r,j,h}=N^{r-h}v_{r,j}.
\]
The index $h$ is the height above zero: $w_{r,j,1}=u_{r,j}$,
$Nw_{r,j,1}=0$, and $Nw_{r,j,h}=w_{r,j,h-1}$ for $h>1$.
For example, a chain of length three has the form
\[
\begin{tikzcd}[column sep=large]
 w_{3,j,3} \arrow[r,mapsto,"N"] &
 w_{3,j,2} \arrow[r,mapsto,"N"] &
 w_{3,j,1} \arrow[r,mapsto,"N"] &0.
\end{tikzcd}
\]
The nodes in this last diagram are vectors, and the arrows show the action
of $N$ on them; it is not a diagram of vector spaces and linear maps.

\begin{theorem}[A basis of nilpotent chains]
\label{p:chain-basis}
The vectors $w_{r,j,h}$, for $1\leq r\leq R$, $1\leq j\leq b_r$ and
$1\leq h\leq r$, form a basis of $V$.
Consequently there is a $k[t]$-module isomorphism
\[
 M\cong\bigoplus_{r=1}^{R}P_r^{\oplus b_r}.
\]
\end{theorem}

\begin{proof}
First prove independence. Suppose that a linear combination of these
vectors is zero. If any coefficient is nonzero, let $H$ be the largest
height $h$ at which a nonzero coefficient occurs. Apply $N^{H-1}$ to the
relation. Every term of height less than $H$ vanishes, and a term
$w_{r,j,H}$ becomes $u_{r,j}$. We obtain a relation among the vectors
$u_{r,j}$ with $r\geq H$. These vectors are independent by
\eqref{p:endpoint-basis}. Thus every coefficient at height $H$ is zero,
contradicting its choice. No nonzero relation exists.

For spanning, prove by induction on $h$ that $K_h$ lies in the span of the
chosen chain vectors. The case $h=0$ is immediate. Given $x\in K_h$,
the vector $N^{h-1}x$ lies in $B_h$. By \eqref{p:endpoint-basis}, write
\[
 N^{h-1}x=\sum_{r\geq h,\,j}c_{r,j}u_{r,j}.
\]
Each $u_{r,j}$ in this sum has the chosen lift $w_{r,j,h}$ at height $h$.
Therefore
\[
 x-\sum_{r\geq h,\,j}c_{r,j}w_{r,j,h}\in K_{h-1}.
\]
The induction hypothesis puts this difference in the required span, and
so puts $x$ there. Since $K_R=V$, spanning follows.

For the module statement, on the copy of $P_r$ indexed by $(r,j)$ define
\[
 \overline{p(t)}\longmapsto p(N)v_{r,j}.
\]
This is a well-defined module map because $N^rv_{r,j}=0$, by the same
argument as Proposition~\ref{p:hom-kernel}. Combining these maps gives a map
from the displayed finite direct sum to $M$. It sends the standard basis
vectors $1,t,\ldots,t^{r-1}$ in each summand to the corresponding chain,
in reverse height order. The basis statement makes this map bijective;
its inverse respects the module operations because the map is bijective
and respects those operations. It is therefore a module isomorphism.
\end{proof}

The choices in this proof are visible. We selected representatives of
quotient bases and then lifted them through powers of $N$. Different
choices can give different chain bases. The spaces $B_r$, the quotients
$E_r$, and their dimensions $b_r$ depend only on $N$.

In particular, the construction proves the precise interpretation
\[
 \boxed{
 \dim\frac{\Ker N\cap\im N^{r-1}}
              {\Ker N\cap\im N^r}
 =\text{the number of chains of length exactly }r.
 }
\]
The quotient removes endpoints which can be lifted farther. A basis of
what remains supplies the new chains of length $r$ in the construction.

\subsection{Reading the operator from its measurements}
\label{p:profile-reconstruction}

In one chain of length $s$, the kernel of $N^r$ is spanned by the bottom
$\min(r,s)$ vectors. The direct-sum decomposition therefore gives
\begin{equation}
\label{p:profile-formula}
 d_r=\sum_{s=1}^{R}b_s\min(r,s),\qquad
 d_r-d_{r-1}=\sum_{s\geq r}b_s.
\end{equation}
The increment counts chains of length at least $r$. Subtracting successive
increments gives exactly the second-difference formula
\eqref{p:second-difference}, which we had already obtained from canonical
subquotients.

\begin{corollary}[Reconstruction by a finite nilpotent profile]
\label{p:reconstruction}
Two finite-dimensional nilpotent $k[t]$-modules $M$ and $M'$ are isomorphic
if and only if
\[
 \dim_k\Hom_{k[t]}(P_r,M)=\dim_k\Hom_{k[t]}(P_r,M')
 \qquad\text{for every }r\geq1.
\]
If a common nilpotence bound $R$ is known, it suffices to use $1\leq r\leq R$.
\end{corollary}

\begin{proof}
A module isomorphism identifies the spaces of maps by postcomposition,
so it preserves all the displayed dimensions. Conversely, equality of
the dimensions means equality of the $d_r$ by
Proposition~\ref{p:hom-kernel}. The formula
$b_r=2d_r-d_{r-1}-d_{r+1}$ then gives the same chain multiplicities for
both modules. Theorem~\ref{p:chain-basis} identifies each with the same
direct sum of the $P_r$, so they are isomorphic.

For a known common bound $R$, put $d_0=0$ and $d_{R+1}=d_R$; nilpotence
justifies the latter equality. The first $R$ measurements therefore
determine every possibly nonzero $b_r$.
\end{proof}

Thus the dimensions of these particular spaces of maps determine the
isomorphism class in this particular category. They do not choose an
isomorphism or a preferred basis. The proof succeeds because we established
a decomposition theorem and then showed that the measurements recover
every multiplicity in that decomposition.

Here are the first distinctions. Write $J_s$ for multiplication by $t$ on
$P_s$; as an operator, it is a single nilpotent chain of length $s$.
In dimension three we obtain
\begin{center}
\begin{tabular}{@{}lll@{}}
\toprule
Operator & $(d_1,d_2,d_3)$ & $(b_1,b_2,b_3)$ \\
\midrule
Zero on $k^3$ & $(3,3,3)$ & $(3,0,0)$ \\
$J_2\oplus J_1$ & $(2,3,3)$ & $(1,1,0)$ \\
$J_3$ & $(1,2,3)$ & $(0,0,1)$ \\
\bottomrule
\end{tabular}
\end{center}
The skipped-attachment operator from Section~\ref{x:skip-layer},
with $N(e_3)=e_1$ and $N(e_1)=N(e_2)=0$, is the middle row:
$(e_3,e_1)$ is a chain of length two and $e_2$ is a chain of length one.
Its first measurement is $d_1=2$, whereas the zero operator has $d_1=3$.
Thus even maps from $P_1$ recover the distinction that both adjacent
two-layer summaries missed.

All three have three simple factors isomorphic to
$S_0=P_1=k[t]/(t)$. Indeed, in a chain the subspaces spanned by its first
$h$ bottom vectors are submodules, and each successive quotient is a copy
of $S_0$. Concatenating these filtrations across the direct sum gives a
composition series. The space $S_0$ is simple because it is
one-dimensional over $k$ and every submodule is a $k$-subspace.
Consequently all three examples have the same class $3[S_0]$ in the
Grothendieck group. Their spaces of maps from $P_1,P_2,P_3$ distinguish them.

This does not contradict the universality of the additive invariant:
these dimensions of spaces of maps need not be additive on short exact
sequences. For example,
\[
 0\longrightarrow P_1\xrightarrow{\ 1\mapsto t\ }P_2
 \xrightarrow{\ \overline p\mapsto\overline p\ }P_1
 \longrightarrow0
\]
is exact. The first map identifies $P_1$ with the submodule $kt$ of
$P_2$, and the second map is reduction modulo $t$, with kernel $kt$.
Yet $\dim\Hom_{k[t]}(P_1,P_2)=1$, whereas the two endpoint measurements
sum to $1+1=2$. A module map $P_1\to P_2$ must send $1$ into $kt$, so its
composition with $P_2\to P_1$ is zero. In particular, the identity map
of the quotient $P_1$ cannot lift to $P_2$. The failure of additivity is
the lifting obstruction from the preceding chapter, appearing in an
explicit measurement.

For a slightly larger calculation, suppose measurements give
\[
 (d_0,d_1,d_2,d_3,d_4,d_5)=(0,3,5,6,7,7)
\]
for an operator with $N^4=0$. The increments are
$3,2,1,1,0$. Their successive differences give
\[
 (b_1,b_2,b_3,b_4)=(1,1,0,1).
\]
Hence $M\cong P_1\oplus P_2\oplus P_4$. The dimensions add to
$1+2+4=7$, and the kernel dimensions of this sum, calculated as
$\min(r,1)+\min(r,2)+\min(r,4)$, recover the original measurements.

The scope matters. These same tests do not reconstruct an arbitrary
operator. If $t$ acts as the identity on a nonzero vector space, every
$N^r$ is the identity and every $\Hom_{k[t]}(P_r,M)$ is zero. Its profile
is therefore the same as that of the zero module. The nilpotence
assumption ensured that our tests eventually reach all of $M$; without
it, an entire part of the object can remain undetected.

We have now proved a concrete reconstruction result. Composition factors
identify the simple layers. The canonical subquotients $E_r$ measure how
far their endpoints can be lifted along the operator. The test modules
$P_r$ turn those lifting depths into dimensions of spaces of maps, and the
chain-basis theorem shows that this information is sufficient to recover
the finite-dimensional nilpotent object up to isomorphism.

% END INPUT 09-probes-and-reconstruction.tex


% BEGIN INPUT 09b-hall-associativity.tex
\section{Associativity by counting filtrations}\label{hall:chapter}

An extension records an attachment. We can also ask how many attachments
of one type occur inside a fixed middle object. This question leads to a
new multiplication. Its associativity will follow from the correspondence
between submodules of a quotient and submodules containing its kernel.
The quotient calculation that began the exposition will thus return as
the mechanism behind an algebra law.

\subsection{A setting where the counts are finite}

Fix a finite field $\kk$ with $q$ elements. We work with finite-dimensional
$\kk[t]$-modules on which $t$ acts nilpotently. Equivalently, our objects
are pairs $(V,N)$ with $N^s=0$ for some $s$, and our maps are linear maps
commuting with the operators. A submodule is exactly an $N$-invariant
subspace. Throughout this chapter an isomorphism must respect $N$.

Submodules and quotients stay in this class because the same power of
$N$ is zero on them. Extensions stay in the class too. Indeed, in an
exact sequence
\[
0\longrightarrow B\longrightarrow E\longrightarrow A\longrightarrow0,
\]
if $t^aA=0$ and $t^bB=0$, then $t^aE\subseteq B$ and hence
$t^{a+b}E=0$. Also $\dim E=\dim A+\dim B$, by the proved dimension
formula for an exact sequence. These are closure checks, not extra
assumptions about the category.

For isomorphism classes $A,B,E$, define the nonnegative integer
\begin{equation}\label{hall:count}
g_{A,B}^{E}
=\#\{U\subseteq E:\ U\text{ is a submodule},\ U\cong B,\ E/U\cong A\}.
\end{equation}
The order matters: $B$ is the lower layer and $A$ is the quotient above it.
This counts actual submodules of $E$. It does not count chosen
isomorphisms $B\to U$, chosen quotient identifications, or marked
extension classes. An isomorphism $E\to E'$ gives a bijection on these
sets by sending $U$ to its image, so the count depends only on the three
isomorphism classes.

The set in \eqref{hall:count} is finite because $E$ has finitely many
elements. For fixed $A,B$, only finitely many classes $E$ can occur.
They have the fixed dimension $d=\dim A+\dim B$. Choosing a basis
identifies such an operator with a $d$ by $d$ matrix over $\kk$, and
there are $q^{d^2}$ matrices altogether. Each isomorphism class has a
representative among them. This argument establishes finiteness without
requiring a classification of nilpotent operators.

Let $\mathcal H$ be the rational vector space with one basis symbol
$u_M$ for each isomorphism class in our collection. Explicitly, its
vectors are finite rational linear combinations of these symbols.
Define a bilinear multiplication by
\begin{equation}\label{hall:product}
u_Au_B=\sum_{[E]}g_{A,B}^{E}u_E.
\end{equation}
The preceding finiteness checks make this an element of $\mathcal H$.
It is the untwisted Hall multiplication for this collection of modules.
We now prove its algebra laws.

\subsection{One flag, two intermediate objects}

\begin{theorem}[Associativity from a flag bijection]\label{hall:associativity}
The multiplication \eqref{hall:product} is associative, with identity $u_0$.
\end{theorem}

\begin{proof}
Fix $A,B,C,E$. Count the set of flags
\begin{equation}\label{hall:flag}
0\subseteq U\subseteq V\subseteq E,
\qquad U\cong C,\quad V/U\cong B,\quad E/V\cong A.
\end{equation}
Inclusions may be equal when one of the specified factors is zero.

First fix $U$. If $E/U\cong D$, there are $g_{D,C}^{E}$ choices of $U$.
For each such $U$, the possible submodules $V/U\subseteq E/U$ with
the required factors are counted by $g_{A,B}^{D}$. The quotient
correspondence gives a bijection
\[
\{V:U\subseteq V\subseteq E\}
\longleftrightarrow
\{W:W\subseteq E/U\},
\quad V\longmapsto V/U,\quad W\longmapsto\pi^{-1}(W).
\]
It also gives $(E/U)/(V/U)\cong E/V$, by the map
$(e+U)+(V/U)\mapsto e+V$. Thus the number of flags is
\[
\sum_{[D]}g_{A,B}^{D}g_{D,C}^{E}.
\]
This is the coefficient of $u_E$ in $(u_Au_B)u_C$.

Alternatively fix $V$. If $V\cong D$, there are $g_{A,D}^{E}$ choices
of $V$, and then $g_{B,C}^{D}$ choices of $U\subseteq V$.
The same flags are therefore counted by
\[
\sum_{[D]}g_{B,C}^{D}g_{A,D}^{E},
\]
the coefficient of $u_E$ in $u_A(u_Bu_C)$. No choice of an isomorphism
to a representative $D$ affects either count: every such isomorphism
induces a bijection of the relevant submodule sets. Both counting
procedures identify their data with precisely the flags in
\eqref{hall:flag}. Equality of all coefficients proves associativity on
basis vectors, and bilinearity proves it on $\mathcal H$.

For the unit, the only submodule of $E$ isomorphic to zero is $0$,
and the only submodule with zero quotient is $E$. Consequently
$g_{A,0}^{E}=g_{0,A}^{E}$ is $1$ if $E\cong A$ and $0$ otherwise.
Hence $u_Au_0=u_0u_A=u_A$.
\end{proof}

The two readings of the flag use different intermediate objects:
\[
\begin{tikzcd}[column sep=large,row sep=large]
0\arrow[r]&U\arrow[r,hook]\arrow[d,equal]&V\arrow[r,two heads]\arrow[d,hook]
  &V/U\arrow[r]\arrow[d,hook]&0\\
0\arrow[r]&U\arrow[r,hook]&E\arrow[r,two heads]&E/U\arrow[r]&0.
\end{tikzcd}
\]
Both rows are exact. The right vertical map sends $v+U$ to $v+U$;
it is injective since its zero class means $v\in U$. The square on the
right commutes on each representative $v$. The quotient of that vertical
inclusion is $E/V$. Thus the intermediate object is either $V$, carrying
the lower two layers, or $E/U$, carrying the upper two layers.

At the beginning of the exposition, associativity of the group operation
was an axiom permitting us to compose translations. Here associativity
of a newly defined operation is a theorem. The theorem is proved by a
bijection: the two bracketings organize exactly the same data. This is
a structural explanation with an explicit inverse correspondence, not
an appeal to a resemblance between diagrams.

\subsection{The first product, computed completely}

Write $S=\kk[t]/(t)$ and $J_r=\kk[t]/(t^r)$. Thus $J_r$ is the single
nilpotent chain of length $r$ constructed in the preceding chapter.
There are two isomorphism classes of two-dimensional nilpotent modules:
$S\oplus S$ and $J_2$. One can see this directly. If $N=0$, the first
description holds. If $N\neq0$, a two-dimensional nilpotent operator has
a chain $Nv\neq0$, $N^2v=0$, and $(Nv,v)$ is a basis. For completeness,
one may obtain such a chain by taking the last two nonzero vectors in
any maximal chain; independence of a chain of length three would
contradict dimension two. The chain basis gives $J_2$.

For $S\oplus S$, every line is a submodule of type $S$ with quotient
$S$. There are $q+1$ lines: the $q$ lines spanned by $(1,a)$ for
$a\in\kk$, and the line spanned by $(0,1)$. In $J_2$, every invariant
line has zero restricted operator, since a nilpotent scalar is zero in
a field. It must therefore lie in $\Ker N$, which is itself one line.
There is exactly one choice. Hence
\begin{equation}\label{hall:square}
u_S^2=(q+1)u_{S\oplus S}+u_{J_2}.
\end{equation}

The nonsplit middle object appears once in this product, although the
marked extension calculation gave $q-1$ nonzero classes when the
endpoints are identified with $S$. The reason is the definition of the
count. Equation~\eqref{hall:square} counts submodules of a fixed middle
object. Marked extension equivalence keeps endpoint identifications;
our Hall basis keeps only the middle object's isomorphism class.
Neither convention can silently replace the other.

\subsection{Three layers and a second count}

The preceding reconstruction theorem gives exactly three nilpotent
types in dimension three:
\[
S^{\oplus3},\qquad J_2\oplus S,\qquad J_3.
\]
The coefficient of $u_E$ in $u_S^3$ counts full invariant flags
$0\subset U\subset V\subset E$, with $\dim U=1$ and $\dim V=2$.
Every factor is $S$, because a one-dimensional nilpotent operator is zero.

For $E=S^{\oplus3}$ all flags are invariant. There are
$(q^3-1)/(q-1)=q^2+q+1$ choices of $U$: the nonzero vectors partition
into sets of $q-1$ vectors spanning the same line. For each $U$, the
planes containing it correspond to the $q+1$ lines in $E/U$.
This gives $(q^2+q+1)(q+1)$ flags.

For $E=J_2\oplus S$, put $L=\im N$ and $K=\Ker N$. Then
$L\subset K$, $\dim L=1$, and $\dim K=2$. An invariant line $U$ lies
in $K$, giving $q+1$ choices. An invariant plane $V$ contains $L$:
the induced operator on the one-dimensional quotient $E/V$ is
nilpotent and hence zero, so $N(E)\subseteq V$. Conversely, every plane
containing $L$ is invariant. If $U=L$, there are $q+1$ possible planes
through it. For each of the other $q$ lines $U\subset K$, the only
plane containing both $U$ and $L$ is $K$. The count is $2q+1$.

For $J_3$, an invariant line must equal $\Ker N$, and an invariant
plane must equal $\im N$ by the same quotient argument and their
dimensions. These two subspaces form a flag, so the count is $1$.
We obtain
\begin{equation}\label{hall:cube}
u_S^3=(q+1)(q^2+q+1)u_{S^{\oplus3}}
       +(2q+1)u_{J_2\oplus S}+u_{J_3}.
\end{equation}
Associativity says that either order of the two successive Hall
multiplications gives this count. The flag description explains why
before any structure constants are expanded.

\subsection{What the multiplication remembers}

Every nonzero coefficient in $u_Au_B$ comes from a short exact sequence,
so it satisfies
\[
[E]=[A]+[B]\quad\text{in }K_0.
\]
This makes Hall multiplication graded by the additive invariant: the
degree of $u_M$ is $[M]$. In our nilpotent collection the only simple
module is $S$. Indeed, if a nonzero nilpotent module is simple, its
nonzero kernel of $t$ is the whole module; then every vector line is a
submodule, forcing dimension one. Writing $\mathcal N_q$ for this
nilpotent collection, we have $K_0(\mathcal N_q)\cong\mathbb Z$, and
this grading is ordinary dimension. Here $K_0(\mathcal N_q)$ is formed
using generators and exact-sequence relations inside $\mathcal N_q$,
as in Section~\ref{k:section}; the closure checks at the beginning of
this chapter make that construction applicable.

Within a single degree, the Hall algebra retains different basis vectors
for different isomorphism classes. The terms in \eqref{hall:square} all
have degree two, but the multiplication distinguishes the split object
from the chain and records their different numbers of subobjects.
This statement is a grading statement; it does not claim that $K_0$ is
a quotient algebra of $\mathcal H$.

We have reached three distinct forms of information. Composition
multiplicities retain what is additive across exact sequences.
Marked extensions retain a specified two-layer attachment. Hall
coefficients retain counts of subobjects and their quotients, and their
associativity packages the compatibility of longer flags. In the
nilpotent setting, the spaces of maps studied earlier recover the whole
isomorphism class. Each construction answers a precise question about
what the preceding description left unresolved.

Hall algebras can be developed much further, including settings where
geometric invariants replace finite counts. That further theory requires
additional constructions. The result established here is the finite
counting algebra and its associativity, proved directly from flags.

% END INPUT 09b-hall-associativity.tex

\clearpage\part{Coherent quotients and deeper structure}

% BEGIN INPUT 10-higher-quotients.tex
\section{Higher quotients: keeping the identifications}
\label{h:chapter}

The exact Mal'cev argument has answered a specific question: which simple
quotient objects survive every choice of finite normal series?  We now ask
what was omitted when we retained only those quotient objects.  A quotient
can remember not only which points have been identified, but also the ways
of identifying them and the compatibility between successive identifications.
This is where the higher framework becomes useful.

The change of foundations in this section is explicit.  We use spaces,
homotopy limits, and a small number of stated theorems about
$\infty$-categories.  Their foundations are not proved here.  A result marked
\emph{Framework theorem, quoted} is an external input; the propositions
following it are deductions with their own proofs.  None of these inputs
is used in the earlier proof of Jordan--H\"older.  Conversely, invoking the
word ``higher'' will not be used to justify an unproved quotient law.

\subsection{One construction of homotopical algebraic objects}

Write $\mathcal S$ for the $\infty$-category of spaces, or homotopy types.
A set is a discrete space.  A space can also have paths, homotopies between
paths, and further coherences.  The mapping object
$\operatorname{Map}(X,Y)$ is consequently a space: its points are maps,
its paths are homotopies, and so on.  An equivalence is the appropriate
replacement for an isomorphism.  Limits and colimits in this setting mean
homotopy-coherent limits and colimits.

An ordinary one-sorted algebraic theory $\mathbb T$ is a small category
with finite products, generated under products by an object $t$.
The objects $t^n$ stand for $n$ inputs.  An arrow $t^n\to t$ stands for
an $n$-ary operation, and composition records substitution.  Equations
between arrows record equations between operations.  A model in a category
with finite products is a finite-product-preserving functor out of
$\mathbb T$.

For groups, multiplication is an arrow $t^2\to t$, inverse is $t\to t$,
and identity is $t^0\to t$.  Associativity is equality of the two arrows
$t^3\to t$ obtained by substitution.  For left $R$-modules we include
addition, zero, inverse, and one unary operation for each $r\in R$.
The equation for multiplication of scalars is an equation between
composites of these unary operations.  Thus the presentation does not
discard the ring when it passes to a common language.

Define
\begin{equation}\label{h:theory-models}
 \mathcal C_{\mathbb T}
   =\operatorname{Fun}^{\times}(\mathbb T,\mathcal S).
\end{equation}
Preservation of products now means preservation up to the specified coherent
equivalences.  This is the promised single construction: change the target
of the theory from sets to spaces, retaining the operations and their
compatibilities.

\begin{proposition}[Recovering the ordinary algebra]\label{h:discrete}
The full subcategory of $0$-truncated objects of $\mathcal C_{\mathbb T}$
is equivalent to the nerve of the ordinary category
$\operatorname{Fun}^{\times}(\mathbb T,\mathbf{Set})$.
\end{proposition}
\begin{proof}
An object $M$ is $0$-truncated when $\operatorname{Map}(P,M)$ is discrete
for every $P$.  If all $M(t^n)$ are discrete, the space of natural
transformations into $M$ is a limit of discrete mapping spaces and is
therefore discrete.  Thus such $M$ is $0$-truncated.

Conversely, the covariant representable
$F=\operatorname{Hom}_{\mathbb T}(t,-)$ preserves products, so belongs to
$\mathcal C_{\mathbb T}$.  Yoneda identifies
$\operatorname{Map}(F,M)$ with $M(t)$: a natural transformation is
determined by its value on $\id_t$, and the value on an arrow
$a:t\to t^n$ is then $M(a)$ applied to that value. This is the
space-valued Yoneda lemma, including its coherent mapping-space statement.
If $M$ is $0$-truncated, $M(t)$
is discrete; product preservation gives $M(t^n)\simeq M(t)^n$, so all its
values are discrete.  Functors between ordinary discrete values have
ordinary natural transformations and the ordinary product-preservation
condition.  This identifies both objects and maps of the displayed
subcategory.
\end{proof}

\begin{theorem}[Framework theorem, quoted: rigidification]
\label{h:rigidification}
For a small ordinary finite-product theory $\mathbb T$, the homotopy theory
of strict simplicial $\mathbb T$-models, with objectwise weak equivalences,
presents $\mathcal C_{\mathbb T}$.
\end{theorem}
This is Lurie's Corollary 5.5.9.3, using Proposition 5.5.9.2, in
\cite{HTT}; see also \cite{Badzioch}.  This theorem is what permits us
to compute with familiar simplicial algebra while interpreting the answer
in \eqref{h:theory-models}.  It is not a claim that all coherences can
simply be omitted.

For group theory, the result is the homotopy theory of simplicial groups.
Two further foundational inputs describe it. The Kan loop-group and
classifying-space adjunction identifies this homotopy theory with pointed
connected spaces \cite[Chapter V, Proposition 6.3 and Corollary 6.4]{GoerssJardine}.
The loop--delooping recognition theorem identifies pointed connected spaces
with grouplike $E_1$-spaces \cite[Theorem 5.2.6.10]{HA}.
An $E_1$ multiplication is associative with its coherent homotopies;
``grouplike'' says that its components form a group. The two quoted
equivalences together supply the asserted identification; rigidification
alone does not supply either comparison.

For module theory, the Dold--Kan theorem identifies simplicial $R$-modules
with nonnegative chain complexes of $R$-modules; weak equivalences become
quasi-isomorphisms \cite[Tag 019G]{Stacks}.  Localizing therefore gives
connective derived $R$-modules.  For a commutative ring $R$, this agrees
with connective modules over its Eilenberg--Mac Lane ring spectrum $HR$
\cite[Theorem 7.1.2.13]{HA}.  The $HR$-module structure is essential:
an arbitrary spectrum, even a connective one, is not thereby an
$R$-module.

The earlier theorem applies to the ordinary models recovered by
Proposition~\ref{h:discrete}, whenever their category satisfies its pointed
exact Mal'cev hypotheses.  Groups and modules are therefore instances
inside a common homotopical construction.  The exactness hypotheses do
real work; product-preserving models in spaces alone do not imply a
composition theorem for every higher object.

\subsection{The action groupoid and the role of associativity}

Let an ordinary group $G$ act on the right on a set $X$.  Form a groupoid
whose objects are elements $x\in X$, and whose arrows are
\[
 (x,g):x\longrightarrow xg.
\]
Composition is
\[
 (x,g)\text{ followed by }(xg,h)=(x,gh).
\]
It is defined because $(xg)h=x(gh)$; it is associative because multiplication
in $G$ is associative; identities and inverses follow from the corresponding
group laws.  This construction retains the actual witness $g$ of an
identification, not only the fact that two points lie in the same orbit.

Its nerve is the simplicial set
\begin{equation}\label{h:action-nerve}
 X_n=X\times G^n.
\end{equation}
In degree $n$, the tuple $(x,g_1,\ldots,g_n)$ means the composable chain
\[
 x\longrightarrow xg_1\longrightarrow xg_1g_2
   \longrightarrow\cdots\longrightarrow xg_1\cdots g_n.
\]
The first face removes $x$ by replacing it with $xg_1$; an interior face
replaces $g_i,g_{i+1}$ by their product; the last face removes $g_n$.
A degeneracy inserts an identity element.  For three consecutive arrows,
the equality of the two ways to take interior faces is precisely
$(g_1g_2)g_3=g_1(g_2g_3)$.  The unit identities similarly supply the
face--degeneracy identities.  The nerve is thus a way to organize, in
one object, the algebraic compatibility behind the elementary maneuvers.

\begin{definition}\label{h:homotopy-quotient}
The homotopy quotient $X/\!/G$ is the space obtained by realizing the
nerve \eqref{h:action-nerve}.  For a point with the trivial action, write
$BG=*/\!/G$.
\end{definition}

\begin{proposition}[What an action quotient retains]\label{h:stabilizers}
The connected components of $X/\!/G$ are the orbits of $X$.  The component
of $x$ is equivalent to $BG_x$, where
$G_x=\{g\in G:xg=x\}$.  In particular a free action has homotopy quotient
equivalent to its discrete orbit set.
\end{proposition}
\begin{proof}
There is an arrow from $x$ to $y$ exactly when $y=xg$ for some $g$.
Thus components are orbits.  In a fixed orbit, the full subgroupoid on
one representative $x$ has automorphism group $G_x$.  Its inclusion is
fully faithful by construction and essentially surjective because every
object in the orbit is isomorphic to $x$.  It is therefore an equivalence
of groupoids.  Realization preserves such equivalences, giving $BG_x$.
For a free action all stabilizers are trivial, so each component is
contractible.
\end{proof}

For the trivial action on a point, the ordinary orbit set is a point,
whereas the homotopy quotient is $BG$ and retains every automorphism in
$G$.  For a subgroup $N\leq G$ acting on $G$ by right translation, the
action is free, so $G/\!/N\simeq G/N$ as spaces.  If $N$ is normal,
multiplication also descends to these cosets.  The two assertions are
different: freeness determines the homotopy type of the quotient;
normality determines its induced group multiplication.

\subsection{Extensions become fiber diagrams}

An ordinary pullback asks for equal images. A homotopy pullback
$X\times^h_ZY$ retains a path between the images: a point is a pair
$(x,y)$ together with a path between $f(x)$ and $g(y)$.
Paths between such data must respect the specified comparison, with the
corresponding higher coherences. In particular, a homotopy fiber over
$z$ remembers both a point $x$ and a path between $f(x)$ and $z$.
The next calculation is a case where this extra datum is a group element.

\begin{proposition}\label{h:extension-fiber}
A short exact sequence of ordinary groups
\[
 1\longrightarrow N\longrightarrow G\xrightarrow{p}Q\longrightarrow1
\]
determines a homotopy pullback square
\[
\begin{tikzcd}
 BN \arrow[r] \arrow[d] & BG \arrow[d,"Bp"]\\
 * \arrow[r] & BQ.
\end{tikzcd}
\]
\end{proposition}
\begin{proof}
The homotopy fiber of the functor between one-object groupoids induced by
$p$ can be modeled by the groupoid with objects $q\in Q$ and arrows
$g:q\to q p(g)$.  Indeed its objects are an object of $BG$ together with
an isomorphism from the chosen object of $BQ$ to its image, and its arrows
are the compatible isomorphisms.  This is the groupoid model of the
homotopy pullback, retaining the isomorphism in the pullback datum.
Surjectivity of $p$ makes it transitive, and the automorphisms of $1\in Q$
are exactly $\ker p=N$.  The argument of
Proposition~\ref{h:stabilizers} identifies this fiber with $BN$.
The inclusion of its representative object gives the usual map $BN\to BG$.
\end{proof}

This square is a fiber description, not a claim about an ordinary
set-theoretic fiber of the map of realizations.  It also is not automatically
a cofiber square.  The distinctions matter because the homotopy theory
of spaces is unstable.

\subsection{Quotients that survive change of base}

An $\infty$-topos is a setting for spaces varying over a site; formally it
is an accessible left exact localization of a presheaf $\infty$-category
of spaces.  The theory-model category $\mathcal C_{\mathbb T}$ need not
be an $\infty$-topos.  The effective-quotient results below apply in spaces
or another specified $\infty$-topos; the subsequent descent equivalence
is stated for spaces.  We need the following two consequences of the
foundational theory of an $\infty$-topos.

\begin{theorem}[Framework theorem, quoted: effective quotients]
\label{h:topos-foundation}
In an $\infty$-topos, colimits are preserved by pullback.  Every groupoid
object is the \v{C}ech nerve of the map from its object of objects to its
realization.
\end{theorem}
These are the universality and effectivity assertions in
\cite[Theorem 6.1.0.6]{HTT}.  A \v{C}ech nerve of $U\to Y$ has degree
$n$ term $U\times_Y\cdots\times_Y U$ with $n+1$ factors.  Effectivity
says that the identification data in a groupoid object are precisely the
data recovered by pulling back its quotient map against itself.

\begin{proposition}[A quotient can be computed after base change]
\label{h:quotient-basechange}
In an $\infty$-topos, let $U_\bullet$ be a groupoid object with realization $Y$, and let
$f:Y'\to Y$.  Then
\[
 \left|Y'\times_Y U_\bullet\right|\simeq Y'.
\]
Moreover $Y'\times_Y U_\bullet$ is the \v{C}ech nerve of
$Y'\times_Y U_0\to Y'$.
\end{proposition}
\begin{proof}
Pullback commutes with the realization colimit by
Theorem~\ref{h:topos-foundation}, so its realization is
$Y'\times_Y|U_\bullet|\simeq Y'\times_Y Y\simeq Y'$.
Effectivity identifies $U_n$ with the $(n+1)$-fold fiber product of $U_0$
over $Y$.  Pulling back this description gives the corresponding
$(n+1)$-fold fiber product of $Y'\times_Y U_0$ over $Y'$, including its
face and degeneracy maps.  These identifications prove the nerve assertion.
\end{proof}

The usefulness is concrete.  A quotient description can be transported to
a new parameter space without recomputing its coherence relations.  If we
make two successive base changes, the two answers agree by the canonical
associativity equivalence of fiber products.  That equivalence is itself
compatible with a third base change: both descriptions are limits of the
same diagram.  Universal properties account for the compatibility that
would otherwise have to be checked separately in each example.

\subsection{Descent: a family over \texorpdfstring{$BG$}{BG} is an action}

\begin{theorem}[Framework theorem, quoted: straightening for spaces]
\label{h:straightening}
For a space $B$, taking homotopy fibers with transport gives an equivalence
\[
 \mathcal S_{/B}\simeq\operatorname{Fun}(B,\mathcal S).
\]
\end{theorem}
This is the groupoid case of straightening and unstraightening
\cite[\S\S2.2,3.2]{HTT}.  The symbol $B$ on the right is regarded as
an $\infty$-groupoid.  Transport along an invertible path is an equivalence;
transport along composable paths includes coherent composition data.

For $B=BG$, a functor $BG\to\mathcal S$ is a space with coherent
$G$-action.  A map $E\to BG$ determines the fiber $X$ over the basepoint,
and loops in $BG$ transport $X$ around the base.  Conversely, the
Grothendieck construction assembles the action into its total space
$X/\!/G\to BG$.  These are the actual forward and inverse constructions
in this instance of the theorem.  In identifying this functor notation
with the usual left-action convention, we use the canonically equivalent
opposite one-object model via inversion: the earlier right action becomes
$g\cdot x=xg^{-1}$.  Our concrete groupoid calculations continue to use
the right-multiplication convention, with first $g$, then $h$, composed
as $gh$.  No change of coefficient structure is involved.

\begin{corollary}[Equivariant information is descent information]
\label{h:action-descent}
The equivalence
\[
 \mathcal S_{/BG}\simeq\operatorname{Fun}(BG,\mathcal S)
\]
identifies maps over $BG$ with coherent equivariant maps.  Its restriction
along $BH\to BG$ for a subgroup $H$ identifies homotopy pullback of families
with restriction of their actions.  Sections of $X/\!/G\to BG$ form the
space of coherent fixed points of $X$.
\end{corollary}
\begin{proof}
The first assertion is the assertion about mapping spaces in the
equivalence.  Taking the fiber of a pulled-back family at a point of $BH$
gives the same fiber as before, and its path transport is the original
transport along the image path in $BG$.  This is precisely restriction
along the functor $BH\to BG$.  Finally a section is a map over $BG$ from
the identity family $BG\to BG$.  That family corresponds to the constant
one-point functor.  Maps from it to the action functor are compatible
choices of a point, with the homotopies and coherences making it fixed.
\end{proof}

For a discrete $G$-set these are ordinary fixed points.  For a general
space the coherences are part of the answer.  We have gained a theory of
families, transport, and descent, not merely another notation for a set
of orbits.

\subsection{Intersections acquire additional components}

Here is a useful calculation that prevents a misleading translation of
the earlier subgroup diagrams.

\begin{proposition}[Double cosets in a homotopy pullback]
\label{h:double-coset}
For subgroups $A,B\leq G$, there is an equivalence of spaces
\[
 BA\times^h_{BG}BB
 \simeq
 \coprod_{[g]\in A\backslash G/B}
 B(A\cap gBg^{-1}).
\]
The component selected by the identity element is $B(A\cap B)$.
\end{proposition}
\begin{proof}
A groupoid model for the homotopy pullback has one object for each
$g\in G$.  An arrow from $g$ to $g'$ is a pair $(a,b)\in A\times B$
such that $g'=a^{-1}gb$.  First $(a,b)$ and then $(a',b')$ compose to
$(aa',bb')$, since $a'^{-1}(a^{-1}gb)b'=(aa')^{-1}g(bb')$.
This is the weak pullback: $g$ records the comparison isomorphism in
$BG$ between the images of the two chosen objects.

Its orbits are the double cosets $AgB$.  The stabilizer of $g$ consists
of pairs with $a=gbg^{-1}$ and is identified with $A\cap gBg^{-1}$ by
projection to $a$.  Apply the transitive-groupoid argument of
Proposition~\ref{h:stabilizers} on each orbit.  For $g=1$ the stabilizer
is $A\cap B$, as claimed.
\end{proof}

The classical intersection is present, but it is one component of a richer
comparison object.  Discarding the other components before doing base
change would discard genuine information.  This is one reason to keep
pointings and transport explicit when moving a subgroup argument into
homotopy theory.

In particular, a proper subgroup inclusion $H\hookrightarrow G$ does not
give a monomorphism $BH\to BG$ in $\mathcal S$.  Its homotopy fiber is the
coset space, which has more than one point.  A monomorphism of spaces has
every fiber either empty or contractible, so this map is not one.
The subgroup inclusion is monic in the category of discrete groups;
its delooping has a different categorical type.

\subsection{Higher normality requires quotient data}

For a higher group map $H\to G$, a normality datum can be specified by a
pointed connected space $Q$ and a homotopy fiber square in the ambient
$\infty$-category of spaces
\[
\begin{tikzcd}
 BH \arrow[r] \arrow[d] & BG \arrow[d]\\
 * \arrow[r] & Q.
\end{tikzcd}
\]
All displayed maps preserve the specified points. Requiring a pullback
in spaces retains the entire homotopy fiber; taking a limit only inside
the subcategory of pointed connected spaces could discard other components.
The loop space $\Omega Q$ then plays the role of the quotient higher
group.  For an ordinary normal inclusion, $Q=B(G/H)$ supplies this datum.
For general higher maps, the datum is additional structure and need not
be unique.

The appropriate isomorphism laws consequently include this structure.
Beardsley and Fox develop normality data for group objects in
$\infty$-topoi.  Their current formulation proves the first and second
Noether theorems, exhibits failures of an unrestricted third theorem,
and proves an iterated-quotient theorem with compatible data
\cite[\S5, especially Theorem 5.14]{BeardsleyFox}.  We import no general
higher Jordan--H\"older theorem from that work.

One can already see why the ordinary definition of a filtration cannot
be transferred without a decision about admissible arrows.

\begin{proposition}[Literal subobject chains force discreteness]
\label{h:monos-discrete}
Suppose a group object in spaces admits a finite chain of categorical
monomorphisms beginning at the trivial group and ending at $G$.  Then
the underlying space of $G$ is discrete up to equivalence.
\end{proposition}
\begin{proof}
Limits of product-preserving theory models are computed pointwise,
because limits commute with products; evaluation on the generator
detects equivalences because its powers give every object of the theory.
Thus limits of group objects are computed on their underlying spaces,
and equivalences are detected there.  A map is monic precisely when its
diagonal is an equivalence.  Hence the forgetful functor detects these
monomorphisms.  Their composite gives a monomorphism $*\to G$ of spaces.
Its fiber over the identity is $\Omega G$, so that loop space is
contractible.  The identity component of $G$ is therefore contractible.
Translation by a point of a group object identifies its component with
the identity component.  All components are contractible, which is
exactly the assertion that $G$ is $0$-truncated.
\end{proof}

Thus a theory of genuinely higher composition objects must specify more
than the old finite chain of monomorphisms.  Our common ordinary theorem
is complete in its stated exact framework.  The higher enlargement gives
coherent quotients, families, and descent; any additional composition
theorem must state and verify the admissible quotient and refinement
conditions it uses.  This is a precise boundary of the proved result,
not a reason to discard the enlargement.

% END INPUT 10-higher-quotients.tex


% BEGIN INPUT 11-linearization-and-derived-structure.tex
\section{Linear objects, extension data, and derived reconstruction}
\label{t:chapter}

The preceding section enlarged a quotient by retaining identifications and
their coherences.  There is a second enlargement with a different purpose:
keep a nonlinear object, and construct a linear category of coefficients
over it.  Modules then emerge from the construction rather than being
imposed on every object at the outset.  We will see this first for ordinary
groups, then for spaces over $BG$, and finally for rings and their
cotangent complexes.

This distinction also clarifies what Jordan--H\"older does not recover.
Two objects can have the same simple factors and different extension
data.  A derived mapping space can retain the missing extension classes
and their automorphisms.  A derived tensor product can retain a relation
that an ordinary tensor product discards.  These are additional structural
theorems with explicit computations, rather than a claim that a change of
language makes the factors determine the original object.

\subsection{Why a full additive replacement cannot work}

\begin{proposition}\label{t:no-full-additive}
There is no fully faithful functor from the category of groups to an
additive category.
\end{proposition}
\begin{proof}
Suppose $F:\mathbf{Grp}\to\mathcal A$ were such a functor.  The trivial
group has just one endomorphism.  Thus $F1$ has one endomorphism, so its
identity equals its zero endomorphism.  It follows that $F1$ is a zero
object: every map into or out of it is zero.  Consequently $F$ carries
zero homomorphisms to zero morphisms, since these factor through $1$.

Let $p:S_3\to C_2$ be the sign quotient.  There are four homomorphisms
$C_2\to S_3$: the trivial one and the three choices of a transposition
as image of the generator.  There are two homomorphisms $C_2\to C_2$.
Postcomposition with $p$ sends only the trivial map to zero, and sends
the three remaining maps to the identity.  In $\mathcal A$, postcomposition
with $Fp$ is a homomorphism of abelian groups
\[
 \operatorname{Hom}_{\mathcal A}(FC_2,FS_3)
 \longrightarrow \operatorname{Hom}_{\mathcal A}(FC_2,FC_2).
\]
Full faithfulness makes its kernel trivial, so it must be injective.  But
it has the four-element and two-element sets just described as its domain
and codomain.  This is a contradiction.
\end{proof}

The obstruction concerns a \emph{full} replacement preserving all maps.
It does not prohibit a linear invariant, nor an encoding with additional
nonlinear structure.  It explains why we retain the original base when
constructing its linear theory.  The homotopy category of a stable
$\infty$-category is additive, so a fully faithful embedding of the nerve
of all groups into a stable $\infty$-category would give the same
contradiction after passage to homotopy categories.

\subsection{Modules already arise as abelian objects over a group}

Fix a group $G$.  An object of $\mathbf{Grp}/G$ is a homomorphism
$p:E\to G$.  An abelian group object in this slice has, in addition,
a zero section $s:G\to E$, fiberwise addition
\[
 \mu:E\times_G E\longrightarrow E,
\]
and fiberwise inverse, satisfying the abelian group axioms.  All these
maps must themselves be group homomorphisms over $G$.  The requirement
relates the old group multiplication to the new fiberwise addition.

\begin{theorem}[Modules over a group as linear objects]
\label{t:beck}
There is an equivalence of ordinary categories
\[
 \operatorname{Ab}(\mathbf{Grp}/G)
 \simeq \mathbb Z[G]\text{-}\mathbf{Mod}.
\]
\end{theorem}
\begin{proof}
Start with an abelian group object $p:E\to G$ and let $M=\ker p$.
The fiberwise addition restricts to a second operation on $M$ with the
same identity as the inherited group multiplication, since $s$ is a
homomorphism and $s(1)=1$.  The fact that $\mu$ is a homomorphism gives
the interchange law
\[
 \mu(aa',bb')=\mu(a,b)\mu(a',b')\qquad(a,a',b,b'\in M).
\]
It follows that
\[
 \mu(a,b)=\mu(a,1)\mu(1,b)=ab,
 \qquad
 \mu(a,b)=\mu(1,b)\mu(a,1)=ba.
\]
Thus $M$ is abelian, and its two operations agree.

The section acts on $M$ by conjugation:
\[
 g\cdot m=s(g)ms(g)^{-1}.
\]
Normality of the kernel keeps this in $M$.  Conjugation is an automorphism
of its abelian group, and $s(gh)=s(g)s(h)$ gives
$(gh)\cdot m=g\cdot(h\cdot m)$.  Extending the action additively gives
a left $\mathbb Z[G]$-module.  The map
\[
 M\rtimes G\longrightarrow E,\qquad (m,g)\longmapsto m s(g)
\]
is an isomorphism: its inverse sends $e$ to
$(e s(p(e))^{-1},p(e))$.  The multiplication formula is
\begin{equation}\label{t:semidirect-law}
 (m,g)(n,h)=(m+g\cdot n,gh).
\end{equation}
Because $\mu$ is a homomorphism, its values on
$(m s(g),n s(g))$ are obtained by multiplying its value on $(m,n)$ with
its value on $(s(g),s(g))$.  The former is $m+n$ and the latter is $s(g)$,
the zero in the fiber.  Thus the fiberwise addition becomes
$((m,g),(n,g))\mapsto(m+n,g)$.

Conversely, let $M$ be a left $\mathbb Z[G]$-module.  Formula
\eqref{t:semidirect-law} is associative because the action is by
automorphisms and respects multiplication in $G$.  Its identity is $(0,1)$
and the inverse of $(m,g)$ is $(-g^{-1}\cdot m,g^{-1})$.
Projection to $G$ has the zero section $g\mapsto(0,g)$.  Fiberwise addition
is the formula above and fiberwise inverse is $(m,g)\mapsto(-m,g)$.
To check addition is a group homomorphism, compare
\[
 (m+g\cdot m')+(n+g\cdot n')
   =(m+n)+g\cdot(m'+n').
\]
This uses commutativity in $M$ and additivity of the action.  Fiberwise
inverse is a homomorphism by the same additivity, and all abelian group
axioms hold fiber by fiber.

A morphism of abelian group objects preserves $p$, $s$, and $\mu$.
Its restriction to $M$ is additive and commutes with conjugation by
$s(g)$, hence is $\mathbb Z[G]$-linear.  Conversely a module map
$u:M\to M'$ gives $(m,g)\mapsto(u(m),g)$, which preserves every structure
map.  These assignments are inverse on objects up to the displayed
isomorphism and inverse on morphisms.  This proves the equivalence.
\end{proof}

No representation has been chosen externally.  The action appeared
because a coefficient object over $G$ must be transported along $G$.
The theorem derives modules as internal linear objects over a nonlinear
base.  It does not identify a simple $\mathbb Z[G]$-module with a simple
quotient group of $G$; these are different questions about the same base.

\subsection{Stabilization retains the base and enriches the coefficients}

A spectrum object can be described by pointed objects $E_0,E_1,\ldots$
and equivalences $E_n\simeq\Omega E_{n+1}$.  The loop object is formed by
a homotopy pullback at the distinguished point.  Stabilization
$\operatorname{Sp}(\mathcal C)$ collects such objects in a suitable
$\infty$-category $\mathcal C$ with finite limits.  We use the standard
theory in \cite[\S\S1.4.2--1.4.3]{HA}, including the standard $t$-structure on
spectra and its heart $\mathbf{Ab}$.  The heart consists of the
Eilenberg--Mac Lane spectra, having only degree-zero homotopy.

The terminology records a degree-by-degree separation. A spectrum is
connective when $\pi_i=0$ for $i<0$, and coconnective when $\pi_i=0$
for $i>0$. The standard $t$-structure supplies compatible truncations
to these ranges. Its heart is their intersection: only $\pi_0$ can
remain, and it is an abelian group. Thus passing to the heart will
recover ordinary linear coefficients from coefficients with higher degrees.

In the slice $\mathcal S_{/BG}$, a pointed object means a family
$E\to BG$ with a section $BG\to E$; that section chooses the basepoint
in each fiber. Looping is performed in each fiber. A spectrum object
is consequently a family with compatible repeated deloopings, so the
base and its transport remain part of the construction.

\begin{theorem}[Stable coefficients on a groupoid]\label{t:spectra-over-bg}
For an ordinary group $G$ there is a natural equivalence
\[
 \operatorname{Sp}(\mathcal S_{/BG})
  \simeq\operatorname{Fun}(BG,\operatorname{Sp}).
\]
Its pointwise standard $t$-structure has heart
$\mathbb Z[G]\text{-}\mathbf{Mod}$.
\end{theorem}
\begin{proof}
Straightening, Theorem~\ref{h:straightening}, first identifies
$\mathcal S_{/BG}$ with $\operatorname{Fun}(BG,\mathcal S)$.
Pointed objects and finite limits in a functor category are computed
pointwise.  Therefore a sequence of pointed functors with equivalences
$E_n\simeq\Omega E_{n+1}$ is exactly a functor from $BG$ to sequences
of pointed spaces satisfying those equivalences.  This identifies the
objects, coherent maps, and their mapping spaces with
$\operatorname{Fun}(BG,\operatorname{Sp})$.

Pointwise truncation gives the stated $t$-structure.  An object in its
heart assigns an Eilenberg--Mac Lane spectrum to the unique object of
$BG$, hence an abelian group $M$.  Its transport maps give a $G$-action
by automorphisms of $M$.  Mapping spaces between heart objects are
discrete, so the coherent conditions reduce exactly to the ordinary
action identities.  Maps are additive maps commuting with the action.
These are left $\mathbb Z[G]$-modules and their linear maps.
\end{proof}

This theorem adds higher coefficient groups and coherent transport while
keeping $BG$ in the construction.  Its restriction to degree zero has
the same module category as Theorem~\ref{t:beck}.  The two descriptions
explain the relationship between ordinary internal linear objects and
stable local coefficient systems.

The ambient homotopical category matters.  We have stabilized
$\mathcal S_{/BG}$, not the ordinary nerve of $\mathbf{Grp}$.
Indeed, in the ordinary pointed category of groups the pullback
$1\times_G1$ is $1$ for every $G$.  Thus its ordinary loop construction
is trivial, and spectrum objects in that unenhanced nerve are trivial.
The homotopy fiber $*\times^h_{BG}*$ instead retains $G$.
Homotopical enhancement must precede this stabilization.

\subsection{A derived connecting map reconstructs an extension}

For this subsection $R$ is an associative unital ring and modules are left
modules.  We use cochain complexes, whose differentials raise degree.
The shift is $(C[1])^j=C^{j+1}$ with negated differential, so a module
$N[1]$ has its only nonzero cohomology in degree $-1$.

We use the standard derived framework: localizing complexes at
quasi-isomorphisms gives a stable $\infty$-category $\mathcal D(R)$;
its heart is $R\text{-}\mathbf{Mod}$; a fiber/cofiber sequence gives the
long exact cohomology sequence; and a bounded projective resolution
computes derived maps and derived tensor products.  These are foundational
inputs, not consequences of composition-factor uniqueness.  References
are \cite[\S1.3]{HA} and \cite[Tags 06XP, 0643]{Stacks}.
Write $[X,Y]$ for $\pi_0\operatorname{Map}_{\mathcal D(R)}(X,Y)$ and set
\[
 \operatorname{Ext}^1_R(L,N)=[L,N[1]].
\]

\begin{theorem}[Extensions are connecting morphisms]
\label{t:extension-classification}
The isomorphism classes of marked short exact sequences
\[
 0\longrightarrow N\xrightarrow{i}E\xrightarrow{p}L\longrightarrow0
\]
are naturally in bijection with $[L,N[1]]$.  Here an isomorphism of marked
extensions induces the identity on both $N$ and $L$.
\end{theorem}
\begin{proof}
The cone of $i:N[0]\to E[0]$ is the complex $N\to E$ in degrees
$-1,0$.  Its cohomology is $L$ in degree zero and zero elsewhere, so
the quotient map identifies it with $L$ in $\mathcal D(R)$.  The cofiber
sequence therefore supplies a connecting map
\[
 N\longrightarrow E\longrightarrow L\xrightarrow{\delta}N[1].
\]
An isomorphism of marked extensions preserves $\delta$.

Conversely, given a map $\delta:L\to N[1]$, form its fiber $E_\delta$.
The long exact cohomology sequence gives
\[
 0\longrightarrow N\longrightarrow H^0(E_\delta)
 \longrightarrow L\longrightarrow0,
\]
and $H^j(E_\delta)=0$ for $j\neq0$.  The canonical truncation maps
therefore identify $E_\delta$ with this module in the derived category.
A path between two choices of $\delta$ induces an equivalence of their
fiber diagrams, hence an isomorphism of marked extensions.

Finally, the fiber of the connecting map of the first construction is
the original $E$, with its given maps.  The connecting map recovered
from the second construction is the original $\delta$, because both
are the same fiber/cofiber diagram.  This uses the universal properties
of the fiber and the stable rotation of a cofiber sequence; the spaces
of choices of each such limit or colimit are contractible.  Thus the
two constructions are inverse on the stated isomorphism classes.
\end{proof}

\begin{proposition}[The automorphisms also survive]
\label{t:extension-automorphisms}
The automorphism group of a marked extension as above is naturally
$\operatorname{Hom}_R(L,N)$.  It agrees with the fundamental group of
$\operatorname{Map}(L,N[1])$ at its class $\delta$.
\end{proposition}
\begin{proof}
For $u:L\to N$, the map $1_E+iup$ preserves both markings and has
inverse $1_E-iup$, since $pi=0$.  Multiplication of these automorphisms
corresponds to addition of $u$.
Conversely, if $f$ preserves the markings, $p(f-1)=0$, so $f-1$ factors
uniquely as $iv$ with $v:E\to N$.  The equality $(f-1)i=0$ and
monicity of $i$ give $vi=0$.  The quotient property of $p$ makes $v=up$
for a unique $u:L\to N$.

In a stable category, looping the mapping space at zero identifies
$\pi_1\operatorname{Map}(L,N[1])$ with $[L,N]$.  Translation in this
grouplike mapping space gives the same group at $\delta$.
The fully faithful inclusion of the heart identifies $[L,N]$ with
$\operatorname{Hom}_R(L,N)$, proving the final assertion.
\end{proof}

The bijection also respects the Baer addition previously
constructed. Naturality of the connecting map makes pullback along
$L\to L\oplus L$ and pushout along $N\oplus N\to N$ act by
precomposition and postcomposition. Applied to $\delta_1\oplus\delta_2$,
their composite is $\delta_1+\delta_2$. Thus the derived notation
$\operatorname{Ext}^1$ identifies the same abelian group as the earlier
extension construction, not only the same underlying set.

The factors $N,L$ are the endpoints, and $\delta$ is the missing
attachment.  Its mapping space also retains the automorphisms of that
attachment.  This is a precise instance of gaining reconstruction data
by retaining maps rather than just a list of objects.

\begin{example}[The same two factors, different attachments]
\label{t:dual-number-extension}
Let $R=k[t]$ for a field $k$, and let $k=R/(t)$ as an $R$-module.
The free resolution
\[
 0\longrightarrow R\xrightarrow{\,t\,}R\longrightarrow k
 \longrightarrow0
\]
is exact: multiplication by $t$ is injective in the polynomial ring,
and its image is the kernel of evaluation at $0$.  Applying
$\operatorname{Hom}_R(-,k)$ makes the differential zero, hence
$\operatorname{Ext}^1_R(k,k)\cong k$.

Every marked extension can also be described without derived notation.
Let $e$ be the specified image of $1$ in its kernel and choose a lift $f$
of $1$ in its quotient.  Then $e,f$ span the underlying vector space:
subtract a multiple of $f$ from any element and the result is in $ke$.
They are independent by applying the quotient map first, then using
injectivity on the kernel.  Necessarily
\[
 te=0,\qquad tf=\lambda e
\]
for a unique $\lambda\in k$.  Replacing $f$ by $f+ce$ does not change
$\lambda$.  Conversely these formulas define an $R$-module extension
for every $\lambda$.  A map preserving the markings preserves this
scalar, and the formulas give an isomorphism whenever the scalars agree.
This is the concrete parameter detected by the derived calculation.

At $\lambda=0$ the module is $k\oplus k$.  At $\lambda\neq0$ it is
isomorphic, after forgetting the markings, to $R/(t^2)$: send $1$ to
$f$ and $t$ to $\lambda e$.  The map is surjective, has kernel $(t^2)$,
and the displayed independent vectors prove injectivity of the induced
map.  Both modules have two factors isomorphic to $k$, but their
$t$-operators are different.  The marked classification distinguishes
all values of $\lambda$, while the middle-object classification merges
all nonzero values.  These classifications must not be conflated.
\end{example}

\subsection{A derived tensor product retains an equation}

For commutative $R$, the derived tensor product is denoted
$M\otimes_R^{\mathbf L}N$.  A projective resolution replaces one input
before tensoring.  This repairs the failure of ordinary tensoring to
preserve exactness by retaining the resulting homology.

\begin{proposition}\label{t:derived-intersection}
For $R=k[t]$ and $k=R/(t)$,
\[
 H^j(k\otimes_R^{\mathbf L}k)=
 \begin{cases}k,&j=0,-1,\\0,&\text{otherwise.}\end{cases}
\]
The ordinary tensor product is only $k$.
\end{proposition}
\begin{proof}
Place the two free terms of the resolution in
Example~\ref{t:dual-number-extension} in degrees $-1,0$.
Tensoring with $k$ gives the complex
\[
 \bigl[\,k\xrightarrow{\,0\,}k\,\bigr]
\]
in those degrees, because $t$ acts as zero on $k$.
Its cohomology is exactly the displayed answer.
Ordinary tensoring gives $R/(t)\otimes_R R/(t)\cong R/(t)$: the map
$a\otimes b\mapsto ab$ has inverse $c\mapsto c\otimes1$.
\end{proof}

In affine geometry this is the derived self-intersection of the origin
inside the affine line.  The degree $-1$ copy of $k$ retains the overlap
of the two imposed copies of the equation $t=0$.  The statement about
cohomology was proved directly by the resolution; the geometric name
does not serve as its proof.  The calculation is already beyond what a
list of simple factors or an ordinary quotient alone records.

\subsection{Cotangent objects: the same method for equations and deformations}

There is a corresponding linear theory for a commutative algebra.
For an ordinary $A$-module $M$, the split square-zero algebra $A\oplus M$
has multiplication
\[
 (a,m)(b,n)=(ab,an+bm).
\]
Its ideal $0\oplus M$ has square zero.  A section of
$A\oplus M\to A$ has the form $a\mapsto(a,d(a))$.  Comparing products
shows that it is a ring map exactly when
$d(ab)=a\,d(b)+b\,d(a)$ and $d$ is additive with $d(1)=0$.
Over a base ring $R$, being an $R$-algebra map adds $d(r)=0$.
Thus derivations arise as maps into linear objects over the algebra.

For the higher statement, the notation retains these markings:
\[
 \operatorname{Der}_A(B,M)
 :=\operatorname{Map}_{(\operatorname{CAlg}_{A/})_{/B}}
       (B,B\oplus M).
\]
The source is the identity map of the $A$-algebra $B$ and the target is
its split square-zero extension, with its projection to $B$.  Thus this
is the space of sections over $B$ in $A$-algebras, including the specified
coherences with the base and the projection.

\begin{theorem}[Framework theorem, quoted: tangent modules and cotangents]
\label{t:cotangent-foundation}
For an $E_\infty$-ring $A$,
$\operatorname{Sp}(\operatorname{CAlg}_{/A})\simeq\operatorname{Mod}_A$.
For a map $A\to B$, derived $A$-derivations of $B$ into a $B$-module $M$
are represented by a cotangent module $L_{B/A}$:
\[
 \operatorname{Der}_A(B,M)
 \simeq\operatorname{Map}_{\operatorname{Mod}_B}(L_{B/A},M).
\]
For $A\to B\to C$ there is a cofiber sequence
\begin{equation}\label{t:transitivity}
 C\otimes_B L_{B/A}\longrightarrow L_{C/A}
 \longrightarrow L_{C/B}.
\end{equation}
\end{theorem}
These are the tangent equivalence, representing property, and transitivity
formalism of \cite[\S7.3; Corollaries 7.3.4.14 and 7.3.3.6]{HA}.
Tensor products here are tensor products of module spectra, hence derived.
In ordinary commutative algebra the corresponding cotangent complexes
and transitivity triangle are constructed using simplicial resolutions
\cite[Tag 08QX]{Stacks}.

The category $\operatorname{CAlg}_{/A}$ contains nonlinear algebras mapping
to $A$.  Its stabilization produces the linear coefficient category at
$A$, just as spaces over $BG$ produced spectrum-valued local coefficients.
The representing property says exactly which measurements the cotangent
module encodes: all derived derivations into any coefficient module.

\begin{proposition}[An explicit obstruction to extending a derivation]
\label{t:derivation-obstruction}
Let $A\to B\to C$ be maps of $E_\infty$-rings and $M$ a $C$-module.
A homotopy class of $A$-derivations $B\to M$ extends to an $A$-derivation
$C\to M$ if and only if its image under the connecting homomorphism
\[
 [C\otimes_B L_{B/A},M]
   \longrightarrow [L_{C/B},M[1]]
\]
vanishes.  The brackets denote morphisms in the homotopy category of
$C$-modules.
\end{proposition}
\begin{proof}
Apply the mapping-spectrum functor into $M$ to
\eqref{t:transitivity}.  Mapping out of a cofiber turns it into a fiber,
giving the exact sequence
\[
 [L_{C/A},M]\longrightarrow[C\otimes_B L_{B/A},M]
   \longrightarrow[L_{C/B},M[1]].
\]
The first arrow is restriction of derivations: naturality of the
representing property identifies the cotangent map with restriction,
and extension of scalars identifies the middle term with
$[L_{B/A},M]$ as a $B$-module mapping group.
Exactness says precisely that a class lies in the image of restriction
if and only if its connecting class is zero.
\end{proof}

This is a theorem about existence of an extension of a homotopy class.
It does not assert uniqueness.  The full mapping fiber retains the
choices of extension, their paths, and their higher compatibilities.
The transitivity triangle organizes all three algebras in one statement;
its connecting arrow identifies the obstruction rather than merely
asserting that some obstruction ought to exist.

For a final explicit computation, use ordinary commutative algebras over
a field $k$ of characteristic zero.  This choice also permits comparison
with commutative differential graded algebras and $E_\infty$-$Hk$-algebras
\cite[Proposition 7.1.4.11]{HA}; no equivalence between those models in
arbitrary positive characteristic is being assumed. The calculation below
is in ordinary derived cotangent theory. Under the rational comparison,
its cotangent object corresponds to the spectral one because both
corepresent the same derived-derivation functor, with the algebra and
coefficient module retained.

\begin{theorem}[Framework theorem, quoted: a hypersurface cotangent]
\label{t:hypersurface-input}
If $P=k[x_1,\ldots,x_n]$ and $A=P/(f)$ with $f$ a nonzero divisor,
the ordinary derived cotangent complex $L_{A/k}$ is represented in
degrees $-1,0$ by
\[
 (f)/(f^2)\longrightarrow\Omega_{P/k}\otimes_P A,
 \qquad [f]\longmapsto df.
\]
\end{theorem}
This regular-quotient computation follows from the cotangent
transitivity triangle and the regular-sequence theorem
\cite[Tags 08SH, 08QX]{Stacks}.  The ordinary module of differentials
alone is its degree-zero cohomology; the complex retains the relation
before taking that quotient.

\begin{example}[The infinitesimal direction of the dual numbers]
\label{t:cotangent-dual-numbers}
Take $A=k[t]/(t^2)$.  The polynomial $t^2$ is a nonzero divisor in
$k[t]$, so Theorem~\ref{t:hypersurface-input} applies.  Both terms are
free rank-one $A$-modules, and the complex is
\[
 L_{A/k}\simeq
 \bigl[\,A e\xrightarrow{\ e\mapsto 2t\,dt\ }A\,dt\,\bigr].
\]
The characteristic assumption makes $2$ invertible.  Multiplication by
$2t$ has image $kt$ and kernel $kt$ in $A$.  Consequently
\[
 H^0(L_{A/k})\cong k\,dt,
 \qquad H^{-1}(L_{A/k})\cong k\,(te).
\]
These assertions can be checked on $a+bt$: multiplication by $2t$
gives $2at$.  The negative-degree term therefore survives; replacing
the complex by its ordinary differentials loses it.

Let $k$ be an $A$-module through $t\mapsto0$.  Since the displayed
cotangent complex is a bounded complex of free modules, its derived
maps into $k$ are computed by the ordinary Hom complex.  That complex
has a copy of $k$ in degrees $0,1$ and zero differential.  Thus
\[
 [L_{A/k},k]\cong k,
 \qquad [L_{A/k},k[1]]\cong k.
\]
The first is the group of derivations into this coefficient module;
the second is the next group available for obstruction calculations
such as Proposition~\ref{t:derivation-obstruction}.  A nonzero possible
obstruction group does not by itself assert that a specified lifting
problem has a nonzero obstruction: that requires computing its
connecting class.
\end{example}

The progression is now visible.  Exact quotients and refinement identify
which simple pieces are invariant.  Coherent quotients retain actions and
transport.  Linear objects over a base produce modules.  Derived maps
retain attachments and automorphisms; derived tensor products and
cotangent complexes retain relations and obstruction classes.  Each
construction answers a new question about information the preceding
one did not retain.  The ordinary Jordan--H\"older theorem remains one
proved consequence of the exact framework, while these higher and
derived results enlarge the kinds of reconstruction that can be asked
and answered.

% END INPUT 11-linearization-and-derived-structure.tex

\clearpage\appendix
% BEGIN INPUT 12-dependencies-and-reading.tex
\section{Dependencies, scope, and further reading}
\label{a:dependencies}

\subsection{What was proved from what}

The ordinary comparison has a single dependency chain:
\[
\begin{tikzcd}[column sep=small,row sep=large]
\text{finite limits and regular images}\arrow[d]\\
\text{associative relation calculus}\arrow[d]\\
\text{Mal'cev difunctionality}\arrow[d]\\
\text{effective support quotients}\arrow[d]\\
\text{saturated descent and refinement cells}\arrow[d]\\
\text{Schreier refinement}\arrow[d]\\
\text{simple-factor rigidity}\arrow[d]\\
\text{Jordan--H\"older}\arrow[d]\\
\text{basis size and additive invariants}.
\end{tikzcd}
\]
Exactness enters at the effective quotients. Pointedness supplies the zero
classes used to concatenate cells into a series. The Mal'cev condition is
used where it is needed: to make every relation difunctional. Ordinary
compatible iterated quotients require regularity and exactness, but do not
require Mal'cev. The written proofs keep these dependencies separate.

The precise ordinary theorem is Theorem~\ref{j:jordan-holder}: in a
pointed Barr-exact Mal'cev category, any two given finite composition
series, retaining their regular quotient maps, have the same factors up
to permutation and isomorphism. No global chain condition is needed to
compare two finite series already supplied. This is a sufficient common
framework, not a claim that these are the logically weakest possible
axioms for every formulation of Jordan--H\"older.

The existence of a finite series is a separate issue. The basis flags
provide the required series for a vector space with a given finite basis.
Later, finite-dimensional module theory uses the now established dimension
function to prove existence of composition series. That later argument
does not feed back into the proof of dimension. One can define $K_0$ by
generators and exact-sequence relations before proving Jordan--H\"older.
What requires the factor comparison is the identification of $K_0$ with
the free abelian group on the simple isomorphism classes. Using that
freeness to prove factor uniqueness would reverse the dependency.

The source of associativity changes as the constructions develop. For the
original algebra it is an axiom equating two composites. For relation
composition, Baer addition, and Hall multiplication, it is a theorem:
both bracketings are obtained from the same complete data of witnesses,
tuples, or flags. For coherent quotients, the same universal-property
principle controls the higher compatibility of successive constructions.
These assertions specify the maps or data that do the work; a resemblance
between diagrams alone would not establish them.

\subsection{Three different comparison claims}

The scope of the page calculations is worth recording alongside the core
theorem. The categorical rectangle works for all objects and finite normal
series satisfying the exact Mal'cev hypotheses. The explicit nonabelian
page construction in Section~\ref{p:chapter} uses composition-series
input in groups with endomorphism operators; that simplicity is used to
prove the normality of differential images. The filtered module calculus
in Section~\ref{s:chapter} allows arbitrary finite filtrations and an
arbitrary filtration-preserving module map. Its convergence records
kernel and cokernel layers, and its staggered identity specialization
pairs simple factors.

Thus the nonabelian page calculation and the module spectral calculation
are compatible where both apply. Neither is asserted to provide a general
spectral sequence for arbitrary filtered groups. The categorical theorem
already supplies the unified Jordan--H\"older statement independently of
those more specialized calculations.

\subsection{The additional hypotheses behind reconstruction}

The universal additive record applies to composition-finite objects, with
the set-size qualification stated in Corollary~\ref{j:universal-monoid}.
For finite-length modules the category is also closed under submodules
and quotients, as explicitly proved. We did not assume that extra closure
for every exact Mal'cev category. Additive records ignore whether an
extension splits; the marked extension group retains its endpoint maps.

The probe theorem reconstructs finite-dimensional \emph{nilpotent}
$k[t]$-modules. Its second-difference formula is proved using canonical
subquotients, and the chain decomposition is constructed before the
reconstruction conclusion is drawn. The same tests fail for an invertible
operator, as the example shows. Hall multiplication uses finite fields
and finite-dimensional nilpotent modules so that every coefficient and
every basis expansion involved is finite. Its associativity counts the
same flags in two orders. Marked extension classes, isomorphism classes
of middle objects, and actual submodules are kept distinct.

\subsection{Where the higher foundations enter}

The higher part has its own declared inputs. Rigidification identifies
product-preserving theory models in spaces with a familiar simplicial
presentation. The Kan loop-group equivalence, loop--delooping recognition,
and Dold--Kan explain the group and module
instances. Effective groupoid quotients and universal colimits make
homotopy quotients compatible with base change. Straightening interprets
families over $BG$ as coherent actions. Stabilization and the derived
category supply mapping spectra, hearts, and fiber/cofiber sequences.
The tangent/cotangent theorems supply linear coefficients for algebras,
the representing object for derivations, and its transitivity sequence.

Those results are explicitly imported from the cited foundational
literature; this exposition does not reprove higher topos theory or the
construction of its model categories. Its deductions include the
double-coset homotopy pullback, the coefficient category over $BG$, the
marked extension classification and its automorphisms, and explicit
derived tensor and cotangent calculations. The ordinary internal-linear
equivalence $\operatorname{Ab}(\mathbf{Grp}/G)\simeq
\mathbb Z[G]\text{-}\mathbf{Mod}$ is proved directly.

The term ``higher'' does not imply an unrestricted higher
Jordan--H\"older theorem. Normality data and compatibility remain part of
the problem, and literal subobject chains from the trivial higher group
force discreteness by Proposition~\ref{h:monos-discrete}. What the higher
framework contributes here is a precise enlargement of the retained data
and several new reconstruction tools. Its ordinary discrete models
recover the common theorem when its exact Mal'cev hypotheses are checked.

\subsection{Audit and reproducibility}

The \href{https://tensorgami.com/notes/quotients-and-reconstruction-audit.html}{public audit report}
inventories all theorem-like statements, their proof or imported-foundation
status, the diagrams, references, and edition fingerprints. It distinguishes
the full review of the preceding edition from the subsequent check of these
small clarifications and the matching web, PDF, and source files.
The ordinary core and the higher chapters received
independent mathematical rereads; integration checks tracked changed
action conventions, quotient maps, normality, degrees, and dependencies.
Supplementary finite group and finite field calculations test examples
of the actual maps and counts. Such computations do not establish the
general theorems. Compilation and inspection of the rendered pages are
separate presentation checks, not proof certification. No proof assistant
formalization or historical priority claim is made.

The earlier edition is preserved separately. This edition reorganizes
the core around relations and exactness, so that the group and module
comparison theorems are genuine instances. It retains the explicit
extension and reconstruction calculations, and places the new higher
constructions after the common ordinary proof.

\begin{thebibliography}{99}
\addcontentsline{toc}{section}{References}
\bibitem{Knop}
Friedrich Knop, \emph{Tensor envelopes of regular categories},
Advances in Mathematics \textbf{214} (2007), 571--617.
\href{https://arxiv.org/pdf/math/0610552}{Author preprint};
Section 5 and Appendices A--B, especially Appendix B2.

\bibitem{HTT}
Jacob Lurie, \emph{Higher Topos Theory}, Princeton University Press, 2009.
\href{https://www.math.ias.edu/~lurie/papers/HTT.pdf}{Author PDF}.
See Corollary 5.5.9.3, Theorem 6.1.0.6, and Sections 2.2 and 3.2.

\bibitem{HA}
Jacob Lurie, \emph{Higher Algebra}, September 18, 2017.
\href{https://www.math.ias.edu/~lurie/papers/HA.pdf}{Author PDF}.
See Sections 1.3, 1.4.2--1.4.3, 5.2.6 and 7.3;
Theorem 7.1.2.13; Proposition 7.1.4.11;
Corollaries 7.3.3.6 and 7.3.4.14.

\bibitem{Badzioch}
Bernard Badzioch, \emph{Algebraic theories in homotopy theory},
Annals of Mathematics \textbf{155} (2002), 895--913.
\href{https://annals.math.princeton.edu/2002/155-3/p06}{Journal article};
\href{https://arxiv.org/abs/math/0110101}{Author preprint}.

\bibitem{GoerssJardine}
Paul G. Goerss and John F. Jardine, \emph{Simplicial Homotopy Theory},
Progress in Mathematics \textbf{174}, Birkh\"auser, 1999.
\href{https://www.sas.rochester.edu/mth/sites/doug-ravenel/otherpapers/Goerss-Jardine2.pdf}{Book PDF};
Chapter V, Proposition 6.3 and Corollary 6.4.

\bibitem{BeardsleyFox}
Jonathan Beardsley and Landon Fox, \emph{Higher Groups and Higher Normality},
\href{https://arxiv.org/html/2407.21210v3}{arXiv:2407.21210v3},
August 13, 2026. In particular Definition 4.1 and Section 5,
including the compatibility condition in Theorem 5.14.

\bibitem{Stacks}
The Stacks Project Authors, \emph{The Stacks Project},
\href{https://stacks.math.columbia.edu/}{online reference}.
The specific tags used are
\href{https://stacks.math.columbia.edu/tag/010I}{010I} (extensions),
\href{https://stacks.math.columbia.edu/tag/012K}{012K} (filtered complexes),
\href{https://stacks.math.columbia.edu/tag/019G}{019G} (Dold--Kan),
\href{https://stacks.math.columbia.edu/tag/06XP}{06XP} (Ext),
\href{https://stacks.math.columbia.edu/tag/0643}{0643} (projective resolutions),
\href{https://stacks.math.columbia.edu/tag/08QX}{08QX} (cotangent transitivity),
and \href{https://stacks.math.columbia.edu/tag/08SH}{08SH} (regular quotients).

\bibitem{Deligne}
Pierre Deligne, \emph{Th\'eorie de Hodge II},
Publications Math\'ematiques de l'IH\'ES \textbf{40} (1971), 5--57.
\href{https://www.numdam.org/item/10.1007/BF02684692.pdf}{Article PDF};
Section 1.2.1 discusses interchange of associated gradings.

\bibitem{Grandis}
Marco Grandis, \emph{Homotopy spectral sequences},
\href{https://arxiv.org/pdf/1007.0632}{arXiv:1007.0632}, 2010.
Sections 2--4 give a broader nonabelian exact-couple setting.
\end{thebibliography}

% END INPUT 12-dependencies-and-reading.tex

\end{document}
